% !TeX program = pdflatex
\documentclass[11pt]{preprint}
\usepackage[T1]{fontenc}
\usepackage[utf8]{inputenc}
\usepackage{amsmath,amssymb,amsthm}
\usepackage{mathtools}
\usepackage{mathrsfs}
\usepackage[a4paper,margin=28mm]{geometry}
\usepackage{enumitem}
\usepackage{graphicx}
\usepackage[hidelinks]{hyperref}
\newcommand{\dd}{\mathrm{d}}
\newcommand{\tr}{\operatorname{tr}}
\newcommand{\Tr}{\operatorname{Tr}}
\newcommand{\id}{\mathbf 1}
\newcommand{\re}{\operatorname{Re}}
\newcommand{\im}{\operatorname{Im}}
\newcommand{\diam}{\operatorname{diam}}
\newcommand{\dist}{\operatorname{dist}}
\DeclarePairedDelimiter{\abs}{\lvert}{\rvert}
\DeclarePairedDelimiter{\norm}{\lVert}{\rVert}
\DeclarePairedDelimiter{\floor}{\lfloor}{\rfloor}
\DeclarePairedDelimiter{\ceil}{\lceil}{\rceil}
\setlength{\parskip}{0.6em}
\emergencystretch=3em
\title{The Renormalisation Group for Four-Dimensional Lattice Gauge Theories: A Self-Contained Account of the Results Used for the Continuum Limit of Wilson Loops}
\author{}
\date{}

\numberwithin{equation}{section}
\newtheorem{theorem}{Theorem}[section]
\newtheorem{proposition}[theorem]{Proposition}
\newtheorem{lemma}[theorem]{Lemma}
\newtheorem{corollary}[theorem]{Corollary}
\theoremstyle{definition}
\newtheorem{definition}[theorem]{Definition}
\theoremstyle{remark}
\newtheorem{remark}[theorem]{Remark}
\theoremstyle{plain}
\renewcommand{\qedsymbol}{\textit{q.e.d.}}
\begin{document}
\maketitle

\begin{abstract}
We give an account of the part of Ba{\l}aban's renormalisation group for lattice gauge theories in four dimensions on which the construction of the continuum limit of Wilson loops rests: the averaging operations, the propagators in a background field, the gauge conditions, the variational problem for the background field, the renormalisation step for small fields, the complete step with large-field regions, and the $\mathbf R$-operation. Each object is defined once, at the generality in which it is used, every constant is named, and no step of a proof is delegated to the original papers, which are cited for attribution only.
\end{abstract}

\setcounter{tocdepth}{2}
{\setlength{\parskip}{0pt}\tableofcontents}

\setcounter{section}{-1}
\section{Introduction}\label{sec:introduction}

Ba{\l}aban's construction of the effective densities of four-dimensional lattice gauge theories is contained in a series of papers written over six years, from the regularity of lattice Green's functions \cite{BalabanRDG} to the renormalisation of the large fields \cite{BalabanLFR1,BalabanLFR2}. A reader who wants to use the construction, and not only its final theorems, meets two difficulties. The first is that the series was not written as one text. The averaging operation which defines the renormalisation transformation is changed between the papers on the linear theory and those on the renormalisation group, because the earlier one ``is not symmetric with respect to lattice Euclidean transformations'' \cite[p.~252]{BalabanRG1}, and the symbol $\abs X$ denotes the operator norm of a matrix in one paper \cite[(19), p.~21]{BalabanAVG} and the Hilbert--Schmidt norm in another \cite[p.~390]{BalabanPBG}. The second difficulty is that many steps are taken over by a reference, to an earlier paper of the series, to the three-dimensional theory \cite{BalabanUV3}, or to papers on other models, with the remark that the argument is the same.

The paper \emph{The Continuum Limit of the Correlation Functions of Wilson Loops in Four-Dimensional Lattice Yang--Mills Theory}, which we call the main paper, uses the inside of the construction throughout. It inserts loops into the densities and follows them through every operation of a step, it compares the expansions of the propagators term by term on two lattices, and it runs the construction with parameters and test functions of its own. For such a use the original statements do not suffice, since what is needed at each place is a definition, an identity, or the dependence of a term on the fields in a region, and these are found in the proofs. The purpose of the present account is to state and to prove what the main paper and its supplement use, in one notation and in the order of the logic, so that neither of them has to be read with the original papers at hand.

Two principles decide what is included. A statement appears here only if a proof in the main paper, in its supplement, or later in this account uses it. Furthermore, every object is defined once, at the widest generality in which it is used: the dimension is four; there is one averaging operation, that of \cite[(0.4), p.~253]{BalabanRG1}; configurations take their values in the complexified group from the start, the gauge fields being the special case; and the parameters, the test functions, and the depth of the $\mathbf R$-operation are those of the construction of the main paper. As a consequence, a result which the original papers prove for real fields and for the earlier average, and which is then needed for complex fields and for the later one, is proved here once, in the form in which it is needed. We do not aim at reproducing the original statements. What a statement has to deliver is fixed by the proofs of the main paper which use it, and where a definition or a bound of the papers can be replaced by one which is simpler to state or to prove and which serves those proofs equally, we use the simpler one.

The statements are written so that they can be used without their proofs. For each object we give its definition, the identities which it satisfies, the set of variables on which it depends, the domain on which it is analytic, and its bounds, in this order and as separate clauses. Every constant has a name and a place in an order of choices, a constant being allowed to depend only on those chosen before it. A statement which holds ``for sufficiently small'' values of a parameter names the threshold.

\emph{The state of this version.} This version contains Section~\ref{sec:setting}, which fixes the lattices, the fields, the norms, and the geometry of the domains, and Section~\ref{sec:averaging}, on the average and the relative average. In Section~\ref{sec:setting}, Lemmata~\ref{lem:norms}, \ref{lem:action}, \ref{lem:covariance}, \ref{lem:distance}, and~\ref{lem:cube-count} are proved, and Lemma~\ref{lem:lie-algebra}, the first inequality of Lemma~\ref{lem:summation}, and Lemma~\ref{lem:contracted-geometry} are stated without proof. The statements of Section~\ref{sec:averaging} are given without proofs. Section~\ref{sec:propagators} contains the definitions of the operators of a sequence of domains and of the regular pairs, the bounds of the four basic operators, and the walk expansion of the scalar operator with the first approximations of the other two, all without proofs. It also contains the operators of the Landau gauge at a pair, the propagator of the fluctuation, and the dependence of the operators on the sequence of domains, with their bounds and identities. Lemma~\ref{lem:composition}, on the composition of bounds between blocks, is proved. The expansions of the block inverse and of the operator $G$ into walks of local factors are stated as an existence theorem, Theorem~\ref{thm:walk-structure}(ii); the walks themselves are not yet defined. Section~\ref{sec:trivial-background} states the momentum representation at the trivial background, without proof. Section~\ref{sec:gauges} defines the regular fields, the two transports, and the axial and the Landau gauge; Lemma~\ref{lem:axial-representative} is proved, and Lemma~\ref{lem:axial-smallness} and Theorem~\ref{thm:landau-gauge} are stated without proof. Section~\ref{sec:variational} states the variational problem and its solution: Lemma~\ref{lem:problem-covariance}, Corollary~\ref{cor:minimiser-regular}, and Lemmata~\ref{lem:functional} and~\ref{lem:support} are proved from the statements before them, and Lemma~\ref{lem:critical}(i), Theorems~\ref{thm:minimiser}, \ref{thm:critical}, \ref{thm:background-field}, and~\ref{thm:support}, and Propositions~\ref{prop:linearising}, \ref{prop:potential}, \ref{prop:second-operators}, and~\ref{prop:fixed-points} are stated without proof. Not yet stated in that section are the minimiser with a part of the data left free, the bound of the plaquettes of the data by the action, and the bounds of the second order of the solutions at a pair whose second entry is not the current of the first. Section~\ref{sec:small-field} is begun: it defines the spaces of pairs of one scale and states the criterion for membership, Proposition~\ref{prop:membership}, without proof. Not yet stated are the walks of the operators of the Landau gauge, the gauge transformation at a complex base, the variants of Theorem~\ref{thm:landau-gauge} without the bound of the current and on cubes, and the covariance of the unit lattice, which belongs to Section~7.

\section{Lattices, Fields, Norms, and Domains}\label{sec:setting}

Nothing in this section is an estimate of the field theory. We fix the group and the norms on its algebra, the lattices of the scales which a renormalisation step relates, the fields on them with the action and the covariant difference operators, and then the two geometric structures in terms of which every later bound is stated: the sequences of domains on which the field lives at different scales, with their weighted distance, and the localisation domains with their linear size.

\subsection{The Gauge Group}\label{sec:group}

Throughout, $N \geq 1$ is an integer, $M_N(\mathbb C)$ is the algebra of the complex $N \times N$ matrices, with unit $\id$, and $G$ is a compact connected Lie subgroup of the unitary group $U(N)$. We follow \cite[Sect.~A]{BalabanAVG} in writing the elements of the group near the unit as $e^{iA}$ with a Hermitian matrix $A$, so that the Lie algebra consists of Hermitian matrices and its bracket is $(A, B) \mapsto i[A, B]$.

\begin{definition}[Algebra, Complexification, and Norms]\label{def:group}
\begin{enumerate}[label=(\roman*)]
\item The \emph{Lie algebra} of $G$ is $\mathfrak g \coloneqq \{ A \in M_N(\mathbb C) \,|\, A^* = A \text{ and } e^{itA} \in G \text{ for all } t \in \mathbb R \}$, and $\mathfrak g^c \coloneqq \{ A + iB \,|\, A, B \in \mathfrak g \}$ is its \emph{complexification}. The \emph{complexified group} $G^c$ is the subgroup of $GL(N, \mathbb C)$ generated by the matrices $e^{iA}$ with $A \in \mathfrak g^c$.
\item The \emph{normalised trace} is $\tr X \coloneqq N^{-1}\Tr X$. On $M_N(\mathbb C)$ we use the scalar product $\langle X, Y\rangle \coloneqq \tr X^*Y$, the \emph{Hilbert norm} $h(X) \coloneqq (\tr X^*X)^{1/2}$, and the operator norm $\norm X$ with respect to the Euclidean norm of $\mathbb C^N$. The \emph{bond norm} $\abs X$ is one of these two norms, chosen once and for all.
\item $\pi$ is the orthogonal projection of $M_N(\mathbb C)$ onto $\mathfrak g^c$ with respect to $\langle\cdot, \cdot\rangle$.
\item For an invertible matrix $\mathbf U$ and $X \in M_N(\mathbb C)$ we write $R(\mathbf U)X \coloneqq \mathbf UX\mathbf U^{-1}$.
\end{enumerate}
\end{definition}

The bond norm is the norm in which the fields are compared with their thresholds, and we do not decide between its two candidates: every statement of this account holds for either choice, the constants being allowed to depend on $N$. What makes this possible is that the two norms differ by the factor $\sqrt N$ at most and share the properties of Lemma~\ref{lem:norms}, which are the only ones that the proofs use. The Hilbert norm is the norm of the action, by \eqref{eq:plaquette-norm} below, and of the scalar products of fields. We need the following facts about the algebra.

\begin{lemma}[The Lie Algebra]\label{lem:lie-algebra}
The set $\mathfrak g$ is a real linear subspace of $M_N(\mathbb C)$, and $i[A, B] \in \mathfrak g$ for all $A, B \in \mathfrak g$. Consequently $\mathfrak g^c$ is a complex linear subspace of $M_N(\mathbb C)$ which is closed under the commutator and under $X \mapsto X^*$, and $\mathfrak g \cap i\mathfrak g = \{ 0 \}$.
\end{lemma}

\begin{lemma}[Norms and the Adjoint Action]\label{lem:norms}
Let $X, Y \in M_N(\mathbb C)$ and $U, V \in U(N)$.
\begin{enumerate}[label=(\roman*)]
\item $\abs{\tr X} \leq h(X) \leq \abs X \leq \norm X \leq \sqrt N\,h(X)$. Each of the three norms $\mathcal N \in \{ h, \abs\cdot, \norm\cdot \}$ satisfies $\mathcal N(XY) \leq \norm X\,\mathcal N(Y)$, $\mathcal N(XY) \leq \mathcal N(X)\norm Y$, and $\mathcal N(X^*) = \mathcal N(X)$.
\item $\mathcal N(UXV) = \mathcal N(X)$ for each of the three norms, and
\begin{equation}\label{eq:plaquette-norm}
    1 - \re\tr U = \frac{1}{2}h(U - \id)^2\;.
\end{equation}
\item For $\mathbf U \in G^c$ the map $R(\mathbf U)$ maps $\mathfrak g^c$ onto itself and commutes with $\pi$. For $U \in G$ it maps $\mathfrak g$ onto itself and preserves the three norms.
\item $\pi$ is complex linear, $h(\pi X) \leq h(X)$, $\abs{\pi X} \leq \sqrt N\abs X$, and $\pi(X^*) = (\pi X)^*$.
\end{enumerate}
\end{lemma}

\begin{proof}
(i) Let $s_1, \dots, s_N$ be the singular values of $X$. Then $h(X)^2 = N^{-1}\sum_is_i^2$ and $\norm X = \max_is_i$, which gives $h(X) \leq \norm X \leq \sqrt N\,h(X)$, and the bond norm is one of $h$ and $\norm\cdot$. Since $h(\id) = 1$, the Cauchy--Schwarz inequality gives $\abs{\tr X} = \abs{\langle\id, X\rangle} \leq h(X)$. Next, $X^*X \leq \norm X^2\id$ as Hermitian matrices, so that $h(XY)^2 = \tr Y^*X^*XY \leq \norm X^2\tr Y^*Y$, while the operator norm is submultiplicative; this is the first product inequality for $h$ and for $\norm\cdot$, hence for the bond norm. The matrices $Z$ and $Z^*$ have the same singular values, so that the three norms are invariant under the adjoint, and the second product inequality is the first applied to $Y^*X^*$.

(ii) The trace is cyclic and $U^*U = V^*V = \id$, so that $\tr (UXV)^*UXV = \tr X^*X$, and the operator norm is invariant because $U$ and $V$ are isometries of $\mathbb C^N$. Finally $h(U - \id)^2 = \tr(2\id - U - U^*) = 2 - 2\re\tr U$.

(iii) For $A \in \mathfrak g^c$ the map $R(e^{iA})$ is the exponential $\sum_n(n!)^{-1}(i\operatorname{ad}A)^n$ of the map $\operatorname{ad}A \colon X \mapsto [A, X]$, which preserves the finite-dimensional space $\mathfrak g^c$ by Lemma~\ref{lem:lie-algebra}. Hence $R(e^{iA})\mathfrak g^c \subset \mathfrak g^c$, the same holds for products of such matrices and for their inverses, and so $R(\mathbf U)\mathfrak g^c = \mathfrak g^c$ for all $\mathbf U \in G^c$. To see that $R(\mathbf U)$ commutes with $\pi$, let $\mathfrak n$ be the orthogonal complement of $\mathfrak g^c$. Since $\mathfrak g^c$ is closed under the adjoint, a matrix $Y$ lies in $\mathfrak n$ if and only if $\tr ZY = 0$ for all $Z \in \mathfrak g^c$. The bilinear form $(Z, Y) \mapsto \tr ZY$ is invariant under $R(\mathbf U)$, and $R(\mathbf U)^{-1} = R(\mathbf U^{-1})$ maps $\mathfrak g^c$ onto itself, so that $R(\mathbf U)\mathfrak n \subset \mathfrak n$. A linear map which preserves both $\mathfrak g^c$ and $\mathfrak n$ commutes with the projection onto the first along the second. For $U \in G$ and $A \in \mathfrak g$ the matrix $UAU^{-1}$ is Hermitian and $e^{itUAU^{-1}} = Ue^{itA}U^{-1} \in G$, so that $R(U)\mathfrak g \subset \mathfrak g$, with equality because the same holds for $U^{-1}$; the norms are preserved by (ii).

(iv) An orthogonal projection onto a complex subspace is complex linear and does not increase the norm of the scalar product, which is $h$; with (i), $\abs{\pi X} \leq \sqrt N\,h(\pi X) \leq \sqrt N\,h(X) \leq \sqrt N\abs X$. If $Y \in \mathfrak n$, then $\tr ZY^* = \overline{\tr YZ^*} = 0$ for all $Z \in \mathfrak g^c$, because $Z^* \in \mathfrak g^c$, so that $Y^* \in \mathfrak n$ by the criterion in the proof of (iii). Writing $X = \pi X + Y$ with $Y \in \mathfrak n$, we obtain $X^* = (\pi X)^* + Y^*$ with $(\pi X)^* \in \mathfrak g^c$ and $Y^* \in \mathfrak n$, which is the last assertion.
\end{proof}

\subsection{The Lattices of a Step}\label{sec:lattices}

A renormalisation step relates fields on lattices whose spacings differ by powers of a fixed odd integer. We describe all of them inside one torus, in the units of \cite[pp.~251--252]{BalabanRG1}, in which the coarsest lattice has the spacing one. The only difference from that description is the position of the sites, which we place at the points of the lattices $a\mathbb Z^4$ and not at the centres of their cells; since the spacings differ by odd factors, the lattices are nested in both descriptions, and a site of a coarse lattice is the centre of its block of finer sites.

\begin{definition}[Lattices and Blocks]\label{def:lattices}
Let $L \geq 3$ be an odd integer and let $k \geq 0$ and $m_{\mathbb T} \geq 1$ be integers. We put
\begin{equation}\label{eq:spacings}
    \eta \coloneqq L^{-k}\;,
    \qquad
    a_j \coloneqq L^j\eta \quad (0 \leq j \leq k)\;,
    \qquad
    \ell_{\mathbb T} \coloneqq 2L^{m_{\mathbb T}}\;.
\end{equation}
\begin{enumerate}[label=(\roman*)]
\item The \emph{torus} is $\mathbb T \coloneqq (\mathbb R/\ell_{\mathbb T}\mathbb Z)^4$, with the coordinate unit vectors $e_0, \dots, e_3$. For $z \in \mathbb T$ we write $\abs z_\infty$ and $\abs z_1$ for the smallest value of the $\ell^\infty$-norm and of the $\ell^1$-norm on the representatives of $z$ in $\mathbb R^4$, and distances of points and of sets are taken with respect to $\abs\cdot_\infty$ unless another norm is named.
\item The \emph{lattice of the scale $j$} is $T^{(j)} \coloneqq (a_j\mathbb Z/\ell_{\mathbb T}\mathbb Z)^4 \subset \mathbb T$, for $0 \leq j \leq k$, and $T_\eta \coloneqq T^{(0)}$ is the \emph{fine lattice}. A length \emph{in the units of the scale $j$} is a length divided by $a_j$.
\item For a number $s$ of the form $L^ra_j$ with integers $0 \leq j \leq k$ and $r \geq 0$ and with $s \leq L^{m_{\mathbb T}}$, the \emph{grid of side $s$} is the family of the closed cubes $\{ z \in \mathbb T \,|\, \abs{z - c}_\infty \leq s/2 \}$ with centres $c \in (s\mathbb Z/\ell_{\mathbb T}\mathbb Z)^4$. The cube of the grid of side $a_j$ with the centre $y \in T^{(j)}$ is denoted by $\Delta_j(y)$.
\item For $y \in T^{(j)}$ and $0 \leq n \leq j$ the \emph{block of order $n$} of $y$ is $B^n(y) \coloneqq \Delta_j(y) \cap T^{(j-n)}$, and $B(y) \coloneqq B^1(y)$.
\end{enumerate}
\end{definition}

Since $L$ is odd, $T^{(j)} \subset T^{(j-1)}$, every cube of the grid of side $Ls$ is the union of $L^4$ cubes of the grid of side $s$, no point of a lattice $T^{(i)}$ lies on the boundary of a cube $\Delta_j(y)$ with $j \geq i$, and
\begin{equation}\label{eq:block}
    B^n(y) = \bigl\{ y + a_{j-n}\textstyle\sum_\mu n_\mu e_\mu \,\big|\, n_\mu \in \mathbb Z,\ \abs{n_\mu} \leq \frac{1}{2}(L^n - 1) \bigr\}\;,
\end{equation}
a set of $L^{4n}$ points, which is the block of \cite[(0.3), p.~252]{BalabanRG1}. The blocks $B^n(y)$, $y \in T^{(j)}$, form a partition of $T^{(j-n)}$. The condition $m_{\mathbb T} \geq 1$ ensures that each lattice has at least six sites in every direction.

\begin{definition}[Bonds, Paths, and Plaquettes]\label{def:bonds}
Let $0 \leq j \leq k$ and $a \coloneqq a_j$.
\begin{enumerate}[label=(\roman*)]
\item A \emph{bond} of $T^{(j)}$ is an ordered pair $b = (b_-, b_+)$ of sites with $b_+ = b_- + ae_\mu$ or $b_+ = b_- - ae_\mu$ for some $\mu$; in the first case it is \emph{positive}. The \emph{reversed bond} is $b^{-1} \coloneqq (b_+, b_-)$.
\item A \emph{path} is a finite sequence $\gamma = (b_1, \dots, b_n)$ of bonds with $(b_i)_+ = (b_{i+1})_-$ for $1 \leq i < n$; it leads from $(b_1)_-$ to $(b_n)_+$, its \emph{length} is $na$, and its reverse is $\gamma^{-1} \coloneqq (b_n^{-1}, \dots, b_1^{-1})$. A path from $x$ to $x'$ is \emph{shortest} if its length is $\abs{x - x'}_1$.
\item For $x \in T^{(j)}$ and $\mu \neq \nu$ the \emph{plaquette} $p_{\mu\nu}(x)$ is the closed path
\[
    \partial p_{\mu\nu}(x) \coloneqq \bigl((x, x + ae_\mu),\ (x + ae_\mu, x + ae_\mu + ae_\nu),\ (x + ae_\mu + ae_\nu, x + ae_\nu),\ (x + ae_\nu, x)\bigr)\;,
\]
which is \emph{positive} if $\mu < \nu$. The path $\partial p_{\nu\mu}(x)$ is the reverse of $\partial p_{\mu\nu}(x)$.
\item For a set $S \subset \mathbb T$ we write $S^{(j)} \coloneqq S \cap T^{(j)}$. The \emph{bonds of $S$ of the scale $j$} are the bonds of $T^{(j)}$ with at least one end in $S$, and the \emph{plaquettes of $S$ of the scale $j$} are the plaquettes with at least one corner in $S$.
\end{enumerate}
\end{definition}

The convention (iv) is that of \cite[p.~224]{BalabanPRT2} and \cite[p.~277]{BalabanVAR}.

\subsection{Configurations and the Action}\label{sec:configurations}

The fields of the theory are functions on the bonds with values in the group. All later analyticity statements concern their extension to the complexified group, and so we define the objects for it at once.

\begin{definition}[Configurations and Forms]\label{def:configurations}
Let $0 \leq j \leq k$.
\begin{enumerate}[label=(\roman*)]
\item A \emph{configuration} on $T^{(j)}$ is a map $\mathbf U$ from the bonds of $T^{(j)}$ to $G^c$ with $\mathbf U(b^{-1}) = \mathbf U(b)^{-1}$. It is a \emph{gauge field} if its values lie in $G$; gauge fields are denoted by $U$, $V$. For a path $\gamma = (b_1, \dots, b_n)$ we put $\mathbf U(\gamma) \coloneqq \mathbf U(b_1)\cdots\mathbf U(b_n)$.
\item A \emph{gauge transformation} is a map $u \colon T^{(j)} \to G^c$. It acts on configurations by
\begin{equation}\label{eq:gauge-transformation}
    \mathbf U^u(b) \coloneqq u(b_-)\mathbf U(b)u(b_+)^{-1}\;,
\end{equation}
so that $\mathbf U^u(\gamma) = u(x)\mathbf U(\gamma)u(x')^{-1}$ for a path $\gamma$ from $x$ to $x'$.
\item A \emph{$0$-form} is a map $\lambda \colon T^{(j)} \to M_N(\mathbb C)$; a \emph{$1$-form} is a family $A = (A_\mu)_{0 \leq \mu \leq 3}$ of $0$-forms, and $A(b) \coloneqq A_\mu(x)$ for the positive bond $b = (x, x + a_je_\mu)$; a \emph{$2$-form} is a family $F = (F_{\mu\nu})_{\mu \neq \nu}$ of $0$-forms with $F_{\nu\mu} = -F_{\mu\nu}$. For a gauge transformation $u$ and a form $f$ we write $R(u)f$ for the form with the components $x \mapsto R(u(x))f(x)$.
\item For two forms of the same degree,
\begin{equation}\label{eq:pairing}
    \langle\lambda, \lambda'\rangle \coloneqq \sum_xa_j^4\tr\lambda(x)\lambda'(x)\;,
    \qquad
    \langle A, A'\rangle \coloneqq \sum_\mu\langle A_\mu, A'_\mu\rangle\;,
    \qquad
    \langle F, F'\rangle \coloneqq \sum_{\mu < \nu}\langle F_{\mu\nu}, F'_{\mu\nu}\rangle\;,
\end{equation}
and $\norm f^2 \coloneqq \sum_xa_j^4\sum h(f(x))^2$, the inner sum over the components $f_\mu$ of a $1$-form and $f_{\mu\nu}$, $\mu < \nu$, of a $2$-form.
\end{enumerate}
\end{definition}

The pairing \eqref{eq:pairing} is bilinear and not sesquilinear, as in \cite[(3.11), p.~392]{BalabanPBG}. On forms with Hermitian values it is the real scalar product whose norm is $\norm\cdot$, and it is the form with respect to which adjoints are holomorphic in the configuration.

\begin{definition}[The Action]\label{def:action}
For a configuration $\mathbf U$ on $T^{(j)}$ and a plaquette $p$ let
\begin{equation}\label{eq:real-imaginary}
    \re\mathbf U(\partial p) \coloneqq \frac{1}{2}\bigl(\mathbf U(\partial p) + \mathbf U(\partial p)^{-1}\bigr)\;,
    \qquad
    \im\mathbf U(\partial p) \coloneqq \frac{1}{2i}\bigl(\mathbf U(\partial p) - \mathbf U(\partial p)^{-1}\bigr)\;.
\end{equation}
The \emph{action} of $\mathbf U$ is
\begin{equation}\label{eq:action}
    A(\mathbf U) \coloneqq \sum_p\bigl(1 - \tr\re\mathbf U(\partial p)\bigr)\;,
\end{equation}
the sum over the positive plaquettes of $T^{(j)}$.
\end{definition}

This is the action \cite[(0.2), p.~252]{BalabanRG1} in dimension four, where it carries no power of the lattice spacing, with the extension to the complexified group of \cite[p.~391]{BalabanPBG}.

\begin{lemma}[The Action]\label{lem:action}
\begin{enumerate}[label=(\roman*)]
\item For a gauge field $U$, $A(U) = \sum_p(1 - \re\tr U(\partial p)) = \frac{1}{2}\sum_ph(U(\partial p) - \id)^2 \geq 0$.
\item $A(\mathbf U^u) = A(\mathbf U)$ for every configuration $\mathbf U$ and every gauge transformation $u$ with values in $G^c$, and $\im\mathbf U^u(\partial p_{\mu\nu}(x)) = R(u(x))\im\mathbf U(\partial p_{\mu\nu}(x))$.
\item The summand of \eqref{eq:action} for $p$ is a polynomial in the entries of the matrices $\mathbf U(b)$ and $\mathbf U(b)^{-1}$ for the four bonds $b$ of $\partial p$.
\end{enumerate}
\end{lemma}

\begin{proof}
For a unitary matrix $W$ we have $W^{-1} = W^*$, so that $\tr\re W = \re\tr W$, and (i) follows from \eqref{eq:plaquette-norm}. By \eqref{eq:gauge-transformation}, $\mathbf U^u(\partial p_{\mu\nu}(x)) = u(x)\mathbf U(\partial p_{\mu\nu}(x))u(x)^{-1}$, and the same holds for the inverse, which gives the second assertion of (ii); the first follows because the trace is invariant under conjugation. Finally (iii) is the definition, since $\mathbf U(\partial p)^{-1}$ is the product of the inverses in the reverse order.
\end{proof}

\subsection{Covariant Differences and the Current}\label{sec:differences}

The quadratic forms of the linear theory and the regularity conditions on the fields are written with difference operators which are covariant with respect to a configuration. We fix a scale $j$, write $a \coloneqq a_j$, and let $\mathbf U$ be a configuration on $T^{(j)}$.

\begin{definition}[Covariant Differences]\label{def:differences}
For a $0$-form $\lambda$ and $0 \leq \mu \leq 3$,
\begin{align}
    (D_{\mathbf U,\mu}\lambda)(x) &\coloneqq a^{-1}\bigl(R(\mathbf U(x, x + ae_\mu))\lambda(x + ae_\mu) - \lambda(x)\bigr)\;,\label{eq:forward-difference}\\
    (D^*_{\mathbf U,\mu}\lambda)(x) &\coloneqq a^{-1}\bigl(R(\mathbf U(x, x - ae_\mu))\lambda(x - ae_\mu) - \lambda(x)\bigr)\;.\label{eq:backward-difference}
\end{align}
With these we define, for a $0$-form $\lambda$, a $1$-form $A$, and a $2$-form $F$:
\begin{enumerate}[label=(\roman*)]
\item the $1$-form $D_{\mathbf U}\lambda$ with the components $D_{\mathbf U,\mu}\lambda$, and the $0$-form $D^*_{\mathbf U}A \coloneqq \sum_\mu D^*_{\mathbf U,\mu}A_\mu$;
\item the $2$-form $D_{\mathbf U}A$ with the components $(D_{\mathbf U}A)_{\mu\nu} \coloneqq D_{\mathbf U,\mu}A_\nu - D_{\mathbf U,\nu}A_\mu$, and the $1$-form $D^*_{\mathbf U}F$ with the components $(D^*_{\mathbf U}F)_\mu \coloneqq \sum_{\nu \neq \mu}D^*_{\mathbf U,\nu}F_{\nu\mu}$;
\item the \emph{covariant gradient} $\nabla_{\mathbf U}A$, the family of the sixteen $0$-forms $(\nabla_{\mathbf U}A)_{\mu\nu} \coloneqq D_{\mathbf U,\mu}A_\nu$, and likewise $\nabla_{\mathbf U}\lambda \coloneqq D_{\mathbf U}\lambda$;
\item the \emph{covariant Laplacian} $\Delta_{\mathbf U} \coloneqq \sum_\mu D^*_{\mathbf U,\mu}D_{\mathbf U,\mu}$, which acts on $0$-forms and, component by component, on $1$-forms.
\end{enumerate}
\end{definition}

These are the operators of \cite[(3.3), (3.4), (3.8), (3.9), pp.~390--392, and (3.23), p.~394]{BalabanPBG}. In that paper a $1$-form is extended to the negative bonds by $A(b^{-1}) = -A(b)$, and the difference of a $1$-form on a plaquette $p_{\mu\nu}(x)$ is written as a sum over its four bonds \cite[(3.4)]{BalabanPBG}; with the transports to the corner $x$ which that formula contains, the sum is the component $(D_{\mathbf U}A)_{\mu\nu}(x)$ above.

\begin{lemma}[Adjoints and Covariance]\label{lem:covariance}
Let $\mathbf U$ be a configuration on $T^{(j)}$ and $u$ a gauge transformation, both with values in $G^c$.
\begin{enumerate}[label=(\roman*)]
\item For $0$-forms $\lambda$, $\lambda'$, a $1$-form $A$, and a $2$-form $F$,
\[
    \langle D_{\mathbf U,\mu}\lambda, \lambda'\rangle = \langle\lambda, D^*_{\mathbf U,\mu}\lambda'\rangle\;,
    \qquad
    \langle D_{\mathbf U}\lambda, A\rangle = \langle\lambda, D^*_{\mathbf U}A\rangle\;,
    \qquad
    \langle D_{\mathbf U}A, F\rangle = \langle A, D^*_{\mathbf U}F\rangle\;.
\]
\item $D_{\mathbf U^u,\mu}R(u)\lambda = R(u)D_{\mathbf U,\mu}\lambda$ and $D^*_{\mathbf U^u,\mu}R(u)\lambda = R(u)D^*_{\mathbf U,\mu}\lambda$, and consequently each of the operators $\mathcal D$ of Definition~\ref{def:differences} satisfies $\mathcal D_{\mathbf U^u}R(u) = R(u)\mathcal D_{\mathbf U}$.
\end{enumerate}
\end{lemma}

\begin{proof}
(i) Write $W \coloneqq \mathbf U(x, x + ae_\mu)$. By the cyclicity of the trace, $\tr(W\lambda(x + ae_\mu)W^{-1}\lambda'(x)) = \tr(\lambda(x + ae_\mu)R(W^{-1})\lambda'(x))$, and $W^{-1} = \mathbf U(x + ae_\mu, x)$. Summing over $x$ and shifting the variable by $ae_\mu$ in this term, we obtain
\[
    \sum_x\tr\bigl((D_{\mathbf U,\mu}\lambda)(x)\lambda'(x)\bigr) = a^{-1}\sum_x\tr\Bigl(\lambda(x)\bigl(R(\mathbf U(x, x - ae_\mu))\lambda'(x - ae_\mu) - \lambda'(x)\bigr)\Bigr)\;,
\]
which is the first identity. The second is its sum over $\mu$ with $\lambda' = A_\mu$. For the third, the first identity gives
\[
    \langle D_{\mathbf U}A, F\rangle = \sum_{\mu < \nu}\bigl(\langle A_\nu, D^*_{\mathbf U,\mu}F_{\mu\nu}\rangle - \langle A_\mu, D^*_{\mathbf U,\nu}F_{\mu\nu}\rangle\bigr)\;.
\]
In the first sum we exchange the names of $\mu$ and $\nu$, and in the second we write $F_{\mu\nu} = -F_{\nu\mu}$. The result is $\sum_\mu\langle A_\mu, \sum_{\nu > \mu}D^*_{\mathbf U,\nu}F_{\nu\mu} + \sum_{\nu < \mu}D^*_{\mathbf U,\nu}F_{\nu\mu}\rangle = \langle A, D^*_{\mathbf U}F\rangle$.

(ii) Since $R$ is multiplicative, $R(\mathbf U^u(x, x + ae_\mu))R(u(x + ae_\mu)) = R(u(x))R(\mathbf U(x, x + ae_\mu))$ by \eqref{eq:gauge-transformation}, and therefore $(D_{\mathbf U^u,\mu}R(u)\lambda)(x) = a^{-1}R(u(x))(R(\mathbf U(x, x + ae_\mu))\lambda(x + ae_\mu) - \lambda(x))$. The same computation with $-e_\mu$ gives the identity for $D^*_{\mathbf U,\mu}$, and the operators of Definition~\ref{def:differences} are sums of products of these two.
\end{proof}

The part of the action which is linear in a fluctuation of the field is the pairing of the fluctuation with a $1$-form built from the plaquettes, and the regularity conditions of the later sections bound this $1$-form separately. For this reason the terms of the effective actions are regarded as functions of a configuration and of an independent variable in the place of this $1$-form \cite[p.~261]{BalabanRG1}.

\begin{definition}[Curvature, Current, and Pairs]\label{def:current}
Let $\mathbf U$ be a configuration on $T^{(j)}$ and $a \coloneqq a_j$.
\begin{enumerate}[label=(\roman*)]
\item The \emph{curvature} of $\mathbf U$ is the $2$-form $F(\mathbf U)$ with the components $F(\mathbf U)_{\mu\nu}(x) \coloneqq a^{-2}\pi\im\mathbf U(\partial p_{\mu\nu}(x))$, and the \emph{current} of $\mathbf U$ is the $1$-form
\begin{equation}\label{eq:current}
    J(\mathbf U) \coloneqq D^*_{\mathbf U}F(\mathbf U)\;.
\end{equation}
\item A \emph{pair} on a set of bonds is a pair $(\mathbf U, \mathbf J)$ of a configuration $\mathbf U$ and a $1$-form $\mathbf J$ with values in $\mathfrak g^c$, both restricted to the set. A gauge transformation acts on pairs by $(\mathbf U, \mathbf J)^u \coloneqq (\mathbf U^u, R(u)\mathbf J)$.
\end{enumerate}
\end{definition}

The family $F(\mathbf U)$ is a $2$-form because $\partial p_{\nu\mu}(x)$ is the reverse of $\partial p_{\mu\nu}(x)$, so that the two matrices $\im\mathbf U(\partial p)$ differ by the sign. Its values lie in $\mathfrak g^c$, and in $\mathfrak g$ if $\mathbf U$ is a gauge field, by Lemma~\ref{lem:norms}(iv). The current \eqref{eq:current} is \cite[(1.8), p.~261]{BalabanRG1}, and the action on pairs is \cite[(1.10), p.~262]{BalabanRG1}. By Lemma~\ref{lem:action}(ii), Lemma~\ref{lem:norms}(iii), and Lemma~\ref{lem:covariance}(ii),
\begin{equation}\label{eq:current-covariance}
    F(\mathbf U^u) = R(u)F(\mathbf U)\;,
    \qquad
    J(\mathbf U^u) = R(u)J(\mathbf U)\;,
\end{equation}
for every gauge transformation $u$ with values in $G^c$, so that $(\mathbf U, J(\mathbf U))^u = (\mathbf U^u, J(\mathbf U^u))$.

\subsection{Sequences of Domains and the Weighted Distance}\label{sec:sequences}

After several steps the field is an integration variable on lattices of different scales in different regions of the torus, the coarser lattices in the regions in which more steps could be performed. The regions are described by a decreasing sequence of domains, and all bounds on propagators decay in a distance which measures a length in each region in the units of its scale.

\begin{definition}[Sequences of Domains]\label{def:sequence}
Let $M$ be a power of $L$ with $M \leq L^{m_{\mathbb T}}$ and let $R \geq 1$ be an integer. A \emph{sequence of domains with the block size $M$ and the separation $R$} is a sequence of subsets
\begin{equation}\label{eq:sequence}
    \mathbb T = \Omega_0 \supset \Omega_1 \supset \dots \supset \Omega_k\;,
    \qquad
    \Omega_{k+1} \coloneqq \emptyset\;,
\end{equation}
of the torus such that, for $1 \leq j \leq k$, the set $\Omega_j$ is a union of cubes of the grid of side $Ma_j$, and
\begin{equation}\label{eq:separation}
    \dist(\Omega_j^c, \Omega_{j+1}) > RMa_j
\end{equation}
for all $1 \leq j < k$ for which neither $\Omega_j^c$ nor $\Omega_{j+1}$ is empty. Its \emph{layers} are the sets
\begin{equation}\label{eq:layers}
    \Gamma_j \coloneqq \Omega_j^{(j)} \setminus \Omega_{j+1} \subset T^{(j)}\;,
    \qquad
    0 \leq j \leq k\;,
\end{equation}
and its \emph{determining set} is the disjoint union $\mathbf B \coloneqq \bigsqcup_{j=0}^{k}\Gamma_j$. For $y \in \Gamma_j$ we call $j(y) \coloneqq j$ the \emph{scale} of $y$, and we put
\begin{equation}\label{eq:cells}
    \Delta(y) \coloneqq B^j(y)\;,
    \qquad
    \tilde\Delta(y) \coloneqq \{ x \in T_\eta \,|\, \abs{x - y}_\infty < a_j \}\;.
\end{equation}
\end{definition}

The domains may be empty or equal to the torus, and the inclusions need not be strict. The sets $\Omega_j \setminus \Omega_{j+1}$ are pairwise disjoint and cover the torus. Both $\Omega_j$ and $\Omega_{j+1}$ are unions of cubes of the grid of side $a_j$, whose boundaries contain no point of the fine lattice, so that the fine points of $\Omega_j \setminus \Omega_{j+1}$ are those of the blocks $B^j(y)$ with $y \in \Gamma_j$. Hence the blocks $\Delta(y)$, $y \in \mathbf B$, form a partition of the fine lattice. For $x \in T_\eta$ we write $y(x)$ for the point of $\mathbf B$ with $x \in \Delta(y(x))$ and $j(x) \coloneqq j(y(x))$. The \emph{bonds of the layer} $\Gamma_j$ are the bonds of $T^{(j)}$ with an end in $\Omega_j$ and no end in $\Omega_{j+1}$. Both ends of such a bond are points of $\mathbf B$ once $RM \geq L$: an end in $\Omega_j$ lies in $\Gamma_j$, and an end outside $\Omega_j$ is a point of $T^{(j)} \subset T^{(j-1)}$ in $\Omega_{j-1}$, hence a point of $\Gamma_{j-1}$. This is the setting of \cite[(2.1)--(2.4), p.~224]{BalabanPRT2} and \cite[p.~277]{BalabanVAR}, with the sequence starting at $\Omega_0 = \mathbb T$, and with the names of \cite[(2.2), p.~255]{BalabanCRE}, which writes $\Gamma_j$ for the layers and $\mathbf B$ for the determining set. The condition \eqref{eq:separation} is imposed in the $\ell^\infty$-distance, which is the smaller one, so that a sequence in our sense satisfies the condition of the papers cited.

The distance in which the propagators decay is defined in \cite[(2.46), p.~231]{BalabanPRT2} as the smallest length of a contour between two points, the part of the contour in a layer being composed of bonds of the scale of that layer and counted in its units. At the surface between two layers a bond of the finer scale does not end at a site of the coarser one, and so we define the distance by paths in the torus instead of lattice contours. The two properties which the later sections use, the cost of crossing a layer and the summability of the exponentials, are Lemma~\ref{lem:distance} and Lemma~\ref{lem:summation}.

\begin{definition}[The Weighted Distance]\label{def:distance}
Let a sequence of domains be given. For $z \in \mathbb T$ let $j(z)$ be the largest $j$ with $z \in \Omega_j$. The \emph{weighted length} of a Lipschitz path $\gamma \colon [0, 1] \to \mathbb T$ is
\begin{equation}\label{eq:weighted-length}
    \ell(\gamma) \coloneqq \int_0^1\frac{\abs{\dot\gamma(t)}_1}{a_{j(\gamma(t))}}\,\dd t\;,
\end{equation}
and the \emph{weighted distance} $d(z, z')$ of two points of the torus is the infimum of $\ell(\gamma)$ over the Lipschitz paths from $z$ to $z'$.
\end{definition}

The function $z \mapsto j(z)$ is upper semicontinuous, because the domains are closed, so that the integrand of \eqref{eq:weighted-length} is measurable. For a point $x$ of the fine lattice, $j(x)$ is the scale of the block which contains $x$, in agreement with the notation introduced above.

\begin{lemma}[The Weighted Distance]\label{lem:distance}
For every sequence of domains with the block size $M$ and the separation $R$:
\begin{enumerate}[label=(\roman*)]
\item $d$ is a metric on $\mathbb T$, and $\abs{z - z'}_1 \leq d(z, z') \leq \eta^{-1}\abs{z - z'}_1$;
\item if $z \in \Omega_j \setminus \Omega_{j+1}$ and $z' \in \Omega_{j'} \setminus \Omega_{j'+1}$, then $d(z, z') \geq RM\max\{ \abs{j - j'} - 1, 0 \}$.
\end{enumerate}
\end{lemma}

\begin{proof}
(i) The weights satisfy $\eta = a_0 \leq a_{j(z)} \leq a_k = 1$. Hence $\ell(\gamma)$ lies between the $\ell^1$-length of $\gamma$ and $\eta^{-1}$ times that length, and the $\ell^1$-length of a path from $z$ to $z'$ is at least $\abs{z - z'}_1$, with equality for a suitable path. This gives the two inequalities, and with them $d(z, z') > 0$ for $z \neq z'$. The symmetry of $d$ follows by reversing the paths, and the triangle inequality by concatenating them.

(ii) We may assume $j' \geq j + 2$. Let $\gamma$ be a Lipschitz path from $z$ to $z'$ and let $j < i < j'$. Then $\gamma(0) = z \notin \Omega_i$, because $\Omega_i \subset \Omega_{j+1}$, and $\gamma(1) = z' \in \Omega_{i+1}$, because $\Omega_{j'} \subset \Omega_{i+1}$. Let $t_1$ be the smallest $t$ with $\gamma(t) \in \Omega_{i+1}$, which exists since $\Omega_{i+1}$ is closed, and let $t_0$ be the supremum of the $t < t_1$ with $\gamma(t) \notin \Omega_i$. Then $\gamma(t_0)$ lies in the closure of $\Omega_i^c$, and $\gamma(t) \in \Omega_i \setminus \Omega_{i+1}$ for $t_0 < t < t_1$, so that $j(\gamma(t)) = i$ on this interval. The $\ell^1$-length of the arc $\gamma|_{[t_0, t_1]}$ is at least $\abs{\gamma(t_1) - \gamma(t_0)}_\infty$, which is at least the distance of $\Omega_{i+1}$ from the closure of $\Omega_i^c$, and this distance is greater than $RMa_i$ by \eqref{eq:separation}. The weighted length of the arc is therefore greater than $RM$. The intervals $(t_0, t_1)$ which belong to different $i$ are disjoint, since $j(\gamma(t))$ takes different values on them, and there are $j' - j - 1$ of them. Hence $\ell(\gamma) > RM(j' - j - 1)$.
\end{proof}

The second property is the summability of the exponentials of the distance over the determining set, uniformly in the sequence and in the number of the scales. It is \cite[Lemma~2.1, p.~234]{BalabanPRT2}, stated here for an arbitrary rate $\delta$, since the lemma is applied with the decay rates of several operators.

\begin{lemma}[Summation over the Determining Set]\label{lem:summation}
For every $\delta > 0$ and every $0 < \alpha < 1$ there are numbers $r_1(\alpha, \delta) \geq 1$ and $c_1(\alpha, \delta) \geq 1$, which depend on $\alpha$ and $\delta$ only, with the following property. Let $L$, $k$, and $m_{\mathbb T}$ be as in Definition~\ref{def:lattices} and let a sequence of domains with the block size $M$ and the separation $R$ be given, with $RM \geq r_1(\alpha, \delta)$. Then
\begin{equation}\label{eq:summation}
    \sup_{y \in \mathbf B}\sum_{y' \in \mathbf B}e^{-\alpha\delta d(y, y')} \leq c_1(\alpha, \delta)\;,
\end{equation}
and for every $n \geq 1$ and all $y_0, y_n \in \mathbf B$,
\begin{equation}\label{eq:chain-summation}
    \sum_{y_1, \dots, y_{n-1} \in \mathbf B}\prod_{i=1}^{n}e^{-\delta d(y_{i-1}, y_i)} \leq c_1(\alpha, \delta)^{n-1}e^{-(1 - \alpha)\delta d(y_0, y_n)}\;.
\end{equation}
\end{lemma}

\begin{proof}[Proof of \eqref{eq:chain-summation} from \eqref{eq:summation}]
By the triangle inequality of Lemma~\ref{lem:distance}(i), $\sum_{i=1}^{n}d(y_{i-1}, y_i) \geq d(y_0, y_n)$, so that the product in \eqref{eq:chain-summation} is at most $e^{-(1 - \alpha)\delta d(y_0, y_n)}\prod_{i=1}^{n}e^{-\alpha\delta d(y_{i-1}, y_i)}$. We bound the factor with $i = n$ by one and sum over $y_{n-1}, y_{n-2}, \dots, y_1$ in this order, each sum being at most $c_1(\alpha, \delta)$ by \eqref{eq:summation}.
\end{proof}

\subsection{Norms of Fields}\label{sec:norms}

The bounds of the linear theory are stated in supremum norms, in H\"older norms, in norms which weight a field by a power of the spacing of the layer, and in $L^2$-norms. They are the norms of \cite[(1.108), (1.109), p.~35]{BalabanPRT1} and \cite[(3.39)--(3.41), p.~397]{BalabanPBG}.

\begin{definition}[Norms]\label{def:norms}
Let a sequence of domains be given, let $\mathbf U$ be a configuration on $T_\eta$, and let $f$ be a $0$-form on $T_\eta$, or a family of $0$-forms such as a $1$-form or the gradient of one.
\begin{enumerate}[label=(\roman*)]
\item The \emph{supremum norm} $\abs f$ is the supremum of $\abs{f_\iota(x)}$ over the sites $x$ and the components $f_\iota$ of $f$; for a subset $S$ of the torus, $\abs f_S$ is the same supremum over $x \in S^{(0)}$.
\item For $0 \leq \beta \leq 1$ and $0 \leq j \leq k$ the \emph{H\"older norm at the scale $j$} is
\begin{equation}\label{eq:holder}
    \norm f_\beta^{(j)} \coloneqq \sup\Bigl(\frac{\abs{x - x'}_1}{a_j}\Bigr)^{-\beta}\abs{R(\mathbf U(\gamma))f_\iota(x') - f_\iota(x)}\;,
\end{equation}
the supremum over the components, over the pairs of sites $x \neq x'$ of $T_\eta$ with $\abs{x - x'}_1 \leq a_j$, and over the shortest paths $\gamma$ from $x$ to $x'$. We write $\norm f_\beta \coloneqq \norm f_\beta^{(k)}$, and $\norm A_{1,\beta} \coloneqq \norm{\nabla_{\mathbf U}A}_\beta$ for a $1$-form $A$.
\item For a real number $\alpha$ the \emph{weighted norm} is
\begin{equation}\label{eq:weighted-norm}
    \abs f_{(\alpha)} \coloneqq \sup_{x \in T_\eta}a_{j(x)}^{-\alpha}\max_\iota\abs{f_\iota(x)}\;.
\end{equation}
\end{enumerate}
\end{definition}

In the papers cited the H\"older norm $\norm\cdot_\beta$ is taken ``in the $\eta$-scale'', which is our scale $k$, and a superscript $\xi = L^{-j}$ indicates that ``functions and distances are on the $\xi$-lattice'' \cite[p.~397]{BalabanPBG}, which is \eqref{eq:holder}. The norm \eqref{eq:holder} depends on the configuration through the transport along $\gamma$, and on scalar functions, for which we use it as well, it does not.

\subsection{Localisation Domains}\label{sec:localisation-domains}

The terms of the effective actions are indexed by connected unions of large cubes, and they decay exponentially in the linear size of the union. We define these notions for the cubes of an arbitrary grid, together with a variant of the linear size in which a part of the domain is not counted, which is needed when a term has a prescribed support.

\begin{definition}[Unions of Cubes and Their Linear Size]\label{def:localisation}
Let $s$ be the side of a grid in the sense of Definition~\ref{def:lattices}(iii).
\begin{enumerate}[label=(\roman*)]
\item Two cubes of the grid \emph{touch} if they have a point in common, and they \emph{share a wall} if their intersection is a three-dimensional cube. A union $X$ of cubes of the grid is \emph{connected} if any two of its cubes are joined by a chain of cubes of $X$ in which consecutive cubes share a wall. The number of the cubes of $X$ is denoted by $\#(X)$.
\item For a union $X$ of cubes, $X^\sim$ is the union of $X$ and of the cubes of the grid which touch a cube of $X$; $X^{\sim n}$ is the result of $n$ such operations; and $X^{\sim-n} \coloneqq ((X^c)^{\sim n})^c$, where $X^c$ denotes the union of the cubes of the grid which are not in $X$.
\item For a connected union $X$ of cubes and a union $D \subset X$ of cubes, possibly empty, the \emph{linear size of $X$ relative to $D$} is
\begin{equation}\label{eq:contracted-size}
    d(X; D) \coloneqq \frac{1}{s}\inf\abs{G \setminus D}\;,
\end{equation}
the infimum of the one-dimensional measure, with respect to the Euclidean length, over the finite unions $G$ of segments and points with $G \subset X$ such that $G$ meets every cube of $X$ and $G \cup D$ is connected. For a union $X$ of cubes with the connected components $X^i$ we put $d(X; D) \coloneqq \sum_id(X^i; D \cap X^i)$. The \emph{linear size} of $X$ is $d(X) \coloneqq d(X; \emptyset)$.
\item Let $M$ be a power of $L$ with $M \leq L^{m_{\mathbb T}}$. For $0 \leq j \leq k$ the family $\pi_j$ is the grid of side $Ma_j$, a \emph{localisation domain} of the scale $j$ is a connected union of cubes of $\pi_j$, the class of these domains is $\mathbf D_j$, and $d_j(X)$ and $d_j(X; D)$ denote the linear sizes \eqref{eq:contracted-size} for the grid $\pi_j$.
\end{enumerate}
\end{definition}

The domains of $\mathbf D_j$ and the size $d_j(X)$ are those of \cite[p.~257]{BalabanRG1}, where the size is the length of a shortest tree graph in $X$ which meets all its cubes, divided by the side. In the papers the number of the cubes of a domain is bounded by a multiple of its linear size \cite[(1.28), p.~8]{BalabanRG2}, which fails for a single cube, whose linear size is zero. The following count holds for every domain.

\begin{lemma}[The Cubes of a Localisation Domain]\label{lem:cube-count}
Let $X$ be a connected union of cubes of a grid of side $s$. Then
\begin{equation}\label{eq:cube-count}
    \#(X) \leq 2^8(2 + d(X)) \leq 2^9(1 + d(X))\;,
    \qquad
    \diam X \leq s(2d(X) + 4)\;,
\end{equation}
with the diameter in the Euclidean distance. Moreover, for every $\kappa \geq 2^9(6\log 2 + 1)$ and every cube $\square$ of the grid,
\begin{equation}\label{eq:domain-sum}
    \sum_{X \supset \square}e^{-\kappa d(X)} \leq 2^{6 \cdot 2^{10}} + 1\;,
\end{equation}
the sum over the connected unions of cubes of the grid which contain $\square$.
\end{lemma}

\begin{proof}
Let $G$ be a set which is admissible in \eqref{eq:contracted-size} for $D = \emptyset$, that is, a connected finite union of segments in $X$ which meets every cube of $X$, and let $\lambda$ be the sum of the $\ell^\infty$-lengths of its segments, which is at most its Euclidean length. We first cover $G$ by few boxes. Regard $G$ as a connected graph whose arcs are segments, and double every arc. All degrees of the resulting graph are even, so that it has a closed walk which traverses every arc once, and parametrising this walk by the $\ell^\infty$-length we obtain a map $\gamma \colon [0, 2\lambda] \to G$ onto $G$ with $\abs{\gamma(t) - \gamma(t')}_\infty \leq \abs{t - t'}$. Let $n \coloneqq \ceil{\lambda/s}$ and $t_i \coloneqq \min(2is, 2\lambda)$ for $0 \leq i \leq n$. Every $t \in [0, 2\lambda]$ lies within $s$ of some $t_i$, so that every point of $G$ lies within the $\ell^\infty$-distance $s$ of one of the $n + 1 \leq \lambda/s + 2$ points $\gamma(t_i)$.

Next we count the cubes near one of these points. The cubes of the grid are the products of intervals $[c_\nu - s/2, c_\nu + s/2]$ with $c_\nu \in s\mathbb Z$, and such an interval meets $[q_\nu - s, q_\nu + s]$ only if $c_\nu$ lies in an interval of length $3s$, which contains at most four points of $s\mathbb Z$. A box $\{ z \,|\, \abs{z - q}_\infty \leq s \}$ therefore meets at most $4^4 = 2^8$ cubes of the grid. Every cube of $X$ contains a point of $G$ and hence meets one of the $n + 1$ boxes around the points $\gamma(t_i)$, so that $\#(X) \leq 2^8(2 + \lambda/s)$, and the infimum over $G$ gives the first bound of \eqref{eq:cube-count}. The second inequality there is $2 + d \leq 2(1 + d)$.

For the diameter, let $z, z' \in X$. They lie in cubes of $X$ which contain points $q$, $q'$ of $G$, and the Euclidean diameter of a cube of side $s$ in four dimensions is $2s$. Since $G$ is connected, $\abs{q - q'}_\infty \leq \lambda$, and the Euclidean distance is at most twice the $\ell^\infty$-distance. Hence $\abs{z - z'} \leq 4s + 2\lambda$, and the infimum over $G$ gives the bound.

For \eqref{eq:domain-sum}, a connected union of $n$ cubes which contains $\square$ has a spanning tree with respect to the relation of sharing a wall, and a closed walk from $\square$ along this tree which traverses every edge twice has $2(n - 1)$ steps, each of which goes to one of the eight cubes which share a wall with the present one. The union is determined by the walk, so that there are at most $8^{2(n-1)} = 2^{6(n-1)}$ such unions. For $n \geq 2^{10}$ the first bound of \eqref{eq:cube-count} gives $d(X) \geq 2^{-8}n - 2 \geq 2^{-9}n$, and then $e^{-\kappa d(X)} \leq 2^{-6n}e^{-n}$ by the assumption on $\kappa$. The sum is therefore at most $\sum_{n=1}^{2^{10}-1}2^{6(n-1)} + \sum_{n \geq 2^{10}}e^{-n}$, in which the first sum is less than $2^{6 \cdot 2^{10}}$ and the second less than one.
\end{proof}

The cluster expansions of the later sections use five inequalities which relate the linear sizes of a domain, of the parts into which it is decomposed, and of the domain of the next scale which covers it. For $D = \emptyset$ the first is \cite[(2.27), p.~18]{BalabanRG2}, with $3$ in the place of $5$, and the second is \cite[(2.32), p.~18]{BalabanRG2} with the constant $18$ in the place of the printed $4$. With $4$ the inequality fails: if $Z$ consists of one cube and of eight rods of $n \geq 2$ cubes each, one attached to each wall of the cube, and the rods are the sets $Y_i$, then $d_k(Y_i) = n - 2$, while a connected set which meets all the cubes of $Z$ has length at least $(n - 1)Ma_k$ beyond each of the eight walls, so that $d_k(Z) \geq \sum_id_k(Y_i) + 8$. For a union $X$ of cubes of $\pi_k$ we write $X'$ for the smallest union of cubes of the grid of side $LMa_k$ which contains $X^\sim$, as in \cite[p.~13]{BalabanRG2}, and we write $d_k$ and $d'$ for the linear sizes \eqref{eq:contracted-size} of the grids of the sides $Ma_k$ and $LMa_k$.

\begin{lemma}[Five Geometric Inequalities]\label{lem:contracted-geometry}
Let $M$ be a power of $L$ with $LM \leq L^{m_{\mathbb T}}$.
\begin{enumerate}[label=(\roman*)]
\item Let $Y_0$ be a connected union of cubes of $\pi_k$, $\mathbf D$ a finite family of connected unions $Y \subset Y_0$ of cubes, and $D \subset Y_0$ a union of cubes with $n_D \geq 0$ connected components, such that $Y_0 = \bigcup_{Y \in \mathbf D}Y \cup D$. Then $\sum_{Y \in \mathbf D}d_k(Y) \geq d_k(Y_0; D) - 3(\abs{\mathbf D} + n_D - 1)$.
\item Let $Z$ be a connected union of cubes of $\pi_k$, let $Y_1, \dots, Y_n \subset Z$ be connected unions of cubes no two of which have a cube in common or contain two cubes which share a wall, and let $D \subset \bigcup_iY_i$ be a union of cubes. Then $\sum_id_k(Y_i; D \cap Y_i) + 18\,\#(Z \setminus \bigcup_iY_i) \geq d_k(Z; D)$.
\item Let $Z$ be a connected union of cubes of $\pi_k$ and let $D \subset Z$ be a nonempty union of cubes. Then $L\,d'(Z'; D') \leq (1 + 486/L)\,d_k(Z; D)$.
\item Let $Z$ be a connected union of cubes of the grid of side $LMa_k$, let $Z_1, \dots, Z_n \subset Z$ be connected unions of cubes and $D_i \subset Z_i$ unions of cubes, of which $n' \geq 1$ are nonempty and $n_\emptyset$ are empty, and put $Z_0 \coloneqq \bigcup_iZ_i$ and $D \coloneqq \bigcup_iD_i$. Then $d'(Z; D) \leq \sum_id'(Z_i; D_i) + 3\,\#(Z \setminus Z_0) + 3(n_\emptyset + n' - 1)$.
\item Let $X$ be a connected union of cubes of a grid and $D \subset X$ a nonempty union of cubes. Then $\#(X) \leq 81\,\#(D) + 324\,d(X; D)$.
\end{enumerate}
\end{lemma}

\section{The Averaging Operations}\label{sec:averaging}

A renormalisation step replaces a gauge field on a lattice by its average on the lattice whose spacing is $L$ times larger. Two operations are needed. The first is the average itself, a map between gauge fields which commutes with the gauge transformations and with the symmetries of the lattice. The second arises when the field is written as a fluctuation times a background: the average of the product is then the average of the background, multiplied by a function of the fluctuation whose linear part enters every quadratic form of the later sections. We define both, and we state for each what the later sections use: on which bonds it depends, how it transforms, where it is analytic, and how large it is.

\subsection{The Chart of the Group}\label{sec:chart}

The averages are exponentials of means of logarithms, so that we first have to say which logarithm is meant. A Hermitian logarithm of an element of $G$ need not lie in $\mathfrak g$, and for this reason the logarithm on the group is defined as the inverse of the exponential on a ball of the algebra, and not through the spectrum.

\begin{definition}[The Chart]\label{def:chart}
Let $\varrho_G \coloneqq \pi/\sqrt N$. The \emph{chart domain} is $N_G \coloneqq \{ e^{iA} \,|\, A \in \mathfrak g,\ h(A) < \varrho_G \}$, and the \emph{chart} $\log \colon N_G \to i\mathfrak g$ is defined by $\log e^{iA} \coloneqq iA$ for $A \in \mathfrak g$ with $h(A) < \varrho_G$. For a matrix $W$ with $\norm{W - \id} < 1$ we put $\operatorname{Log}W \coloneqq \sum_{n \geq 1}(-1)^{n+1}n^{-1}(W - \id)^n$.
\end{definition}

\begin{lemma}[The Chart]\label{lem:chart}
\begin{enumerate}[label=(\roman*)]
\item The map $A \mapsto e^{iA}$ is injective on $\{ A \in \mathfrak g \,|\, h(A) < \varrho_G \}$, so that the chart is well defined, and $\norm A < \pi$ there. The chart domain is invariant under $W \mapsto W^{-1}$ and under conjugation by elements of $G$, with $\log W^{-1} = -\log W$ and $\log(uWu^{-1}) = R(u)\log W$.
\item There is a number $r_G \in (0, 1)$, which depends on $G$ only, such that every $W \in G$ with $\norm{W - \id} \leq r_G$ lies in $N_G$ and satisfies $\log W = \operatorname{Log}W$.
\item For a Hermitian matrix $X$ with $\norm X \leq \pi$ and each of the three norms $\mathcal N$ of Lemma~\ref{lem:norms}, $\mathcal N(e^{iX} - \id) \leq \mathcal N(X) \leq \frac{\pi}{2}\mathcal N(e^{iX} - \id)$.
\item The function $\operatorname{Log}$ is holomorphic on $\{ \norm{W - \id} < 1 \}$, $e^{\operatorname{Log}W} = W$, $\operatorname{Log}(\mathbf uW\mathbf u^{-1}) = R(\mathbf u)\operatorname{Log}W$ for every invertible $\mathbf u$ with $\norm{\mathbf uW\mathbf u^{-1} - \id} < 1$, and $\operatorname{Log}W \in i\mathfrak g^c$ if $W = e^{iX}$ with $X \in \mathfrak g^c$ and $\norm X < \log 2$.
\end{enumerate}
\end{lemma}

\subsection{The Average of One Step}\label{sec:average}

We fix a scale $1 \leq j \leq k$ and write $a \coloneqq a_{j-1}$ for the spacing of the finer lattice $T^{(j-1)}$. The average of a bond $c$ of $T^{(j)}$ is built from the fine paths which join the block of $c_-$ to the block of $c_+$ \cite[p.~252, (0.4), p.~253]{BalabanRG1}.

\begin{definition}[The Paths of a Coarse Bond]\label{def:block-paths}
\begin{enumerate}[label=(\roman*)]
\item For $y \in T^{(j)}$ and $x \in B(y)$, write $x - y = a\sum_\mu n_\mu e_\mu$ with $\abs{n_\mu} \leq \frac{1}{2}(L - 1)$. The family $\mathbf G(y, x)$ consists of the paths of $T^{(j-1)}$ from $y$ to $x$ which traverse the coordinates $\mu$ with $n_\mu \neq 0$ one after the other, $\abs{n_\mu}$ bonds in the direction of the sign of $n_\mu$ for each, one path for every ordering of these coordinates. For $x = y$ it consists of the empty path.
\item For a bond $c$ of $T^{(j)}$ and $x \in B(c_-)$ let $x' \coloneqq x + (c_+ - c_-) \in B(c_+)$, and let $[x, x']$ be the straight path of $L$ bonds of $T^{(j-1)}$ from $x$ to $x'$; in particular $[c] \coloneqq [c_-, c_+]$.
\item For $x \in B(c_-)$, $\Gamma \in \mathbf G(c_-, x)$, and $\Gamma' \in \mathbf G(c_+, x')$, the \emph{loop} $\ell_c(x, \Gamma, \Gamma')$ is the closed path which consists of $\Gamma$, then $[x, x']$, then $\Gamma'^{-1}$, then $[c]^{-1}$. It begins and ends at $c_-$.
\end{enumerate}
\end{definition}

All these paths run through sites of $B(c_-) \cup B(c_+)$.

\begin{definition}[The Average]\label{def:average}
Let $U$ be a gauge field on $T^{(j-1)}$ and $c$ a bond of $T^{(j)}$. Let $\mathcal O_c$ be the set of the gauge fields for which $U(\ell_c(x, \Gamma, \Gamma')) \in N_G$ for all loops of $c$, and for $U \in \mathcal O_c$ let
\begin{equation}\label{eq:average-exponent}
    \beta_c(U) \coloneqq \frac{1}{L^4}\sum_{x \in B(c_-)}\frac{1}{\abs{\mathbf G(c_-, x)}\abs{\mathbf G(c_+, x')}}\sum_{\Gamma, \Gamma'}\frac{1}{i}\log U(\ell_c(x, \Gamma, \Gamma'))\;.
\end{equation}
The \emph{average} of $U$ is the gauge field $M(U)$ on $T^{(j)}$ with
\begin{equation}\label{eq:average}
    M(U)(c) \coloneqq
    \begin{cases}
        e^{i\beta_c(U)}U([c]) & \text{if } U \in \mathcal O_c\;,\\
        U([c]) & \text{otherwise}\;.
    \end{cases}
\end{equation}
The \emph{average of order $n$} is the $n$-fold iterate $M^n$, which maps gauge fields on $T^{(j)}$ to gauge fields on $T^{(j+n)}$.
\end{definition}

On $\mathcal O_c$ this is \cite[(0.4), p.~253]{BalabanRG1}: the straight transport along $c$ is multiplied by the exponential of a mean, with weights which are positive and sum to one, of the logarithms of the holonomies around the loops of $c$. The second line of \eqref{eq:average} makes the average a map on all gauge fields, which the definition of the renormalisation transformation requires, and it plays no role in any estimate, since every estimate concerns fields in the set of the following lemma. The iterates are those of \cite[(43), p.~24]{BalabanAVG}.

\begin{lemma}[Properties of the Average]\label{lem:average}
\begin{enumerate}[label=(\roman*)]
\item \emph{Consistency.} $\mathcal O_{c^{-1}} = \mathcal O_c$ and $M(U)(c^{-1}) = M(U)(c)^{-1}$, so that $M(U)$ is a gauge field, and $\beta_c(U) \in \mathfrak g$.
\item \emph{Locality.} $M(U)(c)$ depends only on the values of $U$ on the bonds of $T^{(j-1)}$ with both ends in $B(c_-) \cup B(c_+)$.
\item \emph{Covariance.} For every gauge transformation $u \colon T^{(j-1)} \to G$, the set $\mathcal O_c$ is invariant and $M(U^u)(c) = u(c_-)M(U)(c)u(c_+)^{-1}$.
\item \emph{Symmetry.} Let $\theta$ be a translation of the torus by a point of $T^{(j)}$, a reflection $x_\mu \mapsto -x_\mu$, or a permutation of the coordinates, and let $(\theta U)(b) \coloneqq U(\theta^{-1}b)$. Then $M(\theta U) = \theta M(U)$.
\item \emph{The domain.} There are numbers $\alpha_M > 0$ and $C_M$, which depend on $L$ and $G$ only, such that if $\abs{U(\partial p) - \id} \leq \alpha \leq \alpha_M$ for all plaquettes $p$ of $T^{(j-1)}$ whose corners lie in $B(c_-) \cup B(c_+)$, then $U \in \mathcal O_c$, every loop of $c$ satisfies $\norm{U(\ell) - \id} \leq r_G$, and $\abs{\beta_c(U)} \leq C_M\alpha$.
\item \emph{Analyticity.} There is a number $r_M > 0$, which depends on $L$ and $G$ only, with the following property. Let $U$ be as in \emph{(v)}, and let $A$ be a $1$-form with values in $\mathfrak g^c$ on the bonds of $B(c_-) \cup B(c_+)$ with $\abs{aA} < r_M$. Then \eqref{eq:average-exponent} and the first line of \eqref{eq:average}, with $\operatorname{Log}$ in the place of $\log$, define a holomorphic function $A \mapsto M(e^{iaA}U)(c)$ with values in $G^c$, which is the average of the gauge field $e^{iaA}U$ when $A$ has values in $\mathfrak g$, and which has the properties \emph{(i)}--\emph{(iv)}, the covariance for gauge transformations with values in $G^c$ which map the configuration into the domain.
\end{enumerate}
\end{lemma}

By (vi) the average, and with it every iterate, is defined for the configurations $e^{iaA}U$ of that statement, and we use the notation $M(\mathbf U)$ for them.

The first quantitative property is that the average of a field with small plaquettes has small plaquettes, in the units of the coarser lattice, and that this is a local statement. It is \cite[Propositions~1, 2, pp.~26]{BalabanAVG}, in the form which the test functions of the construction need, with an arbitrary bound in the place of a fixed threshold.

\begin{proposition}[Averages of Regular Fields]\label{prop:regular-averages}
There are numbers $C_0 \geq 1$ and $c_2 \in (0, \alpha_M]$, which depend on $L$ and $G$ only, with the following properties.
\begin{enumerate}[label=(\roman*)]
\item Let $U$ be a gauge field on $T^{(j-1)}$ and $p'$ a plaquette of $T^{(j)}$ with the first corner $y_0$, and suppose that $\abs{U(\partial p) - \id} \leq \alpha \leq c_2L^{-2}$ for all plaquettes $p$ of $T^{(j-1)}$ whose corners lie in the blocks of the four corners of $p'$. Then
\begin{equation}\label{eq:averaged-plaquette}
    \abs{M(U)(\partial p') - \id} \leq \frac{1}{L^4}\sum_{x \in B(y_0)}\sum_{p \subset p'_x}\abs{U(\partial p) - \id} + C_0L^4\alpha^2 \leq L^2\alpha + C_0L^4\alpha^2\;,
\end{equation}
where $p'_x$ is the translate of $p'$ by $x - y_0$ and the inner sum runs over the $L^2$ plaquettes of $T^{(j-1)}$ which it contains.
\item Let $0 \leq i < j \leq k$, let $U$ be a gauge field on $T^{(i)}$ and $p'$ a plaquette of $T^{(j)}$, and suppose that $\abs{U(\partial p) - \id} \leq \alpha(a_i/a_j)^2$ with $\alpha \leq c_2$ for all plaquettes $p$ of $T^{(i)}$ whose corners lie in the blocks of order $j - i$ of the four corners of $p'$. Then $\abs{M^{j-i}(U)(\partial p') - \id} \leq \alpha + 2C_0\alpha^2 \leq 2\alpha$.
\end{enumerate}
\end{proposition}

\subsection{The Average Relative to a Background}\label{sec:relative-average}

Let the field be $e^{iaA}\mathbf U$, with a background $\mathbf U$ and a fluctuation $A$. We want to write its average as a function of the fluctuation times the average of the background. The quotient $M(e^{iaA}\mathbf U)(c)M(\mathbf U)(c)^{-1}$ does this, but it is not the function which the linear theory uses: under a gauge transformation of the fluctuation which is trivial at the coarse sites it changes, through the paths $\Gamma$ and $\Gamma'$ which join the coarse sites to the points of their blocks. The function used instead is obtained by first transforming the fluctuation, at each coarse site, to the gauge in which the mean of its transports to the points of the block is trivial \cite[pp.~26--31]{BalabanAVG}, \cite[p.~418]{BalabanPBG}. The transports are taken along all the paths of $\mathbf G(y, x)$ at once, through a mean of group elements \cite[(0.5)--(0.11), p.~253]{BalabanRG1}.

\begin{definition}[The Mean of Group Elements and the Transports of a Block]\label{def:group-mean}
\begin{enumerate}[label=(\roman*)]
\item Let $\mathbf W_1, \dots, \mathbf W_n \in G^c$ with $\norm{\mathbf W_i\mathbf W_1^{-1} - \id} \leq \frac{1}{4}r_G$ for all $i$. Their \emph{mean} $\mathsf m(\{ \mathbf W_i \})$ is the element $\mathbf W \in G^c$ with $\norm{\mathbf W\mathbf W_1^{-1} - \id} \leq \frac{1}{2}r_G$ and
\begin{equation}\label{eq:group-mean}
    \sum_{i=1}^{n}\operatorname{Log}(\mathbf W_i\mathbf W^{-1}) = 0\;.
\end{equation}
\item For a configuration $\mathbf U$ on $T^{(j-1)}$, a site $y \in T^{(j)}$, and $x \in B(y)$, the \emph{transport} from $y$ to $x$ is $\mathbf U(y, x) \coloneqq \mathsf m(\{ \mathbf U(\Gamma) \}_{\Gamma \in \mathbf G(y, x)})$, whenever the mean is defined.
\end{enumerate}
\end{definition}

\begin{lemma}[The Mean of Group Elements]\label{lem:group-mean}
Under the condition of Definition~\ref{def:group-mean}(i) the equation \eqref{eq:group-mean} has exactly one solution with $\norm{\mathbf W\mathbf W_1^{-1} - \id} \leq \frac{1}{2}r_G$. It does not depend on the order of the elements, it is a holomorphic function of them, and
\[
    \mathsf m(\{ \mathbf u\mathbf W_i\mathbf v \}) = \mathbf u\,\mathsf m(\{ \mathbf W_i \})\,\mathbf v\;,
    \qquad
    \mathsf m(\{ \mathbf W_i^{-1} \}) = \mathsf m(\{ \mathbf W_i \})^{-1}\;,
\]
for $\mathbf u, \mathbf v \in G^c$ for which both sides are defined. If all $\mathbf W_i$ lie in $G$, then so does the mean. If $\mathbf W_i = e^{iX_i}$ with $X_i \in \mathfrak g^c$ and $\norm{X_i} \leq \varrho$ for a sufficiently small $\varrho$, which depends on $G$ only, then $\norm{\frac{1}{i}\operatorname{Log}\mathsf m(\{ \mathbf W_i \}) - \frac{1}{n}\sum_iX_i} \leq C_{\mathsf m}\varrho^2$ with a constant $C_{\mathsf m}$ which depends on $G$ only.
\end{lemma}

These are the properties \cite[(0.5)--(0.9), p.~253]{BalabanRG1} of the mean \cite[(0.10)]{BalabanRG1}. By the covariance, $\mathbf U^u(y, x) = u(y)\mathbf U(y, x)u(x)^{-1}$.

\begin{definition}[The Relative Average]\label{def:relative-average}
Let $\mathbf U$ be a configuration on $T^{(j-1)}$ and $A$ a $1$-form on $T^{(j-1)}$ with values in $\mathfrak g^c$, and write $\mathbf U_A \coloneqq e^{iaA}\mathbf U$ for the configuration with $\mathbf U_A(b) = e^{iaA(b)}\mathbf U(b)$ on the positive bonds.
\begin{enumerate}[label=(\roman*)]
\item The \emph{site factor} of the fluctuation at $y \in T^{(j)}$ is
\begin{equation}\label{eq:site-factor}
    v_{\mathbf U}(A)(y) \coloneqq \exp\Bigl(\frac{1}{L^4}\sum_{x \in B(y)}\operatorname{Log}\bigl(\mathbf U_A(y, x)\mathbf U(y, x)^{-1}\bigr)\Bigr)\;.
\end{equation}
\item The \emph{relative average} of the fluctuation on a bond $c$ of $T^{(j)}$ is
\begin{equation}\label{eq:relative-average}
    \mathcal M_{\mathbf U}(A)(c) \coloneqq v_{\mathbf U}(A)(c_-)^{-1}\,M(\mathbf U_A)(c)\,v_{\mathbf U}(A)(c_+)\,M(\mathbf U)(c)^{-1}\;,
\end{equation}
so that
\begin{equation}\label{eq:average-of-product}
    M(e^{iaA}\mathbf U)(c) = v_{\mathbf U}(A)(c_-)\,\mathcal M_{\mathbf U}(A)(c)\,M(\mathbf U)(c)\,v_{\mathbf U}(A)(c_+)^{-1}\;.
\end{equation}
\item The \emph{linear part} $Q(\mathbf U)$ and the \emph{remainder} $C(\mathbf U; \cdot)$ of the relative average are defined by
\begin{equation}\label{eq:linear-part}
    \frac{1}{i}\operatorname{Log}\mathcal M_{\mathbf U}(A)(c) = La\bigl(Q(\mathbf U)A\bigr)(c) + C(\mathbf U; aA)(c)\;,
\end{equation}
where $Q(\mathbf U)$ is linear in $A$ and $C(\mathbf U; \cdot)$ vanishes to the second order at zero. The \emph{linear part of the average} is the linear map $Q^{\mathrm{pl}}(\mathbf U)$ with $\frac{1}{i}\operatorname{Log}(M(e^{iaA}\mathbf U)(c)M(\mathbf U)(c)^{-1}) = La(Q^{\mathrm{pl}}(\mathbf U)A)(c) + O(A^2)$.
\item For a $0$-form $\lambda$ on $T^{(j-1)}$, the \emph{block average} is the $0$-form on $T^{(j)}$
\begin{equation}\label{eq:block-average}
    \bigl(Q'(\mathbf U)\lambda\bigr)(y) \coloneqq \frac{1}{L^4}\sum_{x \in B(y)}R(\mathbf U(y, x))\lambda(x)\;.
\end{equation}
\end{enumerate}
\end{definition}

The identity \eqref{eq:average-of-product} is \cite[(3.113), p.~418]{BalabanPBG} with the transports of Definition~\ref{def:group-mean}, and $Q'$ is \cite[(3.19), p.~393]{BalabanPBG}. The factor $La$ in \eqref{eq:linear-part} is the spacing of the coarse lattice, so that $Q(\mathbf U)$ maps $1$-forms to $1$-forms without a change of their dimension.

\begin{proposition}[The Relative Average of One Step]\label{prop:relative-average}
There are numbers $c_4 \in (0, \alpha_M]$, $C_1$, $C_2$, and $C_3$, which depend on $L$ and $G$ only, with the following properties. Let $U$ be a gauge field on $T^{(j-1)}$ with $\abs{U(\partial p) - \id} \leq \alpha_0 \leq c_4$ on the plaquettes with corners in $B(c_-) \cup B(c_+)$, let $A'$ and $A$ be $1$-forms with values in $\mathfrak g^c$ on the bonds of $B(c_-) \cup B(c_+)$ with $\abs{aA'} \leq \alpha_1 \leq c_4$ and $\abs{aA} < c_4$, and let $\mathbf U \coloneqq e^{iaA'}U$.
\begin{enumerate}[label=(\roman*)]
\item \emph{Analyticity and locality.} The transports, the site factors at $c_\pm$, and $\mathcal M_{\mathbf U}(A)(c)$ are defined, they are holomorphic functions of $(A', A)$, and they depend only on the values of $U$, $A'$, and $A$ on the bonds with both ends in $B(c_-) \cup B(c_+)$, the site factor at $y$ only on those in $B(y)$. For $A' = 0$ and $A$ with values in $\mathfrak g$ their values lie in $G$.
\item \emph{Covariance.} For a gauge transformation $u$ of $T^{(j-1)}$ with values in $G$,
\[
    v_{\mathbf U^u}(R(u)A)(y) = u(y)v_{\mathbf U}(A)(y)u(y)^{-1}\;,
    \qquad
    \mathcal M_{\mathbf U^u}(R(u)A)(c) = u(c_-)\mathcal M_{\mathbf U}(A)(c)u(c_-)^{-1}\;,
\]
so that $Q(\mathbf U^u)R(u) = R(u)Q(\mathbf U)$ and $Q'(\mathbf U^u)R(u) = R(u)Q'(\mathbf U)$, with $R(u)$ on the right sides taken at the sites of $T^{(j)}$. The same holds for $u$ with values in $G^c$ near $G$, as long as $\mathbf U^u$ is of the form assumed.
\item \emph{The two linear parts.} Let $\Lambda(\mathbf U)A$ be the $0$-form on $T^{(j)}$ which is the linear part in $A$ of $\frac{1}{i}\operatorname{Log}v_{\mathbf U}(A)$. Then $Q^{\mathrm{pl}}(\mathbf U) = Q(\mathbf U) - D_{M(\mathbf U)}\Lambda(\mathbf U)$, with the covariant difference of the lattice $T^{(j)}$ at the configuration $M(\mathbf U)$.
\item \emph{Gauge directions.} For a $0$-form $\lambda$ on $T^{(j-1)}$ with values in $\mathfrak g^c$,
\begin{gather*}
    Q(\mathbf U)D_{\mathbf U}\lambda = D_{M(\mathbf U)}Q'(\mathbf U)\lambda\;,
    \qquad
    \Lambda(\mathbf U)D_{\mathbf U}\lambda = Q'(\mathbf U)\lambda - \lambda|_{T^{(j)}}\;,\\
    Q^{\mathrm{pl}}(\mathbf U)D_{\mathbf U}\lambda = D_{M(\mathbf U)}\bigl(\lambda|_{T^{(j)}}\bigr)\;.
\end{gather*}
\item \emph{The form of the linear part.} For a bond $b$ of a line $[x, x']$ and $\Gamma \in \mathbf G(c_-, x)$ let $\Gamma_b$ be the path which consists of $\Gamma$ and of the part of the line from $x$ to $b_-$, and let
\begin{equation}\label{eq:straight-average}
    \bigl(Q_0(\mathbf U)A\bigr)(c) \coloneqq \frac{1}{L^5}\sum_{x \in B(c_-)}\sum_{b \in [x, x']}\frac{1}{\abs{\mathbf G(c_-, x)}}\sum_{\Gamma}R(\mathbf U(\Gamma_b))A(b)\;,
\end{equation}
the mean of the values of $A$ on the lines of $c$, transported to $c_-$. Then $Q(\id) = Q_0(\id)$, and $\abs{((Q(\mathbf U) - Q_0(\mathbf U))A)(c)} \leq C_1(\alpha_0 + \alpha_1)\abs A$, the supremum of $A$ being taken over the bonds with both ends in $B(c_-) \cup B(c_+)$.
\item \emph{The remainder.} $\abs{C(\mathbf U; aA)(c)} \leq C_2\abs{aA}^2$, and the derivative of $A \mapsto C(\mathbf U; aA)(c)$ in a direction $\delta A$ is at most $C_3\abs{aA}\abs{a\,\delta A}$, the norms of $A$ and $\delta A$ being suprema over the bonds of $B(c_-) \cup B(c_+)$.
\item \emph{Dependence on the background.} The maps $Q(\mathbf U)$, $Q'(\mathbf U)$, and $C(\mathbf U; \cdot)$ are holomorphic in $A'$, and $\norm{Q(e^{iaA'_1}U) - Q(e^{iaA'_2}U)} \leq C_1\abs{a(A'_1 - A'_2)}$ in the norms of \emph{(v)}, and likewise for $Q'$.
\item \emph{One function for all scales.} Identify the bonds with both ends in $B(c_-) \cup B(c_+)$ with those of a reference bond by the translation and the symmetry of Lemma~\ref{lem:average}\emph{(iv)} which map $c$ to it. Then $\mathcal M_{\mathbf U}(A)(c)$, $v_{\mathbf U}(A)(c_\pm)$, and $M(\mathbf U)(c)$ are the values of functions of the matrices $\mathbf U(b)$ and $aA(b)$ on these bonds which do not depend on $c$, on the scale $j$, on $k$, or on the torus.
\end{enumerate}
\end{proposition}

The parts (v)--(vii) are \cite[Propositions~3--7, (121)--(125)]{BalabanAVG} for the relative average of this section, and (iii), (iv) are the relations \cite[(3.114), (3.115), p.~418]{BalabanPBG}; the last identity of (iv) is the derivative of the covariance of the average, Lemma~\ref{lem:average}(iii). At the trivial configuration the operator \eqref{eq:straight-average} is the average over the straight lines of \cite[(1.18), p.~20]{BalabanPRT1}, with centred blocks. It does not increase the $\ell^1$-, the $\ell^2$-, and the supremum norm of a $1$-form, each taken with the weight of its lattice, because every bond of $T^{(j-1)}$ lies on exactly $L$ of the lines $[x, x']$ of the bonds $c$ of one direction.

The averages of higher order are compositions, and so are their relative versions: the background of the step $l + 1$ is the average of order $l$ of the background \cite[(90), (91), p.~31]{BalabanAVG}, \cite[(3.15), p.~393]{BalabanPBG}.

\begin{definition}[Relative Averages of Higher Order]\label{def:relative-average-order}
Let $0 \leq i < j \leq k$, $n \coloneqq j - i$, let $\mathbf U$ be a configuration on $T^{(i)}$ and $A$ a $1$-form on $T^{(i)}$ with values in $\mathfrak g^c$. Put $\mathbf U^{(l)} \coloneqq M^l(\mathbf U)$, $A^{(0)} \coloneqq A$, and, for $0 \leq l < n$,
\[
    a_{i+l+1}A^{(l+1)} \coloneqq \frac{1}{i}\operatorname{Log}\mathcal M_{\mathbf U^{(l)}}(A^{(l)})\;.
\]
The \emph{linear part} and the \emph{remainder of order $n$} are defined by $A^{(n)} = Q_n(\mathbf U)A + a_j^{-1}C_n(\mathbf U; a_jA)$ with $Q_n(\mathbf U)$ linear and $C_n(\mathbf U; \cdot)$ vanishing to the second order, so that
\begin{equation}\label{eq:composite-linear-part}
    Q_n(\mathbf U) = Q(\mathbf U^{(n-1)})\cdots Q(\mathbf U^{(1)})Q(\mathbf U^{(0)})\;,
\end{equation}
and $Q'_n(\mathbf U) \coloneqq Q'(\mathbf U^{(n-1)})\cdots Q'(\mathbf U^{(0)})$.
\end{definition}

\begin{proposition}[Relative Averages of Higher Order]\label{prop:relative-average-order}
There are numbers $c'_4 \in (0, c_4]$, $C'_2$, and $C'_3$, which depend on $L$ and $G$ only, with the following properties. Let $0 \leq i < j \leq k$ and $n = j - i$, let $c$ be a bond of $T^{(j)}$, and let $\mathsf B_c$ be the set of the bonds of $T^{(i)}$ with both ends in $B^n(c_-) \cup B^n(c_+)$. Let $U$ be a gauge field on $T^{(i)}$ with $\abs{U(\partial p) - \id} \leq \alpha_0(a_i/a_j)^2$ on the plaquettes with corners in $B^n(c_-) \cup B^n(c_+)$, let $A'$ and $A$ be $1$-forms with values in $\mathfrak g^c$ on $\mathsf B_c$ with $\abs{a_jA'} \leq \alpha_1$ and $\abs{a_jA} < c'_4$, where $\alpha_0, \alpha_1 \leq c'_4$, and let $\mathbf U \coloneqq e^{ia_iA'}U$.
\begin{enumerate}[label=(\roman*)]
\item The averages $M^l(e^{ia_iA}\mathbf U)$, $l \leq n$, on the bonds between the blocks concerned, and the forms $A^{(l)}$ are defined, they are holomorphic in $(A', A)$, and they depend only on $U$, $A'$, and $A$ on $\mathsf B_c$. There are gauge transformations $v_n(A)$ of $T^{(j)}$, holomorphic and local in the same sense, with $v_n(0) = \id$, such that
\begin{equation}\label{eq:average-of-product-order}
\begin{aligned}
    M^n(e^{ia_iA}\mathbf U)(c) &= v_n(A)(c_-)\,e^{iX_n(c)}\,M^n(\mathbf U)(c)\,v_n(A)(c_+)^{-1}\;,\\
    X_n(c) &\coloneqq a_j\bigl(Q_n(\mathbf U)A\bigr)(c) + C_n(\mathbf U; a_jA)(c)\;.
\end{aligned}
\end{equation}
\item $\abs{(Q_n(\mathbf U)A)(c)} \leq (1 + C'_2(\alpha_0 + \alpha_1))\abs A_{\mathsf B_c}$ and $\abs{(Q'_n(\mathbf U)\lambda)(y)} \leq (1 + C'_2(\alpha_0 + \alpha_1))\abs\lambda_{B^n(y)}$.
\item $\abs{C_n(\mathbf U; a_jA)(c)} \leq C'_2\abs{a_jA}^2$, and its derivative in a direction $\delta A$ is at most $C'_3\abs{a_jA}\abs{a_j\,\delta A}$, with the suprema over $\mathsf B_c$.
\item $\abs{\frac{1}{i}\operatorname{Log}v_n(A)(y)} \leq C'_2\abs{a_jA}_{B^n(y)}$.
\item The covariance of Proposition~\ref{prop:relative-average}\emph{(ii)} and the relation $Q_n(\mathbf U)D_{\mathbf U} = D_{M^n(\mathbf U)}Q'_n(\mathbf U)$ hold, and $Q_n$, $Q'_n$, and $C_n$ are Lipschitz continuous in $a_jA'$ with a constant $C'_3$ in the norms of \emph{(ii)} and \emph{(iii)}.
\item $\abs{M^n(e^{ia_iA}\mathbf U)(c)M^n(\mathbf U)(c)^{-1} - \id} \leq C'_3\abs{a_jA}_{\mathsf B_c}$.
\end{enumerate}
The constants do not depend on $n$.
\end{proposition}

This is \cite[Propositions~4--7, pp.~38--43]{BalabanAVG} with \cite[(88), (92), (97), (99), pp.~31--32]{BalabanAVG} for the factors $v_n$. The point of the statement is the uniformity in the order $n$. The plaquettes of the averages $M^l(U)$ and the fields $A^{(l)}$ are measured in the units of their scales, in which the hypotheses reproduce themselves from one step to the next by Proposition~\ref{prop:regular-averages} and by (iii), and the complex part of the background is carried along as a second fluctuation: by \eqref{eq:average-of-product-order} for the pair $(U, A')$, each $M^l(\mathbf U)$ is a gauge transform of $e^{ia_{i+l}A'_l}M^l(U)$ with $\abs{a_{i+l}A'_l}$ of the order $\alpha_1a_{i+l}/a_j$. No bound on the differences of $A'$ is assumed.

The gauge transformation which takes a field to the axial gauge of a background \cite[(106), p.~33]{BalabanAVG} and the averages of gauge transformations \cite[Sect.~F]{BalabanAVG} are built from the site factors $v_n$. They are stated in Section~5 together with the gauge conditions, since their statements refer to a sequence of domains.

\section{The Trivial Background}\label{sec:trivial-background}

At the trivial configuration and for the sequence of domains with a single layer, $\Omega_j = \mathbb T$ for all $j \leq k$, the operators of the linear theory commute with the translations of the coarse lattice and are computed in momentum space \cite[Sects.~B--D]{BalabanPRT1}. The main paper uses these formulae for the coefficients of the coupling flow at zero coupling. In this section $\mathbf U = \id$, the covariant differences are the ordinary ones, written $\partial_\mu$, $\partial$, and $\partial^*$, and the operators act on each component of a form with respect to a basis of $\mathfrak g$ separately, so that we may regard the forms as scalar.

\begin{definition}[Momenta and Symbols]\label{def:symbols}
\begin{enumerate}[label=(\roman*)]
\item The \emph{dual tori} are $\tilde T_\eta \coloneqq \{ p \in (2\pi/\ell_{\mathbb T})\mathbb Z^4 \,|\, -\pi/\eta \leq p_\mu < \pi/\eta \}$ and $\tilde T^{(k)} \coloneqq \{ p' \in (2\pi/\ell_{\mathbb T})\mathbb Z^4 \,|\, -\pi \leq p'_\mu < \pi \}$. For a function $f$ on $T_\eta$, $\tilde f(p) \coloneqq \sum_x\eta^4e^{-ip \cdot x}f(x)$ for $p \in \tilde T_\eta$, and likewise with the weight one for functions on $T^{(k)}$ and $p' \in \tilde T^{(k)}$. Every $p \in \tilde T_\eta$ is uniquely $p' + l$ with $p' \in \tilde T^{(k)}$ and $l \in 2\pi\mathbb Z^4$, $l_\mu = 2\pi m_\mu$ with $\abs{m_\mu} \leq \frac{1}{2}(L^k - 1)$. We write $\int\dd p'$ for $\ell_{\mathbb T}^{-4}(2\pi)^4\sum_{p' \in \tilde T^{(k)}}$.
\item The \emph{symbols} are
\begin{gather*}
    \partial_\mu(p) \coloneqq \frac{e^{i\eta p_\mu} - 1}{\eta}\;,
    \qquad
    \Delta(p) \coloneqq \sum_\mu\abs{\partial_\mu(p)}^2\;,\\
    \partial^1_\mu(p') \coloneqq e^{ip'_\mu} - 1\;,
    \qquad
    \Delta_0(p') \coloneqq \sum_\mu\abs{\partial^1_\mu(p')}^2\;,\\
    u_k(p) \coloneqq \prod_\mu\frac{\sin(\frac{1}{2}p_\mu)}{L^k\sin(\frac{1}{2}\eta p_\mu)}\;,
    \qquad
    v_\mu(p) \coloneqq \frac{e^{ip_\mu} - 1}{L^k(e^{i\eta p_\mu} - 1)}\;,
\end{gather*}
where a quotient with vanishing denominator has the value one.
\end{enumerate}
\end{definition}

\begin{proposition}[The Averages and the Minimal Form in Momentum Space]\label{prop:momentum-space}
Let $Q_k \coloneqq Q_k(\id)$ and $Q'_k \coloneqq Q'_k(\id)$ be the averages of order $k$ of Definition~\ref{def:relative-average-order} at the trivial configuration, from forms on $T_\eta$ to forms on $T^{(k)}$.
\begin{enumerate}[label=(\roman*)]
\item $(Q'_k\lambda)(y)$ is the mean of $\lambda$ over $B^k(y)$, $(Q_kA)_\mu(y)$ is the mean over $x \in B^k(y)$ of the mean of $A_\mu$ over the $L^k$ bonds of the straight path from $x$ to $x + e_\mu$, and
\begin{equation}\label{eq:average-symbols}
    \widetilde{Q'_k\lambda}(p') = \sum_lu_k(p' + l)\tilde\lambda(p' + l)\;,
    \qquad
    \widetilde{(Q_kA)_\mu}(p') = \sum_lu_k(p' + l)v_\mu(p' + l)\tilde A_\mu(p' + l)\;.
\end{equation}
Moreover $Q_k\partial\lambda = \partial^1Q'_k\lambda$, with the difference $\partial^1$ of the lattice $T^{(k)}$.
\item $\abs{u_k(p' + l)} \leq \prod_\mu2(1 + \abs{m_\mu})^{-1}$, $\abs{v_\mu(p)} \leq 1$, $\Delta(p) \geq 4\pi^{-2}\abs p^2$, and $4\pi^{-2}\abs{p'}^2 \leq \Delta_0(p') \leq \abs{p'}^2$.
\item For $p' \neq 0$ let
\begin{equation}\label{eq:phi-symbol}
    \phi_\mu(p') \coloneqq \sum_l\frac{\abs{u_k(p' + l)}^2\abs{v_\mu(p' + l)}^2}{\Delta(p' + l)}\;.
\end{equation}
There are absolute constants $0 < \gamma_0 \leq \gamma_1$ with $\gamma_0 \leq \Delta_0(p')\phi_\mu(p') \leq \gamma_1$ for all $k$, $L$, $\ell_{\mathbb T}$, $\mu$, and $p' \neq 0$.
\item For a $1$-form $B$ on $T^{(k)}$ let $\langle B, \Delta_kB\rangle$ be the minimum of $\langle\partial A, \partial A\rangle$ over the $1$-forms $A$ on $T_\eta$ with $Q_kA = B$. Then
\begin{equation}\label{eq:minimal-form}
\begin{aligned}
    \langle B, \Delta_kB\rangle &= \frac{1}{2}\sum_{\mu, \nu}\frac{1}{(2\pi)^4}\int\sigma_{\mu\nu}(p')\bigl\lvert\widetilde{(\partial^1B)_{\mu\nu}}(p')\bigr\rvert^2\,\dd p'\;,\\
    \sigma_{\mu\nu} &\coloneqq \Bigl(\Delta_0\phi_\mu\,\Delta_0\phi_\nu\sum_\lambda\frac{\abs{\partial^1_\lambda}^2}{\Delta_0^2\phi_\lambda}\Bigr)^{-1}\;,
\end{aligned}
\end{equation}
the term $p' = 0$ being absent, and $\gamma'_0\langle\partial^1B, \partial^1B\rangle \leq \langle B, \Delta_kB\rangle \leq \gamma'_1\langle\partial^1B, \partial^1B\rangle$ with absolute constants $0 < \gamma'_0 \leq \gamma'_1$.
\end{enumerate}
\end{proposition}

These are \cite[(1.18), (1.20), (1.31), (1.36), (1.61), (1.62), (1.65)--(1.67)]{BalabanPRT1} for centred blocks, for which the symbol $u_k$ is real. The form $\Delta_k$ is defined here by the minimum, which does not refer to a gauge; by Theorem~\ref{thm:landau-operators}(iv) at the trivial configuration the minimum is attained at $HB$.

\section{Propagators in a Background Field}\label{sec:propagators}

The fluctuation integral of a renormalisation step is Gaussian to leading order. Its quadratic form is the second variation of the action at the background, to which the constraints on the averages and a gauge-fixing term are added, and everything which the later sections compute is expressed through the inverses of a small number of operators built in this way. This section defines the operators, on a sequence of domains and at a pair of a complex configuration and a current, and states their properties: their analyticity and covariance, the decay and the regularity of their kernels, and their expansions into sums over walks of local operators. The expansions are part of the statements and not only of the proofs, because the main paper compares them term by term on two lattices.

\subsection{The Forms and the Operators}\label{sec:forms}

Throughout this section a sequence of domains with the block size $M$ and the separation $R$ is given, with its layers $\Gamma_j$ and its determining set $\mathbf B$, and $\mathbf U$ is a configuration on the fine lattice $T_\eta$ for which the averages $\mathbf U^{(j)} \coloneqq M^j(\mathbf U)$ and the operators $Q_j(\mathbf U)$, $Q'_j(\mathbf U)$ of Definition~\ref{def:relative-average-order} are defined on the bonds and sites concerned. Forms are forms on $T_\eta$ with values in $\mathfrak g^c$, and $D$, $D^*$, $\nabla$, and the covariant Laplacian carry the index $\mathbf U$ unless another configuration is named. In addition a \emph{Dirichlet domain} $\Omega_\mathrm D \subset \mathbb T$ is given, which is either the torus or a union of cubes of the grid of side $M\eta$ which contains $\Omega_1$; forms are then required to vanish at the sites and bonds outside $\Omega_\mathrm D$, and an operator $\mathcal A$ is replaced by its compression $\chi\mathcal A\chi$ with the indicator $\chi$ of $\Omega_\mathrm D$, which we do not write. The case of a proper subset is needed for the local operators of Section~\ref{sec:walks}.

\begin{definition}[The Second Variation of the Action]\label{def:hessian}
The operator $\Delta(\mathbf U)$ on $1$-forms is the operator which is symmetric with respect to the pairing \eqref{eq:pairing} and satisfies
\begin{equation}\label{eq:action-expansion}
    A(e^{i\eta A}\mathbf U) = A(\mathbf U) + \langle A, J(\mathbf U)\rangle + \frac{1}{2}\langle A, \Delta(\mathbf U)A\rangle + O(A^3)
\end{equation}
for $1$-forms $A$ with values in $\mathfrak g^c$, where $e^{i\eta A}\mathbf U$ is the configuration with the values $e^{i\eta A(b)}\mathbf U(b)$ on the positive bonds.
\end{definition}

The expansion \eqref{eq:action-expansion} is \cite[(3.7), (3.12), pp.~391--392]{BalabanPBG}. Its linear term is the pairing with the current \eqref{eq:current}: the first variation of a plaquette term pairs $D_{\mathbf U}A$ with $\eta^{-2}\im\mathbf U(\partial p)$, and since $A$ has values in $\mathfrak g^c$ only the projection $\pi$ of the latter contributes.

\begin{definition}[The Operators of a Sequence]\label{def:operators}
Let $\mathsf a > 0$, and let $\mathsf a_0 \coloneqq \mathsf a_1 \coloneqq \mathsf a$ and $\mathsf a_{j+1} \coloneqq \mathsf a\mathsf a_j/(\mathsf aL^{-2} + \mathsf a_j)$.
\begin{enumerate}[label=(\roman*)]
\item The \emph{averages of the sequence} are the maps $Q$, from $1$-forms on $T_\eta$ to functions on the bonds of the layers, and $Q'$, from $0$-forms on $T_\eta$ to functions on $\mathbf B$, with
\[
    (QA)(c) \coloneqq \bigl(Q_j(\mathbf U)A\bigr)(c) \text{ for a bond } c \text{ of } \Gamma_j\;,
    \qquad
    (Q'\lambda)(y) \coloneqq \bigl(Q'_j(\mathbf U)\lambda\bigr)(y) \text{ for } y \in \Gamma_j\;,
\]
where $Q_0$ and $Q'_0$ are the identity. The \emph{constraint forms} are
\begin{equation}\label{eq:constraint-forms}
\begin{aligned}
    \langle A, Q^*\mathsf aQA\rangle &\coloneqq \mathsf a\sum_{j=0}^{k}\sum_{c}a_j^2\tr\bigl((QA)(c)\bigr)^2\;,\\
    \langle\lambda, Q'^*\mathsf aQ'\lambda\rangle &\coloneqq \sum_{j=0}^{k}\mathsf a_j\sum_{y \in \Gamma_j}a_j^2\tr\bigl((Q'\lambda)(y)\bigr)^2\;,
\end{aligned}
\end{equation}
the inner sum of the first over the positive bonds of $\Gamma_j$, and $Q^*\mathsf aQ$, $Q'^*\mathsf aQ'$ are the symmetric operators of these forms.
\item $\Delta'_{\mathsf a} \coloneqq \Delta_{\mathbf U} + Q'^*\mathsf aQ'$ on $0$-forms, and $G' \coloneqq (\Delta'_{\mathsf a})^{-1}$.
\item $R \coloneqq I - G'Q'^*(Q'G'^2Q'^*)^{-1}Q'G'$ and $P \coloneqq I - R$, operators on $0$-forms, where $Q'^*$ is the transpose of $Q'$ with respect to the pairing \eqref{eq:pairing} and the pairing $\sum_{y \in \mathbf B}a_{j(y)}^4\tr\mu(y)\mu'(y)$ of functions on $\mathbf B$.
\item $\Delta_{\mathsf a} \coloneqq \Delta(\mathbf U) + DRD^* + Q^*\mathsf aQ$ on $1$-forms, and $G \coloneqq (\Delta_{\mathsf a})^{-1}$.
\end{enumerate}
The inverses are taken on the forms which vanish outside the Dirichlet domain, and each operator is defined whenever the inverses in its definition exist.
\end{definition}

These are the operators \cite[(3.16), (3.18)--(3.27), pp.~393--395]{BalabanPBG}. For a gauge field the pairings are real scalar products, $R$ is the orthogonal projection onto the image of the kernel of $Q'$ under the covariant Laplacian \cite[(3.21), (3.25)]{BalabanPBG}, and $DRD^*$ is the quadratic form of the gauge-fixing term. For a complex configuration the formulae are the definitions.

\begin{definition}[Regular Pairs]\label{def:regular-pairs}
\begin{enumerate}[label=(\roman*)]
\item A \emph{cube of the class} of the sequence, of the index $j$, is a cube $\square$ which is a union of cubes of the grid of side $Ma_j$, has a side $s_\square$ with $10Ma_j \leq s_\square \leq 20Ma_j$, is contained in $\Omega_j \setminus \Omega_{j+2}$, and meets $\Omega_j \setminus \Omega_{j+1}$.
\item A gauge field $U$ on $T_\eta$ is \emph{regular with the constant $\alpha_0$} if for every cube $\square$ of the class, of the index $j$, there is a gauge transformation $u$ on $\square$ with values in $G$ such that $U^u = e^{i\eta A}$ on the bonds of $\square$ with a $1$-form $A$ with values in $\mathfrak g$ and
\begin{equation}\label{eq:regular}
    \abs A \leq \frac{s_\square}{a_j}\,\alpha_0a_j^{-1}\;,
    \qquad
    \abs{\nabla_{\id}A} \leq \frac{s_\square}{a_j}\,\alpha_0a_j^{-2}
\end{equation}
on $\square$. It is \emph{regular to the second order} if in addition $\abs{D^*_{\id}D_{\id}A} \leq (s_\square/a_j)\alpha_0a_j^{-3}$ on $\square$.
\item A configuration $\mathbf U$ on $T_\eta$ is \emph{regular with the constants $(\alpha_0, \alpha_1)$} if $\mathbf U = e^{i\eta A'}U$ with a gauge field $U$ which is regular with the constant $\alpha_0$ and a $1$-form $A'$ with values in $\mathfrak g^c$ such that
\begin{equation}\label{eq:regular-complex}
    \abs{A'(b)} \leq \alpha_1a_j^{-1}\;,
    \qquad
    \abs{(\nabla_UA')(x)} \leq \alpha_1a_j^{-2}
\end{equation}
for the bonds $b$ and the sites $x$ of $\Omega_j$, $0 \leq j \leq k$. It is \emph{regular to the second order} if $U$ is and $\abs{D^*_UD_UA'} \leq \alpha_1a_j^{-3}$ on $\Omega_j$.
\item A pair $(\mathbf U, \mathbf J)$ is \emph{regular with the constants $(\alpha_0, \alpha_1, \gamma)$} if $\mathbf U$ is regular with $(\alpha_0, \alpha_1)$ and $\abs{\mathbf J(b)} \leq \gamma a_j^{-3}$ for the bonds $b$ of $\Omega_j$, $0 \leq j \leq k$.
\end{enumerate}
\end{definition}

These are the conditions \cite[(3.35)--(3.38), p.~396]{BalabanPBG}; the operators of this subsection do not involve the current, which enters with the operators of the Landau gauge.

\subsection{Decay and Regularity}\label{sec:propagator-bounds}

\begin{theorem}[The Operators $G'$, $(Q'G'^2Q'^*)^{-1}$, $P$, and $G$]\label{thm:propagators}
Let $\mathsf a = 1$. There are numbers $M_1$, $R_1$, $\delta_0 > 0$, $a_0 > 0$, and $B_0$, and numbers $B_0(\beta)$ for $0 \leq \beta < 1$ and $B'_0(\epsilon, \beta)$ for $0 < \epsilon \leq 1$ and $0 \leq \beta < 1$, all depending on $L$ and $G$ only, with the following properties. Let a sequence of domains with $M \geq M_1$ and $R \geq R_1$ be given, with a Dirichlet domain, and let $\mathbf U = e^{i\eta A'}U$ be regular with constants $(\alpha_0, \alpha_1)$ such that $M\alpha_0 + \alpha_1 \leq a_0$.
\begin{enumerate}[label=(\roman*)]
\item \emph{Existence.} The operators $\Delta'_{\mathsf a}$, $Q'G'^2Q'^*$, and $\Delta_{\mathsf a}$ are invertible, so that $G'$, $R$, $P$, and $G$ are defined. For a gauge field, $\Delta'_{\mathsf a}$, $G'$, $Q'G'^2Q'^*$, $\Delta_{\mathsf a}$, and $G$ are symmetric and positive definite with respect to the scalar products of the forms with values in $\mathfrak g$.
\item \emph{Analyticity.} For a fixed regular gauge field $U$ the four operators are holomorphic functions of $A'$ on the set \eqref{eq:regular-complex}.
\item \emph{Covariance.} For a gauge transformation $u$ of $T_\eta$ with values in $G$, $G'(\mathbf U^u) = R(u)G'(\mathbf U)R(u)^{-1}$, and the same holds for $R$, $P$, $\Delta_{\mathsf a}$, and $G$; the kernel of $(Q'G'^2Q'^*)^{-1}$ is conjugated by the values of $u$ at the points of $\mathbf B$. All the bounds below are invariant under these transformations.
\item \emph{Pointwise bounds.} Let $y \in \Gamma_j$, $y' \in \Gamma_{j'}$, $x \in \Delta(y)$, and let $\lambda$ be a $0$-form supported in $\Delta(y')$. Then
\begin{equation}\label{eq:propagator-bound}
    \abs{(G'\lambda)(x)} \leq B_0a_j^2e^{-\delta_0d(y, y')}\abs\lambda\;,
    \qquad
    \abs{(\nabla G'\lambda)(x)},\ \abs{(G'\nabla^*\lambda)(x)} \leq B_0a_je^{-\delta_0d(y, y')}\abs\lambda\;,
\end{equation}
and $\abs{(\Delta_{\mathbf U}G'\lambda)(x)} \leq B_0e^{-\delta_0d(y, y')}\abs\lambda$.
\item \emph{H\"older bounds.} Let $0 \leq \beta \leq \beta_0 < 1$ and $0 < \epsilon \leq 1$, let $\zeta$ be a real function supported in $\tilde\Delta(y)$, and let $\lambda$ be supported in $\tilde\Delta(y')$. Then
\begin{align}
    \norm{\zeta\nabla G'\lambda}_\beta,\ \norm{\zeta G'\nabla^*\lambda}_\beta &\leq B_0(\beta_0)a_j^{1-\beta}\bigl(\norm\zeta_\beta^{(j)} + \abs\zeta\bigr)e^{-\delta_0d(y, y')}\abs\lambda\;,\label{eq:propagator-holder}\\
    \abs{(\nabla G'\nabla^*\lambda)(x)} &\leq B'_0(\epsilon, 0)e^{-\delta_0d(y, y')}\bigl(\norm\lambda_\epsilon^{(j')} + \abs\lambda\bigr)\;,\label{eq:propagator-second}\\
    \norm{\zeta\nabla G'\nabla^*\lambda}_\beta &\leq B'_0(\epsilon, \beta_0)a_j^{-\beta}\bigl(\norm\zeta_\beta^{(j)} + \abs\zeta\bigr)e^{-\delta_0d(y, y')}\bigl(\norm\lambda_{\beta+\epsilon}^{(j')} + \abs\lambda\bigr)\;,\label{eq:propagator-second-holder}
\end{align}
the last for $\beta + \epsilon \leq 1$.
\item \emph{Bounds in $L^2$.} For a function $h$ supported in $\Delta(y)$ and $\lambda$ supported in $\Delta(y')$, the $L^2$-norms of $hG'\lambda$, of $h\nabla G'\lambda$ and $hG'\nabla^*\lambda$, and of $h\nabla\nabla G'\lambda$, $h\nabla G'\nabla^*\lambda$, and $hG'\nabla^*\nabla^*\lambda$ are at most $B_0\abs he^{-\delta_0d(y, y')}\norm\lambda$ times $a_j^2$, $a_j$, and $1$, respectively.
\item \emph{The block inverse and the projection.} For $y \in \Gamma_j$ and $y' \in \Gamma_{j'}$ the kernel of $(Q'G'^2Q'^*)^{-1}$ with respect to the pairing of Definition~\ref{def:operators}\emph{(iii)} satisfies
\begin{equation}\label{eq:block-inverse-bound}
    \abs{(Q'G'^2Q'^*)^{-1}(y, y')} \leq B_0a_j^{-4}a_{j'}^{-4}e^{-\delta_0d(y, y')}\;,
\end{equation}
and for $x \in \Delta(y)$, $x' \in \Delta(y')$ the kernels of $P$, $\nabla P$, $P\nabla^*$, and $\nabla P\nabla^*$ with respect to \eqref{eq:pairing} are at most $B_0a_{j'}^{-4}e^{-\frac{1}{2}\delta_0d(y, y')}$ times $1$, $a_j^{-1}$, $a_j^{-1}$, and $a_j^{-2}$, respectively.
\item \emph{The operator $G$.} The bounds \emph{(iv)}--\emph{(vi)} hold for $G$ in the place of $G'$ and $1$-forms $J$ with values in $\mathfrak g^c$ in the place of $\lambda$, the gradients being those of Definition~\ref{def:differences}\emph{(iii)}.
\end{enumerate}
In \emph{(iv)}--\emph{(viii)} the norm of a kernel or of a value is the bond norm of the matrix, or the norm of the linear map of $\mathfrak g^c$ which it defines, and $d$ is the weighted distance of the sequence.
\end{theorem}

This combines \cite[Theorems~3.1--3.4, Corollaries~3.5, 3.6, pp.~397--408]{BalabanPBG} with \cite[(3.48), (3.49)]{BalabanPBG}; at the trivial configuration it contains \cite[Propositions~2.2, 2.3, 2.6]{BalabanPRT2}. The constants do not depend on the sequence, on the Dirichlet domain, on $k$, or on the torus, so that the theorem applies to the local sequences of the cubes in Section~\ref{sec:walks} and to sequences with a block size other than that of the localisation domains. A power of the spacing in a bound may be divided between the two points: by Lemma~\ref{lem:distance}(ii), $a_j^\sigma e^{-\delta d(y, y')}$ is at most a constant times $a_{j'}^\sigma e^{-\delta'd(y, y')}$ for every $\delta' < \delta$ once $RM$ is large, depending on $\sigma$, $\delta - \delta'$, and $L$.

\subsection{Walks}\label{sec:walks}

The bounds of Theorem~\ref{thm:propagators} are obtained from the same bounds for operators on single cubes, where the background can be made small by a gauge transformation, and the passage from the cubes to the lattice is an expansion of each inverse into a sum over walks through overlapping cubes \cite[Sect.~C]{BalabanPBG}. We describe the cubes, the operators attached to them, and the expansions. In this subsection the Dirichlet domain of the given sequence is the torus.

\begin{definition}[The Cubes of the Walks and Their Local Sequences]\label{def:walk-cubes}
Let $R_0$ and $M_0$ be powers of $L$, and let $M = KR_0M_0$ with a positive integer $K$ which is a power of $L$.
\begin{enumerate}[label=(\roman*)]
\item For $0 \leq j \leq k$ the family $\mathscr D_j$ consists of the closed cubes of side $2Ma_j$ whose centre is a corner of a cube of the grid of side $Ma_j$, which are contained in $\Omega_j$, and which contain a cube of that grid which is not contained in $\Omega_{j+1}$. We put $\mathscr D \coloneqq \bigcup_j\mathscr D_j$, and for $\square \in \mathscr D_j$ and $n \geq 1$ we write $\tilde\square^n$ for the cube of side $(2 + 2n)Ma_j$ with the centre of $\square$.
\item The \emph{local sequence} of $\square \in \mathscr D_j$ is the sequence of domains $\Omega_0(\square) \supset \dots \supset \Omega_{j+1}(\square)$, and $\Omega_n(\square) \coloneqq \emptyset$ for $n > j + 1$, with
\[
    \Omega_{j+1}(\square) \coloneqq \tilde\square^3 \cap \Omega_{j+1}\;,
    \qquad
    \Omega_j(\square) \coloneqq \tilde\square^4\;,
\]
and, for $0 \leq n < j$, with $\Omega_n(\square)$ the cube with the centre of $\square$ for which
\begin{equation}\label{eq:local-sequence}
    \dist\bigl(\Omega_n(\square)^c, \Omega_{n+1}(\square)\bigr) = 2R_0M_0a_n\;.
\end{equation}
It is a sequence of domains with the block size $M_0$ and the separation $R_0$, and its Dirichlet domain is $\Omega_0(\square)$, which is contained in $\tilde\square^5$.
\item The \emph{local operators} $G'_\square$, $C_\square \coloneqq (Q'G'^2_\square Q'^*)^{-1}$, $P_\square \coloneqq G'_\square Q'^*C_\square Q'G'_\square$, and $G_\square$ are the operators of Definition~\ref{def:operators} for the local sequence of $\square$, at the same configuration $\mathbf U$.
\end{enumerate}
\end{definition}

By \eqref{eq:local-sequence}, $\dist(\Omega_n(\square)^c, \tilde\square^4) = 2R_0M_0a_n(L^{j-n} - 1)/(L - 1) < 2R_0M_0a_j \leq 2Ma_j$, which gives $\Omega_0(\square) \subset \tilde\square^5$, and the boundaries of the cubes $\Omega_n(\square)$ lie on the boundaries of cubes of the grid of side $M_0a_n$, because $M/M_0$ and $L$ are odd. The cubes of $\mathscr D$ cover the torus: a point of $\Omega_j \setminus \Omega_{j+1}$ lies in a cube of side $2Ma_j$ centred at a corner of the grid at the distance at least $Ma_j$ from $\Omega_j^c$, which is contained in $\Omega_j$. A cube of $\mathscr D_j$ may reach into $\Omega_{j+1}$, by one cube of the grid, and it does not meet $\Omega_{j+2}$ once $R \geq 2$. On such a cube the layers of the local sequence are therefore those of the given one, $\square \cap \Omega_{j+1}(\square) = \square \cap \Omega_{j+1}$ and $\square \subset \Omega_j(\square)$, so that the averages of the two sequences coincide on the forms supported in $\square$. The local sequences are those of \cite[p.~408]{BalabanPBG}, and the local operators are controlled by Theorem~\ref{thm:propagators} for these sequences.

\begin{lemma}[A Partition of Unity]\label{lem:partition-of-unity}
For every sequence of domains there are functions $h_\square \colon \mathbb T \to [0, 1]$, $\square \in \mathscr D$, with the following properties.
\begin{enumerate}[label=(\roman*)]
\item $\sum_{\square \in \mathscr D}h_\square^2 = 1$, and $h_\square$ vanishes outside the cube with the centre of $\square$ and the side $\frac{4}{3}Ma_j$, for $\square \in \mathscr D_j$.
\item $h_\square$ is constant on each cube of the grid of side $a_0$, and $\abs{h_\square(x) - h_\square(x')} \leq C_h\abs{x - x'}_\infty(Ma_j)^{-1}$ with an absolute constant $C_h$, with the corresponding bounds $C_h(Ma_j)^{-2}$ for its second differences.
\item $h_\square$ is determined by the domains of the sequence inside $\tilde\square^2$, and the family is covariant under the translations, reflections, and permutations of the coordinates which map every grid to itself: if $\theta$ maps the sequence $\{ \Omega_j \}$ to $\{ \theta\Omega_j \}$, then the functions of the latter are $h_{\theta\square} = h_\square \circ \theta^{-1}$.
\end{enumerate}
\end{lemma}

\begin{definition}[Commutators with the Partition of Unity]\label{def:commutators}
For a function $h$ on $T_\eta$, regarded as a multiplication operator on forms, and on functions on $\mathbf B$ through its values at the points of $\mathbf B$,
\begin{gather*}
    K'(h) \coloneqq [h, \Delta'_{\mathsf a}]\;,
    \qquad
    K(h) \coloneqq [h, \Delta(\mathbf U) + DD^* + Q^*\mathsf aQ]\;,\\
    P_1(h) \coloneqq [DPD^*, h]\;,
    \qquad
    P_{\square,1}(h) \coloneqq [DP_\square D^*, h]\;.
\end{gather*}
\end{definition}

So $\Delta'_{\mathsf a}h = h\Delta'_{\mathsf a} - K'(h)$ and $\Delta_{\mathsf a}h = h\Delta_{\mathsf a} - K(h) - P_1(h)$, which are \cite[(3.88), (3.104)]{BalabanPBG}. The operators $K'(h_\square)$ and $K(h_\square)$ are sums of difference operators of the first order with coefficients which are differences of $h_\square$ and of operators which couple the points of a block.

\begin{theorem}[The Walk Expansions]\label{thm:walks}
There are numbers $R_0$, $M_0$, $K_1$, and $C_{\mathrm w}$, which depend on $L$ and $G$ only, such that for $M = KR_0M_0$ with $K \geq K_1$, for every sequence of domains with the block size $M$, the separation $R \geq R_1$, and the torus as its Dirichlet domain, and for every configuration $\mathbf U$ as in Theorem~\ref{thm:propagators}, the following holds, with the constants $\delta_0$ and $B_0$ of that theorem.
\begin{enumerate}[label=(\roman*)]
\item \emph{The first approximations.} Let $G'_0 \coloneqq \sum_\square h_\square G'_\square h_\square$, $C_0 \coloneqq \sum_\square h_\square C_\square h_\square$, and $G_0 \coloneqq \sum_\square h_\square G_\square h_\square$, the sums over $\square \in \mathscr D$. Then
\begin{align}
    \Delta'_{\mathsf a}G'_0 &= I - R'\;, & R' &\coloneqq \sum_\square K'(h_\square)G'_\square h_\square\;,\label{eq:first-remainder-scalar}\\
    Q'G'^2Q'^*C_0 &= I - R_C\;, & R_C &\coloneqq -\sum_\square\Bigl(Q'(G'^2 - G'^2_\square)Q'^*h_\square + [Q'G'^2_\square Q'^*, h_\square]\Bigr)C_\square h_\square\;,\label{eq:first-remainder-block}\\
    \Delta_{\mathsf a}G_0 &= I - R_G\;, & R_G &\coloneqq \sum_\square\Bigl(K(h_\square) + (DPD^* - DP_\square D^*)h_\square + P_{\square,1}(h_\square)\Bigr)G_\square h_\square\;.\label{eq:first-remainder}
\end{align}
\item \emph{The series.} $G' = G'_0\sum_{n \geq 0}R'^n$, $(Q'G'^2Q'^*)^{-1} = C_0\sum_{n \geq 0}R_C^n$, and $G = G_0\sum_{n \geq 0}R_G^n$, each series converging in the norms of the bounds of Theorem~\ref{thm:propagators} for its operator.
\item \emph{The walks of $G'$.} For a sequence $\omega = (\square_0, \dots, \square_n)$ of cubes of $\mathscr D$ with $\square_{i-1} \cap \square_i \neq \emptyset$, of the length $\abs\omega \coloneqq n$, let
\begin{equation}\label{eq:walk-term}
    G'_\omega \coloneqq h_{\square_0}G'_{\square_0}h_{\square_0}\prod_{i=1}^{n}K'(h_{\square_i})G'_{\square_i}h_{\square_i}\;.
\end{equation}
Then $G' = \sum_\omega G'_\omega$. Each $G'_\omega$ is a holomorphic function of the configuration which depends only on its values on the bonds of $\tilde\square_0^5 \cup \dots \cup \tilde\square_n^5$, and it is covariant under gauge transformations as $G'$ is.
\item \emph{The bound of a factor and of a walk.} For $\square \in \mathscr D_j$, $y, y' \in \square \cap \mathbf B$, $x \in \Delta(y)$, and $\lambda$ supported in $\Delta(y')$,
\begin{equation}\label{eq:walk-factor-bound}
    \abs{(K'(h_\square)G'_\square h_\square\lambda)(x)} \leq C_{\mathrm w}M^{-1}e^{-\delta_0\abs{y - y'}_1/a_j}\abs\lambda\;,
\end{equation}
and for a walk $\omega$, $y \in \square_0 \cap \Gamma_j$, $y' \in \square_n \cap \mathbf B$, $x \in \Delta(y)$, and $\lambda$ supported in $\Delta(y')$,
\begin{equation}\label{eq:walk-bound}
    \abs{(G'_\omega\lambda)(x)} \leq B_0a_j^2\bigl(C_{\mathrm w}M^{-1}\bigr)^{\abs\omega}e^{-\frac{1}{2}\delta_0d(\omega, y, y')}\abs\lambda\;,
\end{equation}
where $d(\omega, y, y')$ is the infimum of $d(y, y_1) + d(y_1, y_2) + \dots + d(y_n, y')$ over the points $y_i \in \square_i \cap \mathbf B$. The corresponding bounds hold in the other norms of Theorem~\ref{thm:propagators}\emph{(iv)}--\emph{(vi)}.
\item \emph{The remainders.} The kernels of $R'$ and of $R_G$ are bounded by $C_{\mathrm w}M^{-1}e^{-\frac{1}{2}\delta_0d(y, y')}$ in the sense of \eqref{eq:walk-factor-bound}, and that of $R_C$ by $C_{\mathrm w}M^{-1}a_{j'}^{-4}e^{-\frac{1}{2}\delta_0d(y, y')}$ for $y \in \Gamma_j$, $y' \in \Gamma_{j'}$.
\end{enumerate}
\end{theorem}

The parts (i)--(iv) are \cite[(3.87)--(3.96), (3.104)--(3.106), Theorem~3.7, Corollary~3.8]{BalabanPBG}. The remainders $R_C$ and $R_G$ are written here in a form in which the cancellation between the given operators and the local ones is visible; they equal those of \cite[(3.95), (3.105)]{BalabanPBG}, which insert the indicator of $\tilde\square^1$ and the function $\zeta_\square$ in order to separate the terms with a small factor from an overall decay from those with a small commutator. The identities \eqref{eq:first-remainder-scalar}--\eqref{eq:first-remainder} use that the given operators and the local ones coincide on forms supported in $\square$, apart from the projections: for example $h_\square\Delta_{\mathsf a}G_\square h_\square = h_\square^2 - h_\square(DPD^* - DP_\square D^*)G_\square h_\square$, since $\Delta_{\mathsf a}$ and the local operator differ by $D(P_\square - P)D^*$ there.

The remainders $R_C$ and $R_G$ still contain the operators $G'$ and $P$ of the whole lattice, and the expansions of the block inverse and of $G$ into walks whose terms are products of local operators only \cite[Theorems~3.9, 3.10]{BalabanPBG} are obtained by inserting the walk expansion of $G'$ into them. For these expansions the papers give a bound of a whole term and not of its factors \cite[p.~413]{BalabanPBG}. The main paper, however, compares the terms of two lattices by replacing one factor at a time, and for this it needs a term to be a product whose bound follows from bounds of the factors. We therefore state the expansions in this form. The bounds of the factors are bounds between three kinds of objects: sources, fields, and data on the determining set.

\begin{definition}[Block Norms and Block Bounds]\label{def:block-bounds}
Fix $0 < \beta < 1$. For $y \in \Gamma_j$ we define the following \emph{block norms}:
\begin{enumerate}[label=(\roman*)]
\item for a form $f$ on $T_\eta$, regarded as a \emph{source}, $\abs f_y$ is the supremum of the components of $f$ over $\Delta(y)$;
\item for a form $u$ on $T_\eta$, regarded as a \emph{field}, $N_y(u)$ is the supremum over $\Delta(y)$ of $a_j^{-1}\abs u + \abs{\nabla u}$, plus $a_j^\beta$ times the supremum of $(\abs{x - x'}_1/a_j)^{-\beta}\abs{R(\mathbf U(\gamma))(\nabla u)(x') - (\nabla u)(x)}$ over the pairs of sites $x \neq x'$ of $\tilde\Delta(y)$ and the shortest paths $\gamma$ between them;
\item for a function $\mu$ on $\mathbf B$, or on the positive bonds of the layers, regarded as \emph{data}, the block norm is $\abs{\mu(y)}$, or the largest of the norms $\abs{\mu(c)}$ over the bonds $c$ with $c_- = y$.
\end{enumerate}
Let $\mathcal T$ be a linear operator which maps objects of one of the three kinds to objects of one of the three kinds, let $X \subset \mathbb T$, and let $b \geq 0$, $p \in \mathbb Z$, and $\delta > 0$. We say that $\mathcal T$ has the \emph{block bound $(b, p, \delta)$ in $X$} if $\mathcal Tf$ depends only on the restriction of $f$ to $X$, if $[\mathcal Tf]_y = 0$ for every $y \in \mathbf B \setminus X$, and if for every $f$ and every $y \in \mathbf B \cap X$
\begin{equation}\label{eq:block-bound}
    [\mathcal Tf]_y \leq b\,a_{j(y)}^p\sum_{y' \in \mathbf B \cap X}e^{-\delta d(y, y')}[f]_{y'}\;,
\end{equation}
where $[\cdot]_y$ denotes the block norm of the kind of the object.
\end{definition}

The exponent $p$ counts the powers of the spacing which the operator carries: a propagator on sources has $p = 1$ as a map into fields, since the norm $N_y$ measures a field through its first differences, and a difference operator of the first order from fields to sources has $p = 0$.

\begin{lemma}[Composition of Block Bounds]\label{lem:composition}
Let $\delta > 0$, $0 < \alpha < 1$, and $S \geq 0$, and let the sequence of domains satisfy $RM \geq r_1(\frac{1}{2}\alpha, \delta)$ and $\alpha\delta RM \geq 2S\log L$, with the number $r_1$ of Lemma~\ref{lem:summation}. Let $\mathcal T_1, \dots, \mathcal T_n$ be operators such that $\mathcal T_i$ has the block bound $(b_i, p_i, \delta)$ in $X_i$ and the product $\mathcal T_1\cdots\mathcal T_n$ is defined, and suppose that $\abs{p_i + p_{i+1} + \dots + p_n} \leq S$ for $2 \leq i \leq n$. Then for every $f$ and every $y \in \mathbf B \cap X_1$,
\begin{equation}\label{eq:composition}
    [\mathcal T_1\cdots\mathcal T_nf]_y \leq \bigl(c_1L^S\bigr)^{n-1}\Bigl(\prod_{i=1}^{n}b_i\Bigr)a_{j(y)}^{p_1 + \dots + p_n}\sum_{y' \in \mathbf B \cap X_n}e^{-(1 - \alpha)\delta d_{\mathbf X}(y, y')}[f]_{y'}\;,
\end{equation}
where $c_1 \coloneqq c_1(\frac{1}{2}\alpha, \delta)$ and $d_{\mathbf X}(y, y')$ is the infimum of $d(y, y_1) + d(y_1, y_2) + \dots + d(y_{n-1}, y')$ over the points $y_i \in \mathbf B \cap X_i \cap X_{i+1}$.
\end{lemma}

\begin{proof}
Applying \eqref{eq:block-bound} to $\mathcal T_1$, then to $\mathcal T_2$, and so on, and writing $y_0 \coloneqq y$, we obtain
\[
    [\mathcal T_1\cdots\mathcal T_nf]_{y_0} \leq \Bigl(\prod_ib_i\Bigr)\sum_{y_1, \dots, y_n}\Bigl(\prod_{i=1}^{n}a_{j(y_{i-1})}^{p_i}e^{-\delta d(y_{i-1}, y_i)}\Bigr)[f]_{y_n}\;,
\]
where $y_i$ runs over $\mathbf B \cap X_i \cap X_{i+1}$ for $i < n$, because the bound of $\mathcal T_i$ contains the block norms of its argument at the points of $X_i$ only and the block norms of $\mathcal T_{i+1}(\cdot)$ vanish outside $X_{i+1}$, and $y_n$ runs over $\mathbf B \cap X_n$. We first collect the powers of the spacings at the point $y_0$. With $s_i \coloneqq p_i + \dots + p_n$,
\[
    \prod_{i=1}^{n}a_{j(y_{i-1})}^{p_i} = a_{j(y_0)}^{s_1}\prod_{i=2}^{n}\Bigl(\frac{a_{j(y_{i-1})}}{a_{j(y_{i-2})}}\Bigr)^{s_i}\;,
\]
as one sees by collecting the exponent of each $a_{j(y_l)}$ on the right side, which is $s_{l+1} - s_{l+2} = p_{l+1}$. The quotient of the spacings of two points is $L^{j(y_{i-1}) - j(y_{i-2})}$, and by Lemma~\ref{lem:distance}(ii) the difference of the scales of two points is at most $1 + d/(RM)$, with $d$ their weighted distance. Since $\abs{s_i} \leq S$ and $S\log L/(RM) \leq \frac{1}{2}\alpha\delta$, each quotient is at most $L^Se^{\frac{1}{2}\alpha\delta d(y_{i-2}, y_{i-1})}$. The summand is therefore at most $L^{S(n-1)}a_{j(y_0)}^{s_1}\prod_ie^{-(1 - \frac{1}{2}\alpha)\delta d(y_{i-1}, y_i)}[f]_{y_n}$. By the definition of $d_{\mathbf X}$, $\sum_id(y_{i-1}, y_i) \geq d_{\mathbf X}(y_0, y_n)$, so that the product of the exponentials is at most $e^{-(1 - \alpha)\delta d_{\mathbf X}(y_0, y_n)}\prod_ie^{-\frac{1}{2}\alpha\delta d(y_{i-1}, y_i)}$. In the last product we bound the factor with $i = n$ by one and sum over $y_{n-1}, \dots, y_1$ in this order, each sum being at most $c_1(\frac{1}{2}\alpha, \delta)$ by \eqref{eq:summation}.
\end{proof}

\begin{theorem}[Walks as Products of Local Factors]\label{thm:walk-structure}
Under the hypotheses of Theorem~\ref{thm:walks}, and with a constant $C_{\mathrm l}$ which depends on $L$, $G$, and $\beta$ only:
\begin{enumerate}[label=(\roman*)]
\item \emph{The factors of the walks of $G'$.} For $\square \in \mathscr D_j$, the operator $h_\square G'_\square h_\square$ maps sources to fields with the block bound $(C_{\mathrm l}, 1, \delta_0)$ in $\square$, and $K'(h_\square)G'_\square h_\square$ maps sources to sources with the block bound $(C_{\mathrm l}M^{-1}, 0, \delta_0)$ in $\square$. If $G'_\square$ is replaced in them by an operator $\mathcal O$ with $N_y(\mathcal Of) \leq \theta a_je^{-\delta d(y, y')}\abs f_{y'}$ for all $y, y' \in \mathbf B \cap \square$ and all sources $f$ supported in $\Delta(y')$, the two operators have the block bounds $(C_{\mathrm l}\theta, 1, \delta)$ and $(C_{\mathrm l}M^{-1}\theta, 0, \delta)$ in $\square$.
\item \emph{The block inverse and the operator on $1$-forms.} There are sets $\mathscr W_C$ and $\mathscr W_G$ and, for each of their elements $\omega$, a number $\abs\omega \geq 0$, a sequence of cubes $\square_0, \dots, \square_m$ of $\mathscr D$ in which consecutive cubes intersect, and an operator $W_\omega$, such that
\begin{equation}\label{eq:walk-expansions}
    (Q'G'^2Q'^*)^{-1} = \sum_{\omega \in \mathscr W_C}W_\omega\;,
    \qquad
    G = \sum_{\omega \in \mathscr W_G}W_\omega\;,
\end{equation}
and such that each $W_\omega$ is a product $\mathcal T_1\cdots\mathcal T_n$ of $n \leq C_{\mathrm l}(\abs\omega + 1)$ operators with the following properties. Each $\mathcal T_i$ has the form $E_i\mathcal O_iE'_i$, where $\mathcal O_i$ is a local operator $G'_X$, $C_X$, or $G_X$ of the local sequence of a cube $X$ which is a cube of the walk or its enlargement $\tilde\square_l^2$, or is absent, and where $E_i$ and $E'_i$ are products of at most $C_{\mathrm l}$ operators of the following kinds: the multiplication by a function which is determined by the sequence of domains; the commutator of such a multiplication with $\Delta'_{\mathsf a}$ or with $\Delta(\mathbf U) + DD^* + Q^*\mathsf aQ$; one of the operators $D$, $D^*$, $Q$, $Q'$, $Q^*$, $Q'^*$; and the difference $\Delta'_{\mathsf a} - \Delta'_{\mathsf a,X}$ of the scalar operators of the given sequence and of the local sequence of $X$. Each $\mathcal T_i$ has a block bound $(b_i, p_i, \frac{1}{2}\delta_0)$ in a cube of the walk, between objects of the three kinds, the sums $p_i + \dots + p_n$ are at most $4$ in absolute value, and
\begin{equation}\label{eq:walk-factor-product}
    \prod_{i=1}^{n}b_i \leq C_{\mathrm l}\bigl(C_{\mathrm w}M^{-1}\bigr)^{\abs\omega}\;,
    \qquad
    \sum_{i=1}^{n}p_i = -4 \text{ for } \mathscr W_C \text{ and } 1 \text{ for } \mathscr W_G\;.
\end{equation}
Moreover, if in one $\mathcal T_i$ the local operator $\mathcal O_i$ is replaced by an operator which satisfies, between the blocks of $X$, $\theta$ times the bounds which Theorem~\ref{thm:propagators} gives for $\mathcal O_i$, with a rate $\delta \geq \frac{1}{2}\delta_0$, then $\mathcal T_i$ has the block bound $(C_{\mathrm l}\theta b_i, p_i, \frac{1}{2}\delta_0)$.
\item \emph{Naturality.} The sets $\mathscr W_C$ and $\mathscr W_G$, the numbers $\abs\omega$, the cubes of a walk, and for each factor its kind, its cube, and its function are determined by the sequence of domains as a sequence of subsets of the torus and by $M$, $R_0$, and $M_0$. They do not depend on the configuration or on the spacing $\eta$ of the fine lattice, and they are mapped to those of the image sequence by the translations, reflections, and permutations of the coordinates which map every grid to itself.
\item \emph{Analyticity and locality.} Each $W_\omega$ is a holomorphic function of the configuration on the set of Theorem~\ref{thm:propagators}, it depends only on its values on the bonds of $\tilde\square_0^5 \cup \dots \cup \tilde\square_m^5$, and it is covariant under gauge transformations as the operator which it expands. The block bounds of the factors hold on that whole set.
\end{enumerate}
\end{theorem}

The term \eqref{eq:walk-term} of a walk of $G'$ is the product of $h_{\square_0}G'_{\square_0}h_{\square_0}$ and of the operators $K'(h_{\square_i})G'_{\square_i}h_{\square_i}$, so that (i) and Lemma~\ref{lem:composition} with $S = 1$ give \eqref{eq:walk-bound} in the supremum norm and in the norm $N_y$, which contains the H\"older norm of the first differences. In the same way (ii) and the lemma with $S = 4$ give the bounds of the terms of the other two expansions \cite[(3.99), (3.108)]{BalabanPBG}, in the distance through the cubes of the walk, the constant $c_1L^4$ of each factor being absorbed into $C_{\mathrm w}$; and they give the same bound, multiplied by $C_{\mathrm l}\theta$, for the product in which one local operator is replaced by an operator $\theta$ times as large. This last statement is the one on which the comparison of two lattices in the main paper rests.

\subsection{The Operators of the Landau Gauge}\label{sec:landau-operators}

The background field of a step is the minimiser of the action with prescribed averages, and to second order the problem is linear: one minimises $\frac{1}{2}\langle A, \Delta(\mathbf U)A\rangle$ among the $1$-forms with $QA = B$. The functional and the constraint are invariant under the gauge transformations $A \mapsto A - D\lambda$ with $Q'\lambda = 0$, by Proposition~\ref{prop:relative-average-order}(v), but only up to terms which contain the current, and the gauge is fixed by the condition $RD^*A = 0$. The solution operator of this problem and of a variant of it, in which the constraint is corrected to the second order, are the operators through which the background field and the covariance of the fluctuation are written in the later sections \cite[Sect.~D]{BalabanPBG}. In the formulae the current of the configuration appears explicitly, and following \cite[p.~261]{BalabanRG1} we define the operators at a pair, with the variable $\mathbf J$ in its place.

\begin{definition}[The Operators at a Pair]\label{def:landau-operators}
Let $(\mathbf U, \mathbf J)$ be a pair on $T_\eta$ for which the operators of Definition~\ref{def:operators} exist. For a $1$-form $A$ let $\lambda_A \coloneqq G'RD^*A$ and $\mathcal PA \coloneqq A - D\lambda_A$.
\begin{enumerate}[label=(\roman*)]
\item $\Delta'_\pi = \Delta'_\pi(\mathbf U, \mathbf J)$ is the symmetric operator on $1$-forms with $\langle A, \Delta'_\pi A\rangle = \langle X(A), \mathbf J\rangle$, where $X(A)$ is the $1$-form with
\begin{equation}\label{eq:gauge-defect}
    X(A)(b) \coloneqq i[\lambda_A(b_-), A(b)] - i[A(b), R(\mathbf U(b))\lambda_A(b_+)] - i[\lambda_A(b_-), (D\lambda_A)(b)]
\end{equation}
on the positive bonds $b$.
\item $G_\pi \coloneqq (\Delta_{\mathsf a} - \Delta'_\pi)^{-1}$ and $H \coloneqq G_\pi Q^*(QG_\pi Q^*)^{-1}$, where $Q^*$ is the transpose of $Q$ with respect to \eqref{eq:pairing} and the pairing $\sum_ca_{j(c)}^4\tr\mu(c)\mu'(c)$ of functions on the positive bonds $c$ of the layers.
\item For a $1$-form $A$ let $\mathsf C^{(2)}(A)$ be the function on the bonds of the layers which on a bond $c$ of $\Gamma_j$ equals $a_j^{-1}$ times the part of the second order in $A$ of the remainder $C_j(\mathbf U; a_jA)(c)$ of Definition~\ref{def:relative-average-order}. Then $\Delta^{(2)} = \Delta^{(2)}(\mathbf U, \mathbf J)$ is the symmetric operator with $\langle A, \Delta^{(2)}A\rangle = 2\langle H\mathsf C^{(2)}(A), \mathbf J\rangle$, and $\Delta^{(2)}_\pi \coloneqq \mathcal P^T\Delta^{(2)}\mathcal P$, with the transpose $\mathcal P^T = I - DRG'D^*$.
\item $G_1 \coloneqq (\Delta_{\mathsf a} - \Delta'_\pi - \Delta^{(2)}_\pi)^{-1}$ and $H_1 \coloneqq G_1Q^*(QG_1Q^*)^{-1}$.
\end{enumerate}
\end{definition}

These are \cite[(3.119)--(3.129), (3.134), (3.135)]{BalabanPBG} with $\mathbf J$ in the place of the current. The map $\mathcal P$ takes a $1$-form to the gauge $RD^*A = 0$ by a transformation with $Q'\lambda = 0$: indeed $Q'G'R = 0$ and $R^2 = R$ by the definition of $R$, so that $Q'\lambda_A = 0$, $\Delta_{\mathbf U}\lambda_A = RD^*A$, and $RD^*\mathcal PA = 0$.

\begin{theorem}[The Operators of the Landau Gauge]\label{thm:landau-operators}
Let $\mathsf a = 1$. There are numbers $a_2 > 0$, $\delta_1 > 0$, $B_1$, and $B_1(\beta)$, $B'_1(\epsilon, \beta)$ as in Theorem~\ref{thm:propagators}, which depend on $L$ and $G$ only, such that for every sequence of domains with $M \geq M_1$ and $R \geq R_1$ and every regular pair $(\mathbf U, \mathbf J)$ with constants $M\alpha_0 + \alpha_1 + \gamma \leq a_2$ the following holds.
\begin{enumerate}[label=(\roman*)]
\item \emph{Existence and analyticity.} The operators of Definition~\ref{def:landau-operators} exist, they are holomorphic in $(A', \mathbf J)$ for a fixed regular gauge field $U$, and
\begin{equation}\label{eq:landau-series}
    G_\pi = \sum_{n \geq 0}G(\Delta'_\pi G)^n\;,
    \qquad
    G_1 = \sum_{n \geq 0}G\bigl((\Delta'_\pi + \Delta^{(2)}_\pi)G\bigr)^n\;,
\end{equation}
with the operator $G$ of Definition~\ref{def:operators}, the series converging in the norms of \emph{(iii)}.
\item \emph{The vertices.} For $1$-forms $A_1$ supported in $\Delta(y)$, $y \in \Gamma_j$, and $A_2$ supported in $\Delta(y')$, $y' \in \Gamma_{j'}$,
\begin{equation}\label{eq:vertex-bound}
    \abs{\langle A_1, \Delta'_\pi A_2\rangle} \leq B_1\gamma\bigl(\norm{D^*A_1}_{L^1} + a_j^{-1}\norm{A_1}_{L^1}\bigr)e^{-\frac{1}{2}\delta_0d(y, y')}\bigl(\abs{D^*A_2} + a_{j'}^{-1}\abs{A_2}\bigr)\;,
\end{equation}
where $\norm f_{L^1} \coloneqq \sum_x\eta^4\abs{f(x)}$, and for a bond $b$ of $\Delta(y)$, $\abs{(\Delta^{(2)}A)(b)} \leq B_1\gamma a_j^{-2}\abs A$ with the supremum of $A$ over the blocks $\Delta(y'')$ with $d(y, y'') \leq 2$.
\item \emph{Bounds.} The operators $G_\pi$ and $G_1$ satisfy the bounds \emph{(iv)}--\emph{(vi)} of Theorem~\ref{thm:propagators} for $1$-forms, with the constants $B_1$, $B_1(\beta)$, $B'_1(\epsilon, \beta)$, and $\delta_1$, except the bound of the covariant Laplacian. For bonds $c$ of $\Gamma_j$ and $c'$ of $\Gamma_{j'}$,
\begin{equation}\label{eq:block-form-inverse}
    \abs{(QG_\pi Q^*)^{-1}(c, c')} \leq B_1a_j^{-2}a_{j'}^{-4}e^{-\delta_1d(c_-, c'_-)}\;,
\end{equation}
and the same holds with $G_1$. For $x \in \Delta(y)$, $y \in \Gamma_j$, a bond $c'$ of $\Gamma_{j'}$, a real function $\zeta$ supported in $\tilde\Delta(y)$, and $0 \leq \beta \leq \beta_0 < 1$, the kernels of $H$ and of $H_1$ satisfy
\begin{equation}\label{eq:minimiser-kernel}
\begin{aligned}
    \abs{H(x, c')} + a_j\abs{\nabla H(x, c')} &\leq B_1a_{j'}^{-4}e^{-\frac{1}{2}\delta_1d(y, c'_-)}\;,\\
    \norm{\zeta\nabla H(\cdot, c')}_\beta &\leq B_1(\beta_0)a_j^{-1-\beta}\bigl(\norm\zeta_\beta^{(j)} + \abs\zeta\bigr)a_{j'}^{-4}e^{-\frac{1}{2}\delta_1d(y, c'_-)}\;.
\end{aligned}
\end{equation}
\item \emph{The current of the configuration.} If $\mathbf J = J(\mathbf U)$, then $\Delta(\mathbf U) - \Delta'_\pi = \mathcal P^T\Delta(\mathbf U)\mathcal P$, the form $\langle A, (\Delta(\mathbf U) - \Delta'_\pi)A\rangle$ is invariant under $A \mapsto A - D\lambda$ for $Q'\lambda = 0$, and
\begin{equation}\label{eq:landau-identity}
    RD^*G_\pi Q^* = 0\;,
    \qquad
    QG_\pi DR = 0\;,
\end{equation}
so that $RD^*H = 0$, while $QH = I$ at every pair by the definition of $H$; the same identities hold for $G_1$ and $H_1$. If moreover $\mathbf U$ is a gauge field, then $G_\pi$, $G_1$, $QG_\pi Q^*$, and $QG_1Q^*$ are symmetric and positive definite, and $HB$ is the unique minimiser of $\frac{1}{2}\langle A, \Delta(\mathbf U)A\rangle$ among the $1$-forms with values in $\mathfrak g$ which satisfy $QA = B$ and $RD^*A = 0$.
\end{enumerate}
\end{theorem}

This is \cite[Theorem~3.12]{BalabanPBG} with \cite[(3.124), (3.126), (3.130)--(3.138)]{BalabanPBG}. The identities \eqref{eq:landau-identity} rest on the gauge invariance of the action and do not hold at a pair whose second entry is not the current of the first. The bound \eqref{eq:block-form-inverse} is \cite[(3.132)]{BalabanPBG}, which is printed there without a proof.

\subsection{The Propagator of the Fluctuation}\label{sec:fluctuation-propagator}

After the background field has been split off, the fluctuation is a $1$-form with vanishing averages in the gauge $RD^*A = 0$, and its Gaussian measure has the quadratic form $\Delta(\mathbf U) - \Delta'_\pi - \Delta^{(2)}_\pi$ on this subspace. Its covariance is expressed through $G_1$ \cite[(3.139)--(3.153)]{BalabanPBG}.

\begin{definition}[The Projection and the Covariance of the Constrained Fluctuation]\label{def:fluctuation-propagator}
At a pair $(\mathbf U, \mathbf J)$ for which the operators of Definition~\ref{def:landau-operators} exist,
\begin{equation}\label{eq:fluctuation-propagator}
    \mathfrak P \coloneqq I - G_1Q^*(QG_1Q^*)^{-1}Q - G_1DRD^*\;,
    \qquad
    \mathfrak G \coloneqq G_1 - G_1DRD^*G_1 - G_1Q^*(QG_1Q^*)^{-1}QG_1\;.
\end{equation}
\end{definition}

\begin{theorem}[The Propagator of the Fluctuation]\label{thm:fluctuation-propagator}
Under the hypotheses of Theorem~\ref{thm:landau-operators}:
\begin{enumerate}[label=(\roman*)]
\item $\mathfrak P$ and $\mathfrak G$ are holomorphic in $(A', \mathbf J)$, and $\mathfrak G$ satisfies the bounds \emph{(iv)}--\emph{(vi)} of Theorem~\ref{thm:propagators} for $1$-forms, except that of the covariant Laplacian, with the constants of Theorem~\ref{thm:landau-operators}.
\item If $\mathbf J = J(\mathbf U)$, then
\begin{equation}\label{eq:gauge-fixing-identities}
    RD^*G_1 = RG'D^*\;,
    \qquad
    G_1DR = DG'R\;,
    \qquad
    RD^*G_1DR = R\;,
\end{equation}
and consequently $Q\mathfrak P = 0$, $RD^*\mathfrak P = 0$, $\mathfrak P^2 = \mathfrak P$, $\mathfrak P A = A$ for every $A$ with $QA = 0$ and $RD^*A = 0$, and $\mathfrak G = \mathfrak PG_1 = G_1\mathfrak P^T$.
\item If moreover $\mathbf U$ is a gauge field, then $\mathfrak G$ is symmetric and positive semidefinite, its range is the space $\mathcal S$ of the $1$-forms with values in $\mathfrak g$ which satisfy $QA = 0$ and $RD^*A = 0$, and for every $1$-form $J$ with values in $\mathfrak g$,
\begin{equation}\label{eq:fluctuation-covariance}
    \int_{\mathcal S}e^{\langle A, J\rangle}\,\dd\mu(A) = e^{\frac{1}{2}\langle J, \mathfrak GJ\rangle}\;,
\end{equation}
where $\mu$ is the Gaussian probability measure on $\mathcal S$ with the density proportional to $\exp(-\frac{1}{2}\langle A, (\Delta(U) - \Delta'_\pi - \Delta^{(2)}_\pi)A\rangle)$ with respect to the Lebesgue measure of $\mathcal S$.
\end{enumerate}
\end{theorem}

This is \cite[Theorem~3.13]{BalabanPBG} with \cite[(3.147)--(3.153)]{BalabanPBG}. In (iii) the quadratic form is positive definite on $\mathcal S$ because on $\mathcal S$ it equals $\langle A, G_1^{-1}A\rangle$.

\subsection{The Dependence on the Sequence}\label{sec:change-of-sequence}

A term of the density is computed with the operators of its own sequence of domains, and different terms have different sequences. Where two sequences agree, their operators agree up to an error which decays in the distance to the region in which the sequences differ \cite[Theorem~3.14, pp.~426--427]{BalabanPBG}. We state this for two sequences which agree on a region of any shape, which is the form in which the localisation of the later sections uses it.

\begin{theorem}[Change of the Sequence]\label{thm:change-of-sequence}
Let $\{ \Omega_j \}$ and $\{ \Omega'_j \}$ be two sequences of domains with the same block size and separation, as in Theorem~\ref{thm:landau-operators}, and let $W \subset \mathbb T$ be a union of cubes of the grid of side $M\eta$ such that $\Omega_j \cap W = \Omega'_j \cap W$ for all $j$. Let $(\mathbf U, \mathbf J)$ be a pair which is regular with the constants of Theorem~\ref{thm:landau-operators} for both sequences. Let $\mathcal A$ and $\mathcal A'$ be the operators of one of the kinds $G'$, $(Q'G'^2Q'^*)^{-1}$, $P$, $G$, $G_\pi$, $G_1$, $(QG_\pi Q^*)^{-1}$, $(QG_1Q^*)^{-1}$, $H$, $H_1$, $\mathfrak G$ for the two sequences. Then, for points $y$, $y'$ of the determining sets which lie in $W$, the difference $\mathcal A - \mathcal A'$ satisfies every bound which Theorems~\ref{thm:propagators}, \ref{thm:landau-operators}, and~\ref{thm:fluctuation-propagator} state for $\mathcal A$ at the localisations $y$, $y'$, with a constant which depends on $L$ and $G$ only in the place of the constant of the bound, with half the decay rate, and with the additional factor
\begin{equation}\label{eq:change-of-sequence}
    \exp\Bigl(-\frac{1}{2}\delta_1\bigl(d(y, W^c) + d(y', W^c)\bigr)\Bigr)\;,
\end{equation}
where $d(y, W^c)$ is the infimum of the weighted distance of $y$ from the points of the complement of $W$, which is the same for both sequences.
\end{theorem}

For $W = \Omega_k \cap \Omega'_k$ and points of the lattice $T^{(k)}$ this is \cite[Theorem~3.14]{BalabanPBG}, whose factor \cite[(3.154)]{BalabanPBG} is the exponential of the smallest length of a path from $y$ to $y'$ through the complement of $W$. A path from a point of $W$ which leaves $W$ and stays in $W$ until then has the same weighted length for both sequences, so that $d(y, W^c)$ is well defined.

\section{Regular Fields and Gauge Conditions}\label{sec:gauges}

The fields over which a step integrates have prescribed averages on the determining set, and two such fields which differ by a gauge transformation which is trivial on the determining set are equivalent. To compute with them one chooses a representative in each class. Two choices are used. The \emph{axial gauge} is algebraic: the representative is found by an explicit formula, block by block, and this is the gauge in which the variables of the construction are defined. It gives, however, no control of the size of a field in terms of its plaquettes beyond a bound which does not improve with the scale. The \emph{Landau gauge} relative to a background is the gauge of the operators of Section~\ref{sec:propagators}, in which a regular field differs from the background by a fluctuation of the natural size, and the passage from the first gauge to the second is the main result of \cite{BalabanREG}. Throughout this section a sequence of domains with the block size $M$ and the separation $R$ is given.

\subsection{Fields with Given Averages}\label{sec:regular-fields}

\begin{definition}[Regular Fields and Data]\label{def:regular-fields}
\begin{enumerate}[label=(\roman*)]
\item For $\alpha_0 > 0$ the set $\mathfrak U(\alpha_0)$ consists of the gauge fields $U$ on $T_\eta$ with
\begin{equation}\label{eq:plaquette-regular}
    \abs{U(\partial p) - \id} \leq \alpha_0\Bigl(\frac{\eta}{a_j}\Bigr)^2\;,
    \qquad
    \abs{J(U)(b)} \leq \alpha_0a_j^{-3}\;,
\end{equation}
for the plaquettes $p$ and the bonds $b$ of $\Omega_j$ of the scale $0$, for $0 \leq j \leq k$.
\item \emph{Data} on the determining set are a gauge field $V$ on the bonds of the layers, that is, a map which assigns an element of $G$ to every bond of every layer, compatibly with the reversal of bonds. The set $\mathfrak V(V)$ consists of the gauge fields $U$ on $T_\eta$ with $M^j(U)(c) = V(c)$ for the bonds $c$ of $\Gamma_j$, $0 \leq j \leq k$.
\item $\mathcal G_{\mathbf B}$ is the group of the gauge transformations $u \colon T_\eta \to G$ with $u(y) = \id$ for $y \in \mathbf B$.
\end{enumerate}
\end{definition}

These are the spaces \cite[(1.7)--(1.14), pp.~77--78]{BalabanREG}, with the current \eqref{eq:current} in the second condition. By Lemma~\ref{lem:average}(iii) and Lemma~\ref{lem:action}(ii), $\mathfrak U(\alpha_0)$ is invariant under all gauge transformations with values in $G$, and $\mathfrak V(V)$ under those of $\mathcal G_{\mathbf B}$.

Inside a block the points are reached from the centre along paths, and the gauge conditions are conditions on the transports along them. The construction uses two kinds of transports: along one fixed path, which is simpler but distinguishes an order of the coordinates, and the mean over the paths of all orders \cite[p.~253]{BalabanRG1}.

\begin{definition}[Transports in a Block]\label{def:transports}
Let $1 \leq n \leq k$, let $\mathbf U$ be a configuration on $T^{(n-1)}$, $y \in T^{(n)}$, and $x \in B(y)$. The \emph{axial transport} $\mathbf U_{\mathrm{ax}}(y, x)$ is $\mathbf U(\Gamma)$ for the path $\Gamma$ of $\mathbf G(y, x)$ which traverses the coordinates in increasing order, and the \emph{mean transport} $\mathbf U_{\mathrm m}(y, x)$ is the transport $\mathbf U(y, x)$ of Definition~\ref{def:group-mean}(ii). A \emph{transport rule} is one of the symbols $\mathrm{ax}$ and $\mathrm m$, fixed throughout an argument and written $\mathfrak t$. For both rules $\mathbf U^u_{\mathfrak t}(y, x) = u(y)\mathbf U_{\mathfrak t}(y, x)u(x)^{-1}$.
\end{definition}

\subsection{The Axial Gauge}\label{sec:axial-gauge}

For $y \in \Gamma_j$ and a point $x_0$ of the block $B^j(y)$ there is a unique chain $x_0, x_1, \dots, x_j = y$ with $x_n \in T^{(n)}$ and $x_n \in B(x_{n+1})$, the \emph{chain} of $x_0$.

\begin{definition}[Axial Gauges]\label{def:axial-gauge}
Let $U_0$ and $U$ be gauge fields on $T_\eta$ for which the averages and the transports below are defined. The field $U$ is \emph{in the axial gauge of $U_0$}, with respect to the rule $\mathfrak t$, if
\begin{equation}\label{eq:axial-gauge}
    M^n(U)_{\mathfrak t}(x_{n+1}, x_n) = M^n(U_0)_{\mathfrak t}(x_{n+1}, x_n)
\end{equation}
for every $1 \leq j \leq k$, every $y \in \Gamma_j$, every chain in $B^j(y)$, and every $0 \leq n < j$. It is \emph{in the axial gauge} if this holds with the right side replaced by $\id$.
\end{definition}

With the axial transports these are the conditions \cite[(1.15), p.~78]{BalabanREG} and \cite[(1.19), (1.20), p.~79]{BalabanREG}; the latter are written there as a product along the path of the relative variables $M^n(U)M^n(U_0)^{-1}$ transported to the centre, which telescopes to \eqref{eq:axial-gauge}.

\begin{lemma}[The Representative in the Axial Gauge]\label{lem:axial-representative}
Let $U_0$ and $U$ be as in Definition~\ref{def:axial-gauge}. There is exactly one $u \in \mathcal G_{\mathbf B}$ for which $U^u$ is in the axial gauge of $U_0$. On the chain of a point $x_0 \in B^j(y)$, $y \in \Gamma_j$, it is given by $u(x_j) = \id$ and
\begin{equation}\label{eq:axial-transformation}
    u(x_n) = M^n(U_0)_{\mathfrak t}(x_{n+1}, x_n)^{-1}\,u(x_{n+1})\,M^n(U)_{\mathfrak t}(x_{n+1}, x_n)\;,
    \qquad
    n = j - 1, \dots, 0\;.
\end{equation}
In particular $u$ on $B^j(y)$ depends only on $U$ and $U_0$ on the bonds with both ends in $B^j(y)$, it depends on them analytically, and if $w$ is a gauge transformation which is constant on each of the blocks $\Delta(y)$, $y \in \mathbf B$, then the representative of $U^w$ in the axial gauge of $U_0^w$ is $(U^u)^w$.
\end{lemma}

\begin{proof}
By Lemma~\ref{lem:average}(iii), $M^n(U^u)(c) = u(c_-)M^n(U)(c)u(c_+)^{-1}$ for the bonds $c$ of $T^{(n)}$, so that $M^n(U^u)$ is the gauge transform of $M^n(U)$ by the restriction of $u$ to $T^{(n)}$, and by Definition~\ref{def:transports} its transport from $x_{n+1}$ to $x_n$ is $u(x_{n+1})M^n(U)_{\mathfrak t}(x_{n+1}, x_n)u(x_n)^{-1}$. The condition \eqref{eq:axial-gauge} for $U^u$ is therefore equivalent to \eqref{eq:axial-transformation}, which determines $u$ at the points of $T^{(n)}$ in the block from its values at the points of $T^{(n+1)}$, beginning with $u(y) = \id$. A point of $T^{(n+1)}$ in the block is the centre of its own block of order one, where the transport is trivial, so that the values assigned to it at the different levels agree. The transformation so defined equals $\id$ at the points of $\mathbf B$. The locality and the analyticity are those of the averages and of the transports. If $w$ is constant on the blocks, then $M^n(U^w)$ and $M^n(U_0^w)$ are the transforms of $M^n(U)$ and $M^n(U_0)$ by the same constant in each block, both sides of \eqref{eq:axial-gauge} are conjugated by it, and the condition is preserved.
\end{proof}

The next statement is the only quantitative property of the axial gauge. It is \cite[Lemma~1, p.~79]{BalabanREG}, for one level and for either transport rule.

\begin{lemma}[Smallness in the Axial Gauge]\label{lem:axial-smallness}
There are numbers $C_{\mathrm{ax}}$ and $c_{\mathrm{ax}} > 0$, which depend on $L$ and $G$ only, with the following property. Let $1 \leq n \leq k$, let $c$ be a bond of $T^{(n)}$, and let $V_0$ and $V$ be gauge fields on $T^{(n-1)}$ such that
\begin{enumerate}[label=(\roman*)]
\item $\abs{V_0(\partial p) - \id} \leq \alpha_0L^{-2}$ and $\abs{V(\partial p) - \id} \leq \alpha_0L^{-2}$ for the plaquettes $p$ of $T^{(n-1)}$ with corners in $B(c_-) \cup B(c_+)$;
\item $V_{\mathfrak t}(y, x) = V_{0,\mathfrak t}(y, x)$ for $y \in \{ c_-, c_+ \}$ and $x \in B(y)$;
\item $\abs{M(V)(c)M(V_0)(c)^{-1} - \id} \leq \alpha_1$;
\end{enumerate}
with $\alpha_0 + \alpha_1 \leq c_{\mathrm{ax}}$. Then $\abs{V(b)V_0(b)^{-1} - \id} \leq C_{\mathrm{ax}}\alpha_0 + \alpha_1$ for every bond $b$ of $T^{(n-1)}$ with both ends in $B(c_-) \cup B(c_+)$.
\end{lemma}

\subsection{The Landau Gauge}\label{sec:landau-gauge}

\begin{definition}[The Landau Gauge and Averaged Gauge Transformations]\label{def:landau-gauge}
Let $\mathbf U_0$ be a configuration on $T_\eta$ for which the operator $R = R(\mathbf U_0)$ of Definition~\ref{def:operators} exists.
\begin{enumerate}[label=(\roman*)]
\item A configuration $e^{i\eta A}\mathbf U_0$, with a $1$-form $A$ with values in $\mathfrak g^c$, is \emph{in the Landau gauge of $\mathbf U_0$} if $RD^*_{\mathbf U_0}A = 0$.
\item For a gauge transformation $v$ of $T^{(n-1)}$, a configuration $\mathbf V_0$ on $T^{(n-1)}$, and $y \in T^{(n)}$, let
\begin{equation}\label{eq:averaged-gauge-transformation}
    \overline v(y) \coloneqq v(y)\exp\Bigl(\frac{1}{L^4}\sum_{x \in B(y)}\operatorname{Log}\bigl(v(y)^{-1}R(\mathbf V_{0,\mathfrak t}(y, x))v(x)\bigr)\Bigr)\;.
\end{equation}
The \emph{average of order $j$} of a gauge transformation $u$ of $T_\eta$ relative to $\mathbf U_0$ is the gauge transformation $N^j_{\mathbf U_0}(u)$ of $T^{(j)}$ obtained by applying \eqref{eq:averaged-gauge-transformation} $j$ times, at the step $n$ with $\mathbf V_0 = M^{n-1}(\mathbf U_0)$. The transformation $u$ is \emph{admissible} if $N^j_{\mathbf U_0}(u)(y) = \id$ for all $y \in \Gamma_j$, $0 \leq j \leq k$.
\end{enumerate}
\end{definition}

The condition (i) is \cite[(1.27), p.~80]{BalabanREG}, and the averages of (ii) are \cite[(78)--(80), p.~30]{BalabanAVG}; the linear part of $\frac{1}{i}\operatorname{Log}N^j_{\mathbf U_0}(e^{i\lambda})$ in $\lambda$ is $Q'_j(\mathbf U_0)\lambda$. The admissible transformations take the place of the group $\mathcal G_{\mathbf B}$: they are not a group, and for this reason the Landau representative is defined by starting from the representative in the axial gauge \cite[pp.~80--81]{BalabanREG}.

An admissible transformation is not trivial on the determining set, so that it changes the averages of a field on the layers. What it preserves is the average relative to the background, up to an explicit correction on the bonds which join two layers, and this is the form in which the constraint of a step is carried to the Landau gauge \cite[(1.30), (1.31), pp.~81--82]{BalabanREG}.

\begin{definition}[The Relative Average of a Sequence and the Constraint Field]\label{def:constraint-field}
Let $RM \geq 2L$.
\begin{enumerate}[label=(\roman*)]
\item Let $\mathbf U$ be a configuration on $T_\eta$ for which the operators of Definition~\ref{def:operators}(i) are defined, and let $A$ be a $1$-form with values in $\mathfrak g^c$. For a bond $c$ of $\Gamma_j$ put $\mathsf C(A)(c) \coloneqq a_j^{-1}C_j(\mathbf U; a_jA)(c)$, with $C_0 \coloneqq 0$, and
\begin{equation}\label{eq:sequence-average}
    \mathsf X(A) \coloneqq QA + \mathsf C(A)\;,
\end{equation}
a function on the bonds of the layers, so that \eqref{eq:average-of-product-order} reads $M^j(e^{i\eta A}\mathbf U)(c) = v_j(A)(c_-)\,e^{ia_j\mathsf X(A)(c)}\,M^j(\mathbf U)(c)\,v_j(A)(c_+)^{-1}$.
\item Let $V_0$ and $V$ be maps which assign an element of $G^c$ to every bond of every layer, compatibly with the reversal of bonds. For $1 \leq j \leq k$ and an end $y$ of a bond of $\Gamma_j$ put $m_j(y) \coloneqq \id$ if $y \in \Omega_j$, and otherwise
\begin{equation}\label{eq:block-mean-factor}
    m_j(y) \coloneqq \exp\Bigl(\frac{1}{L^4}\sum_{x \in B(y)}\operatorname{Log}\bigl(V_{\mathfrak t}(y, x)V_{0,\mathfrak t}(y, x)^{-1}\bigr)\Bigr)\;,
\end{equation}
where $B(y) \subset T^{(j-1)}$ is the block of $y$ and the transports are those of the restrictions of $V$ and $V_0$ to the bonds with both ends in $B(y)$, which are bonds of $\Gamma_{j-1}$; put $m_0 \coloneqq \id$. The \emph{constraint field} of $V$ relative to $V_0$ is the function $\mathsf b(V, V_0)$ on the positive bonds of the layers with
\begin{equation}\label{eq:constraint-field}
    \mathsf b(V, V_0)(c) \coloneqq \frac{1}{ia_j}\operatorname{Log}\bigl(m_j(c_-)^{-1}\,V(c)\,m_j(c_+)\,V_0(c)^{-1}\bigr)
\end{equation}
for a bond $c$ of $\Gamma_j$, defined when the arguments of the logarithms lie in the domain of the chart.
\end{enumerate}
\end{definition}

An end $y$ of a bond of $\Gamma_j$ which is not in $\Omega_j$ is a point of $T^{(j)}$ within the $\ell^\infty$-distance $a_j$ of $\Omega_j$, and its block consists of points of $T^{(j-1)}$ within the distance $\frac{3}{2}a_j = \frac{3}{2}La_{j-1}$ of $\Omega_j$ which do not lie in $\Omega_j$, because $\Omega_j$ is a union of cubes of the grid of side $a_j$; by \eqref{eq:separation} and $RM \geq 2L$ they lie in $\Omega_{j-1}$, so that the bonds between them are bonds of $\Gamma_{j-1}$. On a bond with both ends in $\Gamma_j$ the constraint field is $(ia_j)^{-1}\operatorname{Log}(V(c)V_0(c)^{-1})$.

\begin{theorem}[From the Axial to the Landau Gauge]\label{thm:landau-gauge}
There are numbers $c_5 > 0$, $B_2$, and $B_2(\beta_0)$ for $0 \leq \beta_0 < 1$, which depend on $L$ and $G$ only, with the following property. Let the sequence of domains be as in Theorem~\ref{thm:landau-operators}, let $U_0 \in \mathfrak U(\alpha_0)$ be regular with the constant $\alpha_0$ in the sense of Definition~\ref{def:regular-pairs}, and let $U \in \mathfrak U(\alpha_0)$ be in the axial gauge of $U_0$ and satisfy $\abs{M^j(U)(c)M^j(U_0)(c)^{-1} - \id} \leq \alpha_1$ for the bonds $c$ of $\Gamma_j$, $0 \leq j \leq k$, with $\alpha_0 + \alpha_1 \leq c_5$. Write $U = U'U_0$. Then there is exactly one admissible gauge transformation $u$ with values in $G$ such that the configuration $U_1$ with
\[
    U_1(b) \coloneqq u(b_-)^{-1}U'(b)R(U_0(b))u(b_+)
\]
on the positive bonds is of the form $e^{i\eta A}$ with a $1$-form $A$ with values in $\mathfrak g$ which satisfies $RD^*_{U_0}A = 0$ and, on $\Omega_j$ for $0 \leq j \leq k$ and with $\epsilon \coloneqq \alpha_0 + \alpha_1$,
\begin{equation}\label{eq:landau-bounds}
\begin{gathered}
    \abs A \leq B_2\epsilon a_j^{-1}\;,
    \qquad
    \abs{\nabla_{U_0}A} \leq B_2\epsilon a_j^{-2}\;,
    \qquad
    \norm{\nabla_{U_0}A}_\beta \leq B_2(\beta_0)\epsilon a_j^{-2-\beta} \quad (\beta \leq \beta_0)\;,\\
    \abs{D^*_{U_0}D_{U_0}A} \leq B_2\epsilon a_j^{-3}\;,
    \qquad
    \abs{\Delta_{U_0}A} \leq B_2\epsilon a_j^{-3}\;.
\end{gathered}
\end{equation}
Moreover $U_1U_0$ is the gauge transform of $U$ by $u^{-1}$, and the relative average of $A$ is determined by the averages of $U$ and of $U_0$ on the layers alone: if $V$ and $V_0$ denote the data with $V(c) = M^j(U)(c)$ and $V_0(c) = M^j(U_0)(c)$ for the bonds $c$ of $\Gamma_j$, then
\begin{equation}\label{eq:landau-constraint}
    \mathsf X(A) = \mathsf b(V, V_0)\;,
    \qquad
    \abs{a_j\mathsf b(V, V_0)(c)} \leq 8L\alpha_1 \text{ for the bonds } c \text{ of } \Gamma_j\;,
\end{equation}
with the relative average \eqref{eq:sequence-average} at $\mathbf U = U_0$.
\end{theorem}

This is \cite[Theorem~2, p.~83]{BalabanREG} with the bounds \cite[(1.36)--(1.39)]{BalabanREG}; the statement on the averages is \cite[(1.30), (1.31), (1.37)]{BalabanREG}, where the factors \eqref{eq:block-mean-factor} are written for the transports along the fixed paths. The transformation $u$ is expressed through $U_1$ by the site factors of Proposition~\ref{prop:relative-average-order} \cite[(106), p.~33]{BalabanAVG}.

\section{The Variational Problem and the Background Field}\label{sec:variational}

The effective action of a step is expanded about a configuration on the fine lattice which is determined by the field of the step: the configuration of least action among the regular fields whose averages on the layers are that field. This section states that such a configuration exists, that it is unique up to the gauge transformations which are trivial on the determining set, that it is as regular as its data, and that it depends analytically on the data \cite{BalabanVAR}. The proof is a perturbation argument about a configuration which is already close to the minimiser, and its objects are needed afterwards in their own right, because the terms of the effective action are functions of them: the nonlinear part of the action in the fluctuation, the change of variables which makes the constraint linear, the potential which remains once the quadratic form has been split off, and the two fixed-point equations for the critical point. We therefore define these objects at a pair, as in Section~\ref{sec:landau-operators}, state their bounds, and express the minimiser through them.

\subsection{The Problem}\label{sec:problem}

The data of the problem are a gauge field on the bonds of the layers, and their regularity is measured by their plaquettes. At the surface between two layers a plaquette of the coarser lattice has bonds which are not bonds of its layer, and on these the data are replaced by their averages \cite[p.~278]{BalabanVAR}.

\begin{definition}[Regular Data]\label{def:regular-data}
Let $RM \geq 2L$ and let $V$ be data on the determining set (Definition~\ref{def:regular-fields}). For $0 \leq j \leq k$ a \emph{plaquette of the data} of the layer $\Gamma_j$ is a plaquette $p'$ of $T^{(j)}$ with a corner in $\Omega_j$ and no corner in $\Omega_{j+1}$. For a bond $c$ in the boundary of $p'$ which has no end in $\Omega_j$ we put $V(c) \coloneqq M(V)(c)$, the average of Definition~\ref{def:average} of the restriction of $V$ to the bonds of $T^{(j-1)}$ with both ends in $B(c_-) \cup B(c_+)$. The data are \emph{regular with the constant $\varepsilon_1$} if
\begin{equation}\label{eq:regular-data}
    \abs{V(\partial p') - \id} \leq \varepsilon_1
\end{equation}
for every plaquette of the data.
\end{definition}

The definition makes sense for the following reason. A bond of $p'$ with an end in $\Omega_j$ has no end in $\Omega_{j+1}$, because no corner of $p'$ lies there, so that it is a bond of $\Gamma_j$. A bond $c$ of $p'$ with no end in $\Omega_j$ occurs only for $j \geq 1$, its ends are points of $T^{(j)}$ within the $\ell^\infty$-distance $a_j$ of a corner in $\Omega_j$, and as in the remark after Definition~\ref{def:constraint-field} the bonds of $T^{(j-1)}$ with both ends in $B(c_-) \cup B(c_+)$ are bonds of $\Gamma_{j-1}$, on which $V$ is given. This is the condition \cite[(7), p.~278]{BalabanVAR}.

\begin{definition}[Critical and Minimal Configurations]\label{def:minimal}
Let $V$ be data, let $0 < \varepsilon_0 \leq L^{-2}c'_4$, and let $U \in \mathfrak U(\varepsilon_0) \cap \mathfrak V(V)$, so that the operator $Q = Q(U)$ of Definition~\ref{def:operators} is defined by Proposition~\ref{prop:relative-average-order}.
\begin{enumerate}[label=(\roman*)]
\item $U$ is \emph{critical} if $\langle\delta A, J(U)\rangle = 0$ for every $1$-form $\delta A$ with values in $\mathfrak g$ and $Q\,\delta A = 0$.
\item $U$ is \emph{minimal in $\mathfrak U(\varepsilon_0) \cap \mathfrak V(V)$} if $A(U) \leq A(U')$ for every $U'$ in this set.
\end{enumerate}
\end{definition}

The linearisation of the constraint $M^j(U) = V$ at $U$ is not $Q$ but the operator which Proposition~\ref{prop:relative-average} calls $Q^{\mathrm{pl}}$, and the two differ by the gauge transformation of the site factors. Since the current is orthogonal to the gauge directions, the two operators give the same notion of a critical configuration, which is the content of the first part of the next lemma.

\begin{lemma}[Critical Configurations]\label{lem:critical}
Let $V$, $\varepsilon_0$, and $U$ be as in Definition~\ref{def:minimal}.
\begin{enumerate}[label=(\roman*)]
\item $U$ is critical if and only if $\frac{\dd}{\dd t}A(U_t)|_{t=0} = 0$ for every continuously differentiable curve $t \mapsto U_t$ in $\mathfrak V(V)$ with $U_0 = U$.
\item If $U$ is minimal in $\mathfrak U(\varepsilon_0) \cap \mathfrak V(V)$ and $U \in \mathfrak U(\varepsilon'_0)$ for some $\varepsilon'_0 < \varepsilon_0$, then $U$ is critical.
\end{enumerate}
\end{lemma}

\begin{proof}[Proof of (ii) from (i)]
The functions $U \mapsto U(\partial p)$ and $U \mapsto J(U)(b)$ are continuous, so that a curve as in (i) lies in $\mathfrak U(\varepsilon_0)$ for small $\abs t$. On this part of the curve $A(U_t) \geq A(U)$ by the minimality, and a differentiable function with a minimum at $t = 0$ has a vanishing derivative there.
\end{proof}

\begin{lemma}[Gauge Transformations of the Problem]\label{lem:problem-covariance}
Let $V$ and $\varepsilon_0$ be as in Definition~\ref{def:minimal}, let $u$ be a gauge transformation of $T_\eta$ with values in $G$, and let $V^u$ be the data with $V^u(c) \coloneqq u(c_-)V(c)u(c_+)^{-1}$. Then $U \mapsto U^u$ is a bijection of $\mathfrak U(\varepsilon_0) \cap \mathfrak V(V)$ onto $\mathfrak U(\varepsilon_0) \cap \mathfrak V(V^u)$ which preserves the action and maps critical configurations to critical ones and minimal configurations to minimal ones, and $V^u$ is regular with the same constant as $V$. In particular the critical and the minimal configurations in $\mathfrak U(\varepsilon_0) \cap \mathfrak V(V)$ form unions of orbits of $\mathcal G_{\mathbf B}$.
\end{lemma}

\begin{proof}
The action and the set $\mathfrak U(\varepsilon_0)$ are invariant under $u$ by Lemma~\ref{lem:action}(ii) and \eqref{eq:current-covariance}, the norm being invariant under conjugation by elements of $G$ (Lemma~\ref{lem:norms}). By Lemma~\ref{lem:average}(iii), $M^j(U^u)(c) = u(c_-)M^j(U)(c)u(c_+)^{-1}$ for every bond $c$ of $T^{(j)}$, so that $U \in \mathfrak V(V)$ if and only if $U^u \in \mathfrak V(V^u)$; the inverse map is the transformation by $u^{-1}$. The same identity for one step, applied to the data, shows that the plaquettes of $V^u$, with the convention of Definition~\ref{def:regular-data}, are the conjugates of those of $V$ by the values of $u$ at a corner. Minimality is preserved because the action is. For criticality, $J(U^u) = R(u)J(U)$ by \eqref{eq:current-covariance} and $Q(U^u)R(u) = R(u)Q(U)$ by Proposition~\ref{prop:relative-average-order}(v), where $R(u)$ acts on a $1$-form at a bond $b$ and on a function on the bonds of the layers at a bond $c$ by the adjoint action of $u(b_-)$ and of $u(c_-)$. Hence $\delta A \mapsto R(u)\delta A$ maps the kernel of $Q(U)$ onto that of $Q(U^u)$, and $\langle R(u)\delta A, J(U^u)\rangle = \langle\delta A, J(U)\rangle$ by the invariance of the trace. The last statement is the case $u \in \mathcal G_{\mathbf B}$, for which $V^u = V$.
\end{proof}

For $u$ constant on the blocks $\Delta(y)$ this contains the covariance \cite[(181), p.~307]{BalabanVAR} of the minimisers in the axial gauge, by the last statement of Lemma~\ref{lem:axial-representative}.

\begin{theorem}[The Minimal Orbit]\label{thm:minimiser}
Let $M_1$ be the number of Theorem~\ref{thm:propagators}, increased if necessary to a power of $L$. There are numbers $R_2 \geq R_1$, $a^{\mathrm v}_0 > 0$, $a^{\mathrm v}_1 > 0$, $B_3 \geq 1$ with $B_3a^{\mathrm v}_1 \leq a^{\mathrm v}_0 \leq L^{-2}c'_4$, and $B_4(\beta_0)$ for $0 \leq \beta_0 < 1$, which depend on $L$ and $G$ only, with the following properties. Let a sequence of domains with the block size $M_1$ and a separation $R \geq R_2$ be given, and let $V$ be data which are regular with a constant $\varepsilon_1 \leq a^{\mathrm v}_1$.
\begin{enumerate}[label=(\roman*)]
\item \emph{Existence.} There is a gauge field $U \in \mathfrak U(B_3\varepsilon_1) \cap \mathfrak V(V)$ which is minimal in $\mathfrak U(a^{\mathrm v}_0) \cap \mathfrak V(V)$.
\item \emph{Uniqueness.} The critical configurations in $\mathfrak U(a^{\mathrm v}_0) \cap \mathfrak V(V)$ form exactly one orbit of $\mathcal G_{\mathbf B}$, the orbit of $U$.
\item \emph{Regularity.} Let $0 \leq j \leq k$, and let $\square$ be a cube which is a union of cubes of the grid of side $M_1a_j$, of a side $s_\square$ with
\begin{equation}\label{eq:minimiser-cube}
    2R_2M_1a_j \leq s_\square\;,
    \qquad
    s_\square\varepsilon_1 \leq 2a^{\mathrm v}_1R_2M_1a_j\;,
\end{equation}
which meets $\Omega_j$ and for which the set of the points at an $\ell^\infty$-distance of at most $2R_2M_1a_j$ from $\square$ is contained in $\Omega_{j-1} \setminus \Omega_{j+1}$, where $\Omega_{-1} \coloneqq \Omega_0$. Then there is a gauge transformation $u$ on $\square$ with values in $G$ such that $U^u = e^{i\eta A}$ on the bonds of $\square$, with a $1$-form $A$ with values in $\mathfrak g$ which satisfies on $\square$
\begin{equation}\label{eq:minimiser-regularity}
\begin{gathered}
    \abs A \leq \frac{s_\square}{a_j}\,\frac{B_3\varepsilon_1}{2}\,a_j^{-1}\;,
    \qquad
    \abs{\nabla_{\id}A} \leq \frac{s_\square}{a_j}\,\frac{B_3\varepsilon_1}{2}\,a_j^{-2}\;,\\
    \abs{D^*_{\id}D_{\id}A} \leq \frac{s_\square}{a_j}\,\frac{B_3\varepsilon_1}{2}\,a_j^{-3}\;,
    \qquad
    \abs{\Delta_{\id}A} \leq \frac{s_\square}{a_j}\,\frac{B_3\varepsilon_1}{2}\,a_j^{-3}\;,
\end{gathered}
\end{equation}
and, for $0 \leq \beta_0 < 1$ and all sites $x$, $x'$ of $\square$ with $\abs{x - x'}_1 \leq a_j$,
\begin{equation}\label{eq:minimiser-holder}
    \abs{(\nabla_{\id}A)(x') - (\nabla_{\id}A)(x)} \leq \frac{s_\square}{a_j}\,B_4(\beta_0)\varepsilon_1a_j^{-2}\Bigl(\frac{\abs{x - x'}_1}{a_j}\Bigr)^{\beta_0}\;.
\end{equation}
\end{enumerate}
In \emph{(iii)} a cube is a cube of $\mathbb R^4$, and the grids, the domains, and the field are replaced by their preimages and lifts under the covering map $\mathbb R^4 \to \mathbb T$.
\end{theorem}

This is \cite[Theorem~1, p.~279]{BalabanVAR}, with the class of cubes of its proof \cite[p.~300]{BalabanVAR}; the number $s_\square/(2a_j)$ is the size $M$ of the cube there, and the second condition \eqref{eq:minimiser-cube} is $M \leq M(\varepsilon_1)$. Four points of the statement are chosen for the use which the later sections make of it. First, a sequence of domains with a block size $M \geq M_1$ and a separation $R$ is also a sequence with the block size $M_1$ and the separation $RM/M_1$, so that the theorem applies to it when $RM \geq R_2M_1$; the operators of Section~\ref{sec:propagators} depend on the domains and not on the block size by which they are described. Secondly, minimality is asserted on the whole of $\mathfrak U(a^{\mathrm v}_0) \cap \mathfrak V(V)$, and not only near $U$. Thirdly, the H\"older bound is stated for the exponents less than one: the Laplacian of $A$ has jumps across the faces of the blocks, where the constraints act, and the second differences of $A$ are not bounded uniformly in $\eta$. Lastly, the cubes are taken in the covering space, so that the statement has a content for the cubes which are larger than the torus.

In the form in which the propagators of Section~\ref{sec:propagators} require it, the regularity reads as follows.

\begin{corollary}[The Minimiser Is Regular]\label{cor:minimiser-regular}
Let a sequence of domains with a block size $M \geq LR_2M_1$ and a separation $R \geq 27$ be given, let $V$ be data which are regular with a constant $\varepsilon_1 \leq a^{\mathrm v}_1$ such that $40M\varepsilon_1 \leq a^{\mathrm v}_1R_2M_1$, and let $U$ be the configuration of Theorem~\ref{thm:minimiser}. Then for every cube $\square$ of the class of the sequence, of an index $j$ (Definition~\ref{def:regular-pairs}(i), in the covering space), there is a gauge transformation $u$ on $\square$ with values in $G$ such that $U^u = e^{i\eta A}$ on the bonds of $\square$ and
\begin{equation}\label{eq:minimiser-regular}
    \abs A \leq \frac{s_\square}{a_j}\,B_3\varepsilon_1a_j^{-1}\;,
    \qquad
    \abs{\nabla_{\id}A} \leq \frac{s_\square}{a_j}\,B_3\varepsilon_1a_j^{-2}\;,
    \qquad
    \abs{D^*_{\id}D_{\id}A} \leq \frac{s_\square}{a_j}\,B_3\varepsilon_1a_j^{-3}
\end{equation}
on $\square$. In particular $U$ is regular to the second order with the constant $B_3\varepsilon_1$.
\end{corollary}

\begin{proof}
The sequence has the separation $RM/M_1 \geq R_2$ with respect to the block size $M_1$, so that Theorem~\ref{thm:minimiser} applies. Let $\square$ be a cube of the class of the index $j$: a union of cubes of the grid of side $Ma_j$ with $10Ma_j \leq s_\square \leq 20Ma_j$, contained in $\Omega_j \setminus \Omega_{j+2}$, which contains a point $x \notin \Omega_{j+1}$. We write $m \coloneqq 2R_2M_1$, so that $m \leq 2M/L$, and $\tilde\square$ for the set of the points within the $\ell^\infty$-distance $ma_j$ of $\square$. Since $M_1$ divides $M$, the cube $\square$ is a union of cubes of the grid of side $M_1a_j$.

Suppose first that $\tilde\square$ does not meet $\Omega_{j+1}$. The points of $\tilde\square$ are within the distance $ma_j \leq 2Ma_{j-1}$ of $\square \subset \Omega_j$, which is less than $RMa_{j-1}$, so that $\tilde\square \subset \Omega_{j-1}$ by \eqref{eq:separation}. The conditions \eqref{eq:minimiser-cube} hold because $s_\square \geq 10Ma_j \geq ma_j$ and $s_\square\varepsilon_1 \leq 20M\varepsilon_1a_j \leq 2a^{\mathrm v}_1R_2M_1a_j$. Hence Theorem~\ref{thm:minimiser}(iii) applies to $\square$ with the index $j$ and gives \eqref{eq:minimiser-regular} with half the constant.

Suppose next that $\tilde\square$ contains a point $z \in \Omega_{j+1}$, so that $j < k$. Let $\square'$ be a cube which is a union of cubes of the grid of side $M_1a_{j+1}$, contains $\square$ and $z$, and has the smallest side $s'$ for which this is possible. The set $\square \cup \{ z \}$ lies in a box with sides of at most $s_\square + ma_j$, and a box with sides of at most $\ell$ lies in a cube of the grid of side $g$ whose side is the smallest multiple of $g$ not less than $\ell + g$; hence
\[
    s_\square \leq s' \leq s_\square + ma_j + 2M_1a_{j+1} \leq s_\square + \frac{2M}{L}a_j + 2M_1La_j \leq s_\square + 4Ma_j \leq \frac{7}{5}s_\square\;,
\]
where we used $M_1L \leq M$ and $s_\square \geq 10Ma_j$. We verify the hypotheses of Theorem~\ref{thm:minimiser}(iii) for $\square'$ with the index $j + 1$. The cube meets $\Omega_{j+1}$ at $z$. The first condition \eqref{eq:minimiser-cube} holds since $s' \geq 10Ma_j \geq 10LR_2M_1a_j = 5ma_{j+1}$, and the second since $s' \leq 24Ma_j = 24ML^{-1}a_{j+1}$ and $12M\varepsilon_1 \leq a^{\mathrm v}_1R_2M_1$. Every point of the set $\tilde\square'$ of the points within the distance $ma_{j+1} \leq 2Ma_j$ of $\square'$ is within the distance $s' + 2Ma_j \leq 26Ma_j$ of $z$ and of $x$. Since $z \in \Omega_{j+1}$ and $26M < RM$, the set $\tilde\square'$ lies in $\Omega_j$ by \eqref{eq:separation}; and since $x \notin \Omega_{j+1}$ and $26Ma_j < RMa_{j+1}$, it does not meet $\Omega_{j+2}$. The theorem gives a gauge transformation on $\square'$ and a $1$-form with the bounds \eqref{eq:minimiser-regularity} in which $s_\square/a_j$ is replaced by $s'/a_{j+1}$ and $a_j^{-p}$ by $a_{j+1}^{-p}$, $p = 1, 2, 3$. Their ratio to the bounds \eqref{eq:minimiser-regular} is $\frac{1}{2}(s'/s_\square)L^{-1-p} \leq \frac{7}{10}L^{-2} \leq 1$, and we restrict the transformation to $\square$.

In both cases the form has values in $\mathfrak g$, which is the condition of Definition~\ref{def:regular-pairs}(ii) with the constant $B_3\varepsilon_1$, to the second order.
\end{proof}

\subsection{The Expansion about a Pair}\label{sec:expansion}

Throughout this subsection and the next a sequence of domains as in Theorem~\ref{thm:landau-operators} is given, and $(\mathbf U, \mathbf J)$ is a regular pair with constants $(\alpha_0, \alpha_1, \gamma)$ which satisfy the hypothesis of that theorem, so that the operators of Sections~\ref{sec:forms}--\ref{sec:fluctuation-propagator} exist at the pair. The arguments of the functionals below are $1$-forms on $T_\eta$ with values in $\mathfrak g^c$, which we write $A$, $A'$, and $\mathfrak A$ as in \cite{BalabanVAR}; the complex part of $\mathbf U$ in Definition~\ref{def:regular-pairs} does not appear explicitly in this section. A function of such a form and of the pair is called \emph{holomorphic} if it is holomorphic in the form and, for a fixed regular gauge field, in the complex part of $\mathbf U$ and in $\mathbf J$, as in Theorem~\ref{thm:landau-operators}(i).

\begin{definition}[Scaled Norms]\label{def:scaled-norms}
For a $1$-form $A$, a point $y \in \Gamma_j$, and a function $\mu$ on the bonds of the layers,
\begin{equation}\label{eq:scaled-norms}
\begin{gathered}
    \abs A_* \coloneqq \max\bigl\{ \abs A_{(-1)}, \abs{\nabla_{\mathbf U}A}_{(-2)} \bigr\}\;,
    \qquad
    \norm A_y \coloneqq a_j\abs A_{\Delta(y)} + a_j^2\abs{\nabla_{\mathbf U}A}_{\Delta(y)}\;,\\
    \abs\mu_{(\alpha)} \coloneqq \sup_ca_{j(c)}^{-\alpha}\abs{\mu(c)}\;,
\end{gathered}
\end{equation}
with the weighted norms of Definition~\ref{def:norms}(iii), and with $j(c) \coloneqq j$ for a bond $c$ of $\Gamma_j$.
\end{definition}

The first object is the change of variables which makes the constraint linear. The relative average $\mathsf X(A) = QA + \mathsf C(A)$ of \eqref{eq:sequence-average} has a remainder of the second order, and the remainder is removed by subtracting from the field the minimiser of the quadratic form whose average is the remainder \cite[Sect.~C]{BalabanVAR}.

\begin{proposition}[The Linearising Transformation]\label{prop:linearising}
There are numbers $a_3 > 0$, $\delta_2 \in (0, \delta_1]$, and $C_5$, which depend on $L$ and $G$ only, such that for every $1$-form $A'$ with values in $\mathfrak g^c$ and $\abs{A'}_{(-1)} \leq a_3$ the following holds.
\begin{enumerate}[label=(\roman*)]
\item There is exactly one function $\mathsf D = \mathsf D(A')$ on the bonds of the layers with $\abs{\mathsf D}_{(-1)} \leq a_3$ and
\begin{equation}\label{eq:linearising-equation}
    \mathsf D(A') = \mathsf C\bigl(A' - H\mathsf D(A')\bigr)\;.
\end{equation}
It is holomorphic, it vanishes on the bonds of $\Gamma_0$, and $\abs{\mathsf D(A')}_{(-1)} \leq C_5\abs{A'}_{(-1)}^2$.
\item The \emph{linearising transformation} $T(A') \coloneqq A' - H\mathsf D(A')$ satisfies $\mathsf X(T(A')) = QA'$ and $\abs{T(A') - A'}_* \leq C_5\abs{A'}_{(-1)}^2$. It is injective on the set of the $A'$ with $\abs{A'}_{(-1)} \leq a_3$, and for every $A$ with $\abs A_{(-1)} \leq \frac{1}{4}a_3$ there is an $A'$ in this set with $T(A') = A$, $\abs{A'}_{(-1)} \leq 2\abs A_{(-1)}$, and $\abs{A'}_* \leq 2\abs A_*$.
\item Let $\mathfrak D(A')$ be the derivative of $\mathsf D$ at $A'$, a linear map from $1$-forms to functions on the bonds of the layers. For a $1$-form $\delta A$ supported in $\Delta(y)$, $y \in \mathbf B$, and a bond $c$ of the layers,
\begin{equation}\label{eq:linearising-derivative}
    a_{j(c)}\abs{(\mathfrak D(A')\delta A)(c)} \leq C_5\abs{A'}_{(-1)}e^{-\delta_2d(c_-, y)}a_{j(y)}\abs{\delta A}\;,
\end{equation}
and the derivative of $\mathsf D_3 \coloneqq \mathsf D - \mathsf C^{(2)}$ satisfies the same bound with $\abs{A'}_{(-1)}^2$ in the place of $\abs{A'}_{(-1)}$.
\end{enumerate}
\end{proposition}

This is \cite[Proposition~3, p.~289]{BalabanVAR} with \cite[(49)--(62), (73)]{BalabanVAR}, in the normalisation of Definition~\ref{def:landau-operators}, in which $QH = I$; the function $D$ and the operator $H$ of \cite{BalabanVAR} are $a_{j(\cdot)}\mathsf D$ and $Ha_{j(\cdot)}^{-1}$. The identity $\mathsf X(T(A')) = QA'$ is \eqref{eq:linearising-equation}: $QT(A') = QA' - \mathsf D(A')$ and $\mathsf C(T(A')) = \mathsf D(A')$.

Next we split the action. The second variation $\Delta(\mathbf U)$ and the form $\Delta_\pi \coloneqq \Delta(\mathbf U) - \Delta'_\pi$ agree on the forms in the Landau gauge when the second entry of the pair is the current, and \cite{BalabanVAR} writes $V$ for the potential formed with either. Their gradients differ, however, and the two fixed-point equations of Section~\ref{sec:critical-point} need one each, so that we keep both.

\begin{definition}[The Potentials]\label{def:potential}
\begin{enumerate}[label=(\roman*)]
\item $\mathsf V_0(A) \coloneqq A(e^{i\eta A}\mathbf U) - A(\mathbf U) - \langle A, J(\mathbf U)\rangle - \frac{1}{2}\langle A, \Delta(\mathbf U)A\rangle$, an entire function of the $1$-form $A$ which vanishes to the third order at $A = 0$ by Definition~\ref{def:hessian}.
\item For $\abs{A'}_{(-1)} \leq a_3$,
\begin{equation}\label{eq:potential}
    \mathsf V(A') \coloneqq -\langle H\mathsf D_3(A'), \mathbf J\rangle - \langle A', \Delta(\mathbf U)H\mathsf D(A')\rangle + \frac{1}{2}\langle H\mathsf D(A'), \Delta(\mathbf U)H\mathsf D(A')\rangle + \mathsf V_0\bigl(T(A')\bigr)\;,
\end{equation}
and $\mathsf V_\pi(A')$ is defined by the same formula with $\Delta_\pi$ in the place of $\Delta(\mathbf U)$.
\item The \emph{gradients} $\mathsf N(A')$ and $\mathsf N_\pi(A')$ are the $1$-forms with $\langle\delta A, \mathsf N(A')\rangle = \frac{\dd}{\dd t}\mathsf V(A' + t\,\delta A)|_{t=0}$ for all $1$-forms $\delta A$, and likewise for $\mathsf V_\pi$.
\end{enumerate}
\end{definition}

The potential $\mathsf V_\pi$ is \cite[(80), p.~290]{BalabanVAR}, and $\mathsf V$ is the functional which appears, under the same letter, in \cite[(127), (128), p.~297]{BalabanVAR} and in the equations derived from them.

\begin{lemma}[The Action in the New Variable]\label{lem:functional}
Let $\abs{A'}_{(-1)} \leq a_3$ and $A \coloneqq T(A')$. Then
\begin{equation}\label{eq:functional}
    A(e^{i\eta A}\mathbf U) + \langle A, \mathbf J - J(\mathbf U)\rangle = A(\mathbf U) + \langle A', \mathbf J\rangle + \frac{1}{2}\bigl\langle A', \bigl(\Delta(\mathbf U) - \Delta^{(2)}\bigr)A'\bigr\rangle + \mathsf V(A')\;,
\end{equation}
and
\begin{equation}\label{eq:potential-difference}
    \mathsf V(A') - \mathsf V_\pi(A') = \langle A', \Delta'_\pi H\mathsf D(A')\rangle - \frac{1}{2}\langle H\mathsf D(A'), \Delta'_\pi H\mathsf D(A')\rangle\;.
\end{equation}
If $\mathbf J = J(\mathbf U)$ and $RD^*A' = 0$, then $RD^*A = 0$, the right side of \eqref{eq:potential-difference} vanishes, and $\langle A', \Delta^{(2)}A'\rangle = \langle A', \Delta^{(2)}_\pi A'\rangle$.
\end{lemma}

\begin{proof}
By the definition of $\mathsf V_0$ the left side of \eqref{eq:functional} equals $A(\mathbf U) + \langle A, \mathbf J\rangle + \frac{1}{2}\langle A, \Delta(\mathbf U)A\rangle + \mathsf V_0(A)$. We insert $A = A' - H\mathsf D(A')$ and $\mathsf D = \mathsf C^{(2)} + \mathsf D_3$. The linear term becomes $\langle A', \mathbf J\rangle - \langle H\mathsf C^{(2)}(A'), \mathbf J\rangle - \langle H\mathsf D_3(A'), \mathbf J\rangle$, and $\langle H\mathsf C^{(2)}(A'), \mathbf J\rangle = \frac{1}{2}\langle A', \Delta^{(2)}A'\rangle$ by Definition~\ref{def:landau-operators}(iii). The quadratic term becomes $\frac{1}{2}\langle A', \Delta(\mathbf U)A'\rangle - \langle A', \Delta(\mathbf U)H\mathsf D(A')\rangle + \frac{1}{2}\langle H\mathsf D(A'), \Delta(\mathbf U)H\mathsf D(A')\rangle$, because $\Delta(\mathbf U)$ is symmetric. The sum is the right side of \eqref{eq:functional}, by \eqref{eq:potential}. The difference \eqref{eq:potential-difference} is read off from the two definitions, since $\Delta(\mathbf U) - \Delta_\pi = \Delta'_\pi$.

Let now $\mathbf J = J(\mathbf U)$ and $RD^*A' = 0$. Then $RD^*H = 0$ by Theorem~\ref{thm:landau-operators}(iv), so that $RD^*A = 0$, and the gauge functions $\lambda_{A'} = G'RD^*A'$ and $\lambda_{H\mu} = G'RD^*H\mu$ of Definition~\ref{def:landau-operators} vanish. Every term of the form \eqref{eq:gauge-defect} contains a factor $\lambda$, and $\lambda_{A_1 + A_2} = \lambda_{A_1} + \lambda_{A_2}$, so that the bilinear form of $\Delta'_\pi$, which is obtained from $\langle X(A), \mathbf J\rangle$ by polarisation, vanishes on two forms whose gauge functions vanish. This gives the second statement. For the third, $\mathcal PA' = A' - D\lambda_{A'} = A'$, hence $\langle A', \Delta^{(2)}_\pi A'\rangle = \langle\mathcal PA', \Delta^{(2)}\mathcal PA'\rangle = \langle A', \Delta^{(2)}A'\rangle$.
\end{proof}

At $\mathbf J = J(\mathbf U)$ the left side of \eqref{eq:functional} is the action of $e^{i\eta A}\mathbf U$, and \eqref{eq:functional} is its expansion in the variable in which the constraint $\mathsf X(A) = \mathfrak b$ is the linear condition $QA' = \mathfrak b$. This is \cite[(74)--(81), pp.~289--290]{BalabanVAR}.

\begin{proposition}[The Potentials]\label{prop:potential}
The numbers $a_3$ and $\delta_2$ of Proposition~\ref{prop:linearising} can be chosen, and there are numbers $C_4$ and $C_N$ which depend on $L$ and $G$ only, such that for every $1$-form $A'$ with values in $\mathfrak g^c$ and $\abs{A'}_* \leq a_3$:
\begin{enumerate}[label=(\roman*)]
\item $\mathsf V$, $\mathsf V_\pi$, $\mathsf N$, and $\mathsf N_\pi$ are holomorphic, and $\mathsf V$ and $\mathsf V_\pi$ vanish to the third order at $A' = 0$.
\item $\abs{\mathsf N(A')}_{(-3)} \leq C_4\abs{A'}_*^2$.
\item For a bond $b$ of $\Delta(y)$, $y \in \Gamma_j$,
\begin{equation}\label{eq:gradient-kernel}
    a_j^3\abs{\mathsf N(A')(b)} \leq C_N\abs{A'}_*\sum_{y'' \in \mathbf B}e^{-\delta_2d(y, y'')}\norm{A'}_{y''}\;.
\end{equation}
\item For a $1$-form $\mathfrak A$ supported in $\tilde\Delta(y')$, $y' \in \mathbf B$, and a bond $b$ of $\Delta(y)$, $y \in \Gamma_j$, the derivative of $\mathsf N$ at $A'$ in the direction $\mathfrak A$ satisfies
\begin{equation}\label{eq:gradient-derivative}
    a_j^3\Bigl\lvert\Bigl(\frac{\dd}{\dd t}\mathsf N(A' + t\mathfrak A)\Big|_{t=0}\Bigr)(b)\Bigr\rvert \leq C_N\abs{A'}_*e^{-\delta_2d(y, y')}\abs{\mathfrak A}_*\;.
\end{equation}
\end{enumerate}
The bounds \emph{(ii)}--\emph{(iv)} hold for $\mathsf N_\pi$ as well.
\end{proposition}

The bound (ii) is \cite[Proposition~4, p.~292]{BalabanVAR}, and (iv) is \cite[(189), p.~308]{BalabanVAR}, which is printed there without its proof. The form (iii), in which one factor of the field is kept at its place, is not in \cite{BalabanVAR}; it is what the localisation of the minimiser needs.

\subsection{The Equations of the Critical Point}\label{sec:critical-point}

A critical point of the action under the constraint and the gauge condition is a critical point of the functional \eqref{eq:functional} on an affine subspace, and the equation for it becomes a fixed-point equation once the quadratic form is inverted on that subspace. This can be done in two ways. In the first the subspace is that of the forms with given averages in the Landau gauge, the quadratic form is that of $G_1$, and the inverse is the propagator $\mathfrak G$ of the fluctuation \cite[(99)--(111)]{BalabanVAR}. In the second the gauge condition is dropped from the subspace, which is legitimate because the action is gauge invariant, and the inverse is formed with the operator $G$ of Definition~\ref{def:operators}, corrected by the form $\Delta^{(2)}$ \cite[(127)--(143), (179), (180)]{BalabanVAR}. The second equation is the one with the better regularity, because $G$ is bounded together with its covariant Laplacian, and it is the one through which the later sections write the background field.

\begin{definition}[The Operators of the Second Equation]\label{def:second-operators}
With the operator $G$ of Definition~\ref{def:operators} and the form $\Delta^{(2)}$ of Definition~\ref{def:landau-operators},
\begin{equation}\label{eq:second-operators}
    H_0 \coloneqq GQ^*(QGQ^*)^{-1}\;,
    \qquad
    G_{\mathrm c} \coloneqq \bigl(\Delta_{\mathsf a} - \Delta^{(2)}\bigr)^{-1}\;,
    \qquad
    \tilde G \coloneqq G_{\mathrm c} - G_{\mathrm c}Q^*(QG_{\mathrm c}Q^*)^{-1}QG_{\mathrm c}\;.
\end{equation}
\end{definition}

By their definitions $QH_0 = I$, $Q\tilde G = 0$, and $\tilde GQ^* = 0$. In \cite[(129), p.~297]{BalabanVAR} the operator $H_0$ is formed with $G$, and in \cite[(180), p.~306]{BalabanVAR} the operator $\tilde G$ is formed with the corrected operator, which is there again called $G$.

\begin{proposition}[The Operators of the Second Equation]\label{prop:second-operators}
Under the hypotheses of Theorem~\ref{thm:landau-operators}, with smaller numbers $a_2$ and $\delta_1$ and larger constants $B_1$, $B_1(\beta)$, $B'_1(\epsilon, \beta)$ if necessary:
\begin{enumerate}[label=(\roman*)]
\item The operators \eqref{eq:second-operators} exist and are holomorphic in the pair, and $G_{\mathrm c} = \sum_{n \geq 0}G(\Delta^{(2)}G)^n$.
\item $G_{\mathrm c}$ and $\tilde G$ satisfy the bounds \emph{(iv)}--\emph{(vi)} of Theorem~\ref{thm:propagators} for $1$-forms, the bound of the covariant Laplacian included, and $(QGQ^*)^{-1}$ and $(QG_{\mathrm c}Q^*)^{-1}$ satisfy \eqref{eq:block-form-inverse}.
\item The kernel of $H_0$ satisfies \eqref{eq:minimiser-kernel} and, for $x \in \Delta(y)$, $y \in \Gamma_j$, and a bond $c'$ of $\Gamma_{j'}$,
\begin{equation}\label{eq:second-minimiser-laplacian}
    a_j^2\abs{(\Delta_{\mathbf U}H_0)(x, c')} \leq B_1a_{j'}^{-4}e^{-\frac{1}{2}\delta_1d(y, c'_-)}\;.
\end{equation}
\item If $\mathbf J = J(\mathbf U)$ and $\mathbf U$ is a gauge field in $\mathfrak U(\alpha_0)$, then for every function $\mu$ on the bonds of the layers
\begin{equation}\label{eq:minimiser-second-order}
    \abs{D^*DH\mu}_{(-3)} + \abs{\Delta_{\mathbf U}H\mu}_{(-3)} \leq B_1\abs\mu_{(-1)}\;,
\end{equation}
and $G_{\mathrm c}$, $\tilde G$, and $QG_{\mathrm c}Q^*$ are symmetric.
\end{enumerate}
\end{proposition}

These are \cite[(129), (130), (137)--(140), pp.~297--299]{BalabanVAR} and the remark \cite[p.~306]{BalabanVAR} that the corrected operator ``has exactly the same properties'' as $G$. The bound \eqref{eq:minimiser-second-order} is derived there from the identity $RD^*H = 0$, and we do not state it at a pair whose second entry is not the current of the first.

\begin{proposition}[The Two Fixed Points]\label{prop:fixed-points}
There are numbers $a_4 \in (0, \frac{1}{2}a_3]$ and $C_6 \geq 1$, which depend on $L$ and $G$ only, with the following property. Let $\gamma \leq a_4/C_6$, and let $\mathfrak A$ be a $1$-form with values in $\mathfrak g^c$ and $\abs{\mathfrak A}_* \leq a_4$. Then each of the maps
\begin{equation}\label{eq:fixed-point-maps}
    \Phi_1(A) \coloneqq -\mathfrak G\mathbf J - \mathfrak G\mathsf N_\pi(A + \mathfrak A)\;,
    \qquad
    \Phi_0(A) \coloneqq -\tilde G\mathbf J + \tilde G\Delta^{(2)}\mathfrak A - \tilde G\mathsf N(A + \mathfrak A)
\end{equation}
maps the set of the $1$-forms $A$ with values in $\mathfrak g^c$ and $\abs A_* \leq a_4$ into itself and satisfies $\abs{\Phi(A) - \Phi(A'')}_* \leq \frac{1}{2}\abs{A - A''}_*$ on it. Consequently each has exactly one fixed point in this set, $\mathcal A_1(\mathfrak A)$ and $\mathcal A_0(\mathfrak A)$, respectively. The fixed points are holomorphic in $\mathfrak A$ and in the pair, and
\begin{equation}\label{eq:fixed-point-bound}
    \abs{\mathcal A_1(\mathfrak A)}_* \leq C_6\bigl(\gamma + \abs{\mathfrak A}_*^2\bigr)\;,
    \qquad
    \abs{\mathcal A_0(\mathfrak A)}_* + \abs{\Delta_{\mathbf U}\mathcal A_0(\mathfrak A)}_{(-3)} \leq C_6\bigl(\gamma + \abs{\mathfrak A}_*^2\bigr)\;.
\end{equation}
Moreover $Q\mathcal A_0(\mathfrak A) = 0$, and if $\mathbf J = J(\mathbf U)$, then $Q\mathcal A_1(\mathfrak A) = 0$ and $RD^*\mathcal A_1(\mathfrak A) = 0$.
\end{proposition}

This is \cite[Proposition~6, p.~295]{BalabanVAR} with \cite[(116)--(121)]{BalabanVAR} for the first map, and the statement of \cite[p.~300]{BalabanVAR} that the second equation ``has all the properties'' of the first. The terms $\mathfrak G\mathbf J$ and $\tilde G\mathbf J$ are kept in the equations, since the functions defined by them are used at pairs at which they do not vanish.

\begin{definition}[The Two Solutions]\label{def:background-field}
Let $B_H$ be a number such that $\abs{H_0\mu}_* \leq B_H\abs\mu_{(-1)}$ and $\abs{H_1\mu}_* \leq B_H\abs\mu_{(-1)}$ for all functions $\mu$ on the bonds of the layers, which exists and depends on $L$ and $G$ only by \eqref{eq:minimiser-kernel}, Lemma~\ref{lem:summation}, and the remark after Theorem~\ref{thm:propagators} on the powers of the spacing. For $\gamma \leq a_4/C_6$ and a function $\mathfrak b$ on the bonds of the layers with values in $\mathfrak g^c$ and $B_H\abs{\mathfrak b}_{(-1)} \leq a_4$ we put
\begin{equation}\label{eq:background-field}
    \mathcal H_1(\mathfrak b) \coloneqq T\bigl(\mathcal A_1(H_1\mathfrak b) + H_1\mathfrak b\bigr)\;,
    \qquad
    \mathcal H(\mathfrak b) \coloneqq T\bigl(\mathcal A_0(H_0\mathfrak b) + H_0\mathfrak b\bigr)\;,
\end{equation}
with the linearising transformation of Proposition~\ref{prop:linearising}. Both are functions of the pair as well, and we write $\mathcal H(\mathbf U, \mathbf J; \mathfrak b)$ where the pair has to be named.
\end{definition}

The arguments of $T$ have the norm $\abs\cdot_{(-1)}$ at most $2a_4 \leq a_3$, so that both functions are defined, and they are holomorphic in $\mathfrak b$ and in the pair. By Proposition~\ref{prop:linearising}(ii) and $QH_0 = QH_1 = I$ they satisfy $\mathsf X(\mathcal H(\mathfrak b)) = \mathfrak b$ at every pair, and $\mathsf X(\mathcal H_1(\mathfrak b)) = \mathfrak b$ when $\mathbf J = J(\mathbf U)$. The representations \eqref{eq:background-field} are \cite[(174), (175), p.~305]{BalabanVAR} and \cite[(179), (180), p.~306]{BalabanVAR}.

\begin{theorem}[The Critical Configurations near a Regular Field]\label{thm:critical}
There are numbers $a_5 > 0$ and $B_5$, which depend on $L$ and $G$ only, with the following properties. Let $U_0 \in \mathfrak U(\alpha_0)$ be a gauge field which is regular to the second order with the constant $\alpha_0$, where $(M + 1)\alpha_0 \leq a_5$, and let the operators and the functions of this section be taken at the pair $(U_0, J(U_0))$.
\begin{enumerate}[label=(\roman*)]
\item For every function $\mathfrak b$ on the bonds of the layers with values in $\mathfrak g$ and $\abs{\mathfrak b}_{(-1)} \leq a_5$, the forms $\mathcal H(\mathfrak b)$ and $\mathcal H_1(\mathfrak b)$ are equal, have values in $\mathfrak g$, and satisfy $RD^*\mathcal H(\mathfrak b) = 0$ and $\mathsf X(\mathcal H(\mathfrak b)) = \mathfrak b$. The gauge field $e^{i\eta\mathcal H(\mathfrak b)}U_0$ is critical, and
\begin{equation}\label{eq:critical-bounds}
    \abs{\mathcal H(\mathfrak b)}_* + \abs{D^*D\mathcal H(\mathfrak b)}_{(-3)} + \abs{\Delta_{U_0}\mathcal H(\mathfrak b)}_{(-3)} \leq B_5\bigl(\alpha_0 + \abs{\mathfrak b}_{(-1)}\bigr)\;.
\end{equation}
\item Conversely, if $A$ is a $1$-form with values in $\mathfrak g$ such that $RD^*A = 0$, $\abs A_* \leq a_5$, and $e^{i\eta A}U_0$ is critical, then $A = \mathcal H(\mathsf X(A))$.
\item If $U_0$ is critical, then $\mathfrak GJ(U_0) = 0$ and $\tilde GJ(U_0) = 0$, and $\mathcal H(0) = 0$.
\end{enumerate}
\end{theorem}

This collects \cite[Propositions~2, 5, and~7]{BalabanVAR}, in a form in which the configurations in the Landau gauge are not required to have prescribed averages: a field $e^{i\eta A}U_0$ is critical for the data which are its own averages, in the sense of Definition~\ref{def:minimal}, and by Lemma~\ref{lem:problem-covariance} this property does not depend on the gauge. The statement (iii) is \cite[(170), (171), (178)]{BalabanVAR}, and it follows from the definitions: if $U_0$ is critical, then $J(U_0)$ is orthogonal to the kernel of $Q$, which contains the ranges of $\mathfrak P$ and of the projection $I - G_{\mathrm c}Q^*(QG_{\mathrm c}Q^*)^{-1}Q$, while $\mathfrak G = G_1\mathfrak P^T$ by Theorem~\ref{thm:fluctuation-propagator}(ii) and $\tilde G$ is $G_{\mathrm c}$ times the transpose of that projection; and then $A = 0$ is the fixed point of $\Phi_0$ at $\mathfrak A = 0$, because $\mathsf N(0) = 0$.

\subsection{The Background Field as a Function of the Data}\label{sec:background-field}

We can now say how the minimiser moves when the data are changed. The minimisers of two nearby sets of data differ by a fluctuation which, in the Landau gauge of the first, is the function $\mathcal H$ at the constraint field of the two sets of data, and this function extends to complex data \cite[Sect.~G]{BalabanVAR}.

\begin{theorem}[The Background Field]\label{thm:background-field}
There are numbers $C_{\mathrm d} \geq 1$, $C'_{\mathrm d} \geq 1 + 4C_{\mathrm d}$, $B_6$, and $B_6(\beta_0)$ for $0 \leq \beta_0 < 1$, which depend on $L$ and $G$ only, such that the following holds, with smaller numbers $a^{\mathrm v}_1$ and $\delta_2$ in Theorem~\ref{thm:minimiser} and Proposition~\ref{prop:linearising} if necessary. Let the sequence of domains, with its block size $M$ and its separation $R$, be as in Corollary~\ref{cor:minimiser-regular}, let $V_0$ be data which are regular with a constant $\varepsilon_1$ such that $C'_{\mathrm d}\varepsilon_1$ satisfies the two conditions which that corollary imposes on $\varepsilon_1$, let $U_0$ be a minimal configuration for $V_0$ as in Theorem~\ref{thm:minimiser}, and let the functions of Section~\ref{sec:critical-point} be taken at the pair $(U_0, J(U_0))$. For a function $B$ on the positive bonds of the layers with values in $\mathfrak g^c$ and $\abs{B(c)} \leq C_{\mathrm d}\varepsilon_1$ for all $c$, let $V_B$ be the map with $V_B(c) \coloneqq e^{iB(c)}V_0(c)$, and $\mathfrak b_B \coloneqq \mathsf b(V_B, V_0)$ its constraint field \eqref{eq:constraint-field}.
\begin{enumerate}[label=(\roman*)]
\item \emph{The constraint field.} $\mathfrak b_B$ is defined and holomorphic in $B$, it vanishes at $B = 0$, and $a_j\abs{\mathfrak b_B(c)} \leq 8L\sup\abs B$ for a bond $c$ of $\Gamma_j$. The value $\mathfrak b_B(c)$ depends on $B(c)$ and, if an end $y$ of $c$ does not lie in $\Omega_j$, on $B$ on the bonds with both ends in $B(y)$.
\item \emph{The fluctuation.} $\mathcal H(\mathfrak b_B)$ is defined and holomorphic in $B$, it equals $\mathcal H_1(\mathfrak b_B)$, it vanishes at $B = 0$, and
\begin{equation}\label{eq:background-bounds}
    \abs{\mathcal H(\mathfrak b_B)}_* + \abs{D^*D\mathcal H(\mathfrak b_B)}_{(-3)} + \abs{\Delta_{U_0}\mathcal H(\mathfrak b_B)}_{(-3)} \leq B_6\sup\abs B\;,
    \qquad
    RD^*\mathcal H(\mathfrak b_B) = 0\;.
\end{equation}
\item \emph{The minimiser.} There is a gauge transformation $w_B$ of $T_\eta$ with values in $G^c$, holomorphic in $B$, with $w_0 = \id$ and with values in $G$ for $B$ with values in $\mathfrak g$, such that the configuration
\begin{equation}\label{eq:moved-minimiser}
    U(B) \coloneqq \bigl(e^{i\eta\mathcal H(\mathfrak b_B)}U_0\bigr)^{w_B}
\end{equation}
is in the axial gauge of $U_0$ and satisfies $M^j(U(B))(c) = V_B(c)$ for the bonds $c$ of $\Gamma_j$, $0 \leq j \leq k$. For $B$ with values in $\mathfrak g$ the transformation $w_B$ is admissible (Definition~\ref{def:landau-gauge}), the data $V_B$ are regular with the constant $C'_{\mathrm d}\varepsilon_1$, and $U(B)$ is the configuration of Theorem~\ref{thm:minimiser} for $V_B$ which is in the axial gauge of $U_0$.
\item \emph{The derivative.} Let $c'$ be a bond of $\Gamma_{j'}$, and let $K$ be the derivative of $\mathcal H(\mathfrak b_B)$ with respect to $B$ in a direction $\delta B$ which is supported on $c'$. Then, for $x \in \Delta(y)$, $y \in \Gamma_j$, a real function $\zeta$ supported in $\tilde\Delta(y)$, and $0 \leq \beta \leq \beta_0 < 1$,
\begin{equation}\label{eq:background-derivative}
\begin{aligned}
    a_j\abs{K(x)} + a_j^2\abs{(\nabla_{U_0}K)(x)} &\leq B_6e^{-\delta_2d(y, c'_-)}\abs{\delta B}\;,\\
    a_j^3\abs{(D^*DK)(x)} + a_j^3\abs{(\Delta_{U_0}K)(x)} &\leq B_6e^{-\delta_2d(y, c'_-)}\abs{\delta B}\;,\\
    a_j^{2+\beta}\norm{\zeta\nabla_{U_0}K}_\beta &\leq B_6(\beta_0)\bigl(\norm\zeta_\beta^{(j)} + \abs\zeta\bigr)e^{-\delta_2d(y, c'_-)}\abs{\delta B}\;.
\end{aligned}
\end{equation}
\end{enumerate}
\end{theorem}

This is \cite[Proposition~9, p.~309]{BalabanVAR}: the smallness of the data is \cite[(172)]{BalabanVAR}, the bound \eqref{eq:background-bounds} is \cite[(173)]{BalabanVAR} with \cite[(19)--(21), p.~281]{BalabanVAR}, in the form which is linear in $B$ and which follows from the bound by a constant and the vanishing at $B = 0$, and \eqref{eq:background-derivative} is \cite[(190), p.~308]{BalabanVAR}. For real $B$ the action of the minimiser for $V_B$ is $A(e^{i\eta\mathcal H(\mathfrak b_B)}U_0)$, by (iii) and the invariance of the action, so that the least action is a real analytic function of the data. The expansion of $\mathcal H(\mathfrak b)$ in powers of $\mathfrak b$ is obtained by iterating the fixed-point equations; by Theorem~\ref{thm:critical}(iii) and the vanishing of $\mathsf N$ to the second order it begins with $H_1\mathfrak b$ \cite[(176), (177), p.~306]{BalabanVAR}.

\subsection{The Problem on the Support of a Determining Set}\label{sec:support}

A field with given averages is given itself on the bonds of the layer $\Gamma_0$, that is, on every bond which has no end in $\Omega_1$. The variational problem is therefore a problem for the field near $\Omega_1$ only, and its hypotheses should concern the data near $\Omega_1$ only. This is how the problem is used: the test functions of a step and the spaces of Section~7 are defined through the minimisers of determining sets which are built inside a cube or a localisation domain, and outside it no field is given. Following \cite[p.~255]{BalabanCRE} we call a neighbourhood of $\Omega_1$ a support of the determining set.

\begin{definition}[Supports]\label{def:support}
Let a sequence of domains with the block size $M_1$ be given, where $M_1 \geq L^2$. A \emph{support} is a set $W \subset \mathbb T$ which is a union of cubes of the grid of side $M_1\eta$ and contains every point at an $\ell^\infty$-distance of at most $M_1\eta$ from $\Omega_1$. The bonds and the plaquettes \emph{of $W$} are those with all ends and all corners in $W$. With these in the place of the bonds and the plaquettes of $T_\eta$, and with the bonds of the layers with both ends in $W$:
\begin{enumerate}[label=(\roman*)]
\item $\mathfrak U_W(\varepsilon_0)$ is the set of the gauge fields $U$ on the bonds of $W$ which satisfy \eqref{eq:plaquette-regular} for the plaquettes of $W$ and for the bonds of $W$ all of whose plaquettes are plaquettes of $W$;
\item \emph{data on $W$}, their plaquettes, and their regularity are defined as in Definitions~\ref{def:regular-fields} and~\ref{def:regular-data}, and $\mathfrak V_W(V)$ is the set of the gauge fields $U$ on the bonds of $W$ with $M^j(U)(c) = V(c)$ for the bonds $c$ of $\Gamma_j$ with both ends in $W$;
\item $A_W(U)$ is the sum \eqref{eq:action} over the plaquettes of $W$, and $U \in \mathfrak U_W(\varepsilon_0) \cap \mathfrak V_W(V)$ is \emph{critical}, or \emph{minimal}, on $W$ if it satisfies the conditions of Definition~\ref{def:minimal} with $A_W$ and with the $1$-forms $\delta A$ on the bonds of $W$.
\end{enumerate}
\end{definition}

A bond of a layer $\Gamma_j$ with $j \geq 1$ has an end in $\Omega_j \subset \Omega_1$, its other end is within the distance $a_j$ of $\Omega_j$, and the blocks $B^j$ of its ends are within $\frac{3}{2}a_j$ of $\Omega_j$; for $j = 1$ this is less than $M_1\eta$, and for $j \geq 2$ these blocks lie in $\Omega_{j-1} \subset \Omega_1$ by \eqref{eq:separation}. So the bonds of the layers $\Gamma_j$, $j \geq 1$, are bonds of $W$, and their averages are determined by the field on $W$ (Lemma~\ref{lem:average}). A $1$-form $\delta A$ with $Q\,\delta A = 0$ vanishes on the bonds of $\Gamma_0$, so that it is supported on the bonds with an end in $\Omega_1$; the plaquettes of these bonds have their corners within the distance $2\eta$ of $\Omega_1$ and are plaquettes of $W$. The pairing $\langle\delta A, J(U)\rangle$ and the operator $Q$ on such forms are therefore defined by the field on $W$.

\begin{lemma}[Restriction to a Support]\label{lem:support}
Let $W$ be a support, let $V$ and $\varepsilon_0$ be as in Definition~\ref{def:minimal}, and let $U \in \mathfrak U(\varepsilon_0) \cap \mathfrak V(V)$. Then the restriction $U|_W$ lies in $\mathfrak U_W(\varepsilon_0) \cap \mathfrak V_W(V|_W)$, the map $U' \mapsto U'|_W$ is a bijection of $\mathfrak U(\varepsilon_0) \cap \mathfrak V(V)$ onto $\mathfrak U_W(\varepsilon_0) \cap \mathfrak V_W(V|_W)$ along which $A(U') - A_W(U'|_W)$ is constant, and $U$ is critical, or minimal, if and only if $U|_W$ is critical, or minimal, on $W$.
\end{lemma}

\begin{proof}
The conditions which define $\mathfrak U_W(\varepsilon_0)$ and $\mathfrak V_W(V|_W)$ are among those which define $\mathfrak U(\varepsilon_0)$ and $\mathfrak V(V)$, by the remark before the lemma; this gives the first statement. Every $U' \in \mathfrak V(V)$ equals $V$ on the bonds with no end in $\Omega_1$, since $M^0$ is the identity, and a bond which is not a bond of $W$ has no end in $\Omega_1$. Hence $U'$ is determined by $U'|_W$, and the restriction is injective. Conversely let $U''$ lie in $\mathfrak U_W(\varepsilon_0) \cap \mathfrak V_W(V|_W)$, and let $U'$ equal $U''$ on the bonds of $W$ and $V$ on the others, which are bonds of $\Gamma_0$. Then $U' \in \mathfrak V(V)$, because every bond of a layer is a bond of $W$ or one of these. A plaquette which is not a plaquette of $W$ has a corner outside $W$, hence no corner in $\Omega_1$ and no bond with an end in $\Omega_1$, so that on its bonds $U' = V = U$; the same holds for the plaquettes of a bond which has a plaquette that is not a plaquette of $W$. The conditions \eqref{eq:plaquette-regular} for these plaquettes and bonds therefore hold for $U'$ because they hold for $U$, and for the others they hold because $U'' \in \mathfrak U_W(\varepsilon_0)$. So $U' \in \mathfrak U(\varepsilon_0)$, and the restriction is onto. The difference $A(U') - A_W(U'|_W)$ is the sum over the plaquettes which are not plaquettes of $W$, on which $U' = V$. Minimality is thus the same condition on both sides, and so is criticality, by the remark before the lemma.
\end{proof}

\begin{theorem}[The Variational Problem on a Support]\label{thm:support}
Lemma~\ref{lem:problem-covariance}, Theorem~\ref{thm:minimiser}, Corollary~\ref{cor:minimiser-regular}, and Theorems~\ref{thm:critical} and~\ref{thm:background-field} hold, with the same constants, for a sequence of domains with a support $W$, when the fields and the data are given on $W$, the sets $\mathfrak U(\varepsilon_0)$ and $\mathfrak V(V)$ and the action are replaced by $\mathfrak U_W(\varepsilon_0)$, $\mathfrak V_W(V)$, and $A_W$, the gauge transformations are those of the sites of $W$, the operators of Section~\ref{sec:propagators} are taken with the Dirichlet domain $W$, and the cubes of Theorem~\ref{thm:minimiser}\emph{(iii)} and of Corollary~\ref{cor:minimiser-regular} are those for which the neighbourhood named there is contained in $W$.
\end{theorem}

For $W = \mathbb T$ these are the statements of the preceding subsections. By Lemma~\ref{lem:support} the minimiser on a support of data which are regular on the whole torus is the restriction of the minimiser on the torus. The point of the theorem is that nothing is assumed outside $W$; this is the generality of \cite{BalabanVAR}, whose domain $\Omega_0$ is the support.

\section{The Small-Field Step}\label{sec:small-field}

\subsection{The Spaces of Pairs}\label{sec:pair-spaces}

The terms of an effective action are functions of the background field, and the induction needs them on complex neighbourhoods of the real backgrounds, as analytic functions with bounds. The neighbourhoods have to be large enough to contain every configuration at which a later step evaluates a term, and among these are the minimisers of the averages of a configuration: a step replaces the argument $\mathbf U$ of an old term by the configuration of least action with the averages of $\mathbf U$, on each of the scales below that of the term. The definition of the spaces therefore contains, besides bounds on the pair itself, bounds on these derived configurations \cite[Definition, p.~262]{BalabanRG1}. We first define the determining sets inside a domain and the derived configurations.

\begin{definition}[The Minimal Sequences of a Domain]\label{def:minimal-sequence}
Let $M_{\mathrm c}$ be the smallest power of $L$ with $M_{\mathrm c} \geq LR_2M_1$, with the numbers $M_1$ and $R_2$ of Theorem~\ref{thm:minimiser}. Let $1 \leq j \leq k$ and let $X \subset \mathbb T$ be a union of cubes of the grid of side $M_{\mathrm c}a_j$. We define $\Omega_1(X) \supset \dots \supset \Omega_j(X)$ recursively: $\Omega_1(X)$ is the union of the cubes of the grid of side $M_{\mathrm c}a_1$ which lie in $X$ at an $\ell^\infty$-distance of at least $M_1\eta$ from the complement of $X$, and, for $1 \leq n < j$, $\Omega_{n+1}(X)$ is the union of the cubes of the grid of side $M_{\mathrm c}a_{n+1}$ which lie in $\Omega_n(X)$ at an $\ell^\infty$-distance greater than $27M_{\mathrm c}a_n$ from the complement of $\Omega_n(X)$. For $1 \leq n \leq j$ the \emph{minimal sequence of $X$ of the order $n$} is the sequence of domains
\begin{equation}\label{eq:minimal-sequence}
    \mathbb T \supset \Omega_1(X) \supset \dots \supset \Omega_n(X) \supset \emptyset \supset \dots \supset \emptyset\;,
\end{equation}
which has the block size $M_{\mathrm c}$, the separation $27$, and the support $X$; its determining set is denoted by $\mathbf B_n(X)$.
\end{definition}

These are the determining sets $\mathbf B_j(\Omega)$ of \cite[p.~256]{BalabanCRE} and the sequences of the condition (iv) of \cite[p.~262]{BalabanRG1}, with the block size and the separation for which Corollary~\ref{cor:minimiser-regular} holds. The passage from $X$ to $\Omega_1(X)$ removes a layer of the width $M_{\mathrm c}a_1$, and the passage from $\Omega_n(X)$ to $\Omega_{n+1}(X)$ a layer of a width of at most $(27 + L)M_{\mathrm c}a_n$, so that $\Omega_j(X)$ contains every point of $X$ at a distance of at least $2(27 + L)M_{\mathrm c}a_{j-1}$ from the complement of $X$.

\begin{definition}[The Minimiser of the Averages]\label{def:derived-minimiser}
Let $X$, $j$, and $1 \leq n \leq j$ be as in Definition~\ref{def:minimal-sequence}. Let $U$ be a gauge field on the bonds of $X$ whose averages on the layers of $\mathbf B_n(X)$, the data $V_0$ with $V_0(c) \coloneqq M^i(U)(c)$ for the bonds $c$ of the layer $i$, are regular with a constant $\varepsilon_1$ for which Theorem~\ref{thm:background-field} holds on the support $X$, and let $U_0$ be the minimal configuration for $V_0$ on $X$ which is in the axial gauge. Let $\mathbf U$ be a configuration on the bonds of $X$ such that
\begin{equation}\label{eq:data-deviation}
    B(c) \coloneqq \frac{1}{i}\operatorname{Log}\bigl(M^i(\mathbf U)(c)\,V_0(c)^{-1}\bigr)
\end{equation}
is defined for the bonds $c$ of the layers and $\abs{B(c)} \leq C_{\mathrm d}\varepsilon_1$. Then
\begin{equation}\label{eq:derived-minimiser}
    U_{n,X}(\mathbf U; U) \coloneqq U(B)\;,
    \qquad
    J_{n,X}(\mathbf U; U) \coloneqq J\bigl(U(B)\bigr)\;,
\end{equation}
with the configuration \eqref{eq:moved-minimiser} of Theorem~\ref{thm:background-field}\emph{(iii)} at the base $U_0$, and with the current \eqref{eq:current}.
\end{definition}

For a gauge field $\mathbf U$ the configuration $U_{n,X}(\mathbf U; U)$ is the minimal configuration with the averages of $\mathbf U$ which is in the axial gauge, by Theorem~\ref{thm:background-field}(iii), and it does not depend on $U$. For a complex $\mathbf U$ it is the value of the analytic continuation of the minimiser, as a function of the data, from the real data $V_0$.

\begin{definition}[The Space of Pairs of a Scale]\label{def:pair-space}
Let $1 \leq j \leq k$, let $M \geq M_{\mathrm c}$ be the block size of the localisation domains, let $X \in \mathbf D_j$, and let $\alpha_0, \alpha_1, \gamma_0 > 0$ and $B_{\mathrm r} \geq 1$. The set $\mathfrak U^{c}_j(X; \alpha_0, \alpha_1, \gamma_0)$ consists of the pairs $(\mathbf U^u, R(u)\mathbf J)$, where $u$ is a gauge transformation of the sites of $X$ with values in $G^c$ and $(\mathbf U, \mathbf J)$ is a pair on the bonds of $X$ of the form $\mathbf U = e^{i\eta A'}U$ which satisfies the following four conditions.
\begin{enumerate}[label=(\roman*)]
\item $U$ is a gauge field with $\abs{U(\partial p) - \id} \leq \alpha_0(\eta/a_j)^2$ for the plaquettes $p$ of $X$, and for every cube $\square \subset X$ which is a union of cubes of $\pi_j$ and has a side $s_\square \leq 20LMa_j$ there is a gauge transformation $v$ on $\square$ with values in $G$ such that $U^v = e^{i\eta A}$ on the bonds of $\square$ with
\begin{equation}\label{eq:pair-space-real}
    a_j\abs A \leq \frac{s_\square}{a_j}\,B_{\mathrm r}\alpha_0\;,
    \qquad
    a_j^2\abs{\nabla_{\id}A} \leq \frac{s_\square}{a_j}\,B_{\mathrm r}\alpha_0\;.
\end{equation}
\item $A'$ is a $1$-form with values in $\mathfrak g^c$, and $a_j\abs{A'} \leq \alpha_1$, $a_j^2\abs{\nabla_UA'} \leq \alpha_1$ on $X$.
\item $\abs{\mathbf U(\partial p) - \id} \leq \alpha_0(\eta/a_j)^2$ for the plaquettes $p$ of $X$, and $a_j^3\abs{\mathbf J(b)} \leq \gamma_0$ for the bonds $b$ of $X$.
\item For $1 \leq n \leq j$ the configurations $U_{n,X}(\mathbf U; U)$ and $U_{n,X}(U; U)$ of Definition~\ref{def:derived-minimiser} are defined, and each of them, written $\mathbf U_n$, satisfies
\begin{equation}\label{eq:pair-space-derived}
    \abs{\mathbf U_n(\partial p) - \id} \leq \alpha_0\Bigl(\frac{\eta}{a_j}\Bigr)^2\;,
    \qquad
    a_n^3\abs{J(\mathbf U_n)(b)} \leq \alpha_0\Bigl(\frac{a_n}{a_j}\Bigr)^2
\end{equation}
for the plaquettes $p$ and the bonds $b$ of $X^{\sim-2}$, where the operation ${}^\sim$ is that of $\pi_j$.
\end{enumerate}
We write $\mathfrak U^c_j(X; \alpha_0, \alpha_1) \coloneqq \mathfrak U^c_j(X; \alpha_0, \alpha_1, \alpha_0)$.
\end{definition}

This is the Definition of \cite[p.~262]{BalabanRG1} in the units of the lattice $T^{(k)}$: the lattice of spacing $\xi = L^{-j}$ there is $T_\eta$ in the units of the scale $j$, so that $\xi = \eta/a_j$, a field $A'$ and its derivative there are $a_jA'$ and $a_j^2\nabla A'$ here, and a current there is $a_j^3\mathbf J$ here; the current of the condition (iv) is taken there in the units of the scale $n$. The constant $B_{\mathrm r}$ is the constant $B$ of \cite[(1.12)]{BalabanRG1}, which is fixed with the constants of the construction. By definition the set is invariant under the gauge transformations with values in $G^c$, acting on pairs by \cite[(1.10)]{BalabanRG1}, and it increases with each of $\alpha_0$, $\alpha_1$, $\gamma_0$ when the conditions on the regularity of the data in Definition~\ref{def:derived-minimiser} are kept fixed.

The condition (iv) is not explicit, and the spaces are used through simpler sets which they contain \cite[p.~263]{BalabanRG1}.

\begin{proposition}[A Criterion for Membership]\label{prop:membership}
There are numbers $c_F \in (0, 1]$ and $a_6 > 0$, which depend on $L$, $G$, $M$, and $B_{\mathrm r}$ only, with the following property. Let $\alpha_0, \alpha_1 \leq a_6$, $X \in \mathbf D_j$, and let $(\mathbf U, \mathbf J)$ be a pair on the bonds of $X$ which satisfies the conditions \emph{(i)} and \emph{(iii)} of Definition~\ref{def:pair-space} with $c_F\alpha_0$ in the place of $\alpha_0$, the condition \emph{(ii)}, and $a_j^3\abs{\mathbf J} \leq \gamma_0$. Then $(\mathbf U, \mathbf J) \in \mathfrak U^c_j(X; \alpha_0, \alpha_1, \gamma_0)$. In particular, if $U$ is a gauge field on $X$ which satisfies \emph{(i)} with $c_F\alpha_0$ and $a_j^3\abs{J(U)} \leq \alpha_0$, then $(U, J(U)) \in \mathfrak U^c_j(X; \alpha_0, \alpha_1)$.
\end{proposition}

In \cite[(1.17), p.~263]{BalabanRG1} this is derived in one line from the regularity of the minimisers, applied to the averages of the complex configuration, and it is said to be obvious that the condition (iv) follows. The statement is used with the full radius $\alpha_1$ of the complex part, and only the plaquettes are required to be smaller.

A term which is localised near a region of large fields is a function of a background whose determining set has several layers in the domain of the term, and its radii of analyticity depend on the layer. For such terms the spaces are defined with respect to the sequence of domains of the background \cite[(2.34)--(2.39), p.~261]{BalabanCRE}. The derived configurations are now the minimisers of determining sets which follow the given sequence up to a scale and are minimal near the boundary of the domain.

\begin{definition}[The Localised Sequences]\label{def:localised-sequence}
Let a sequence of domains $\{ \Omega_n \}$ with a block size $M \geq M_{\mathrm c}$ and a separation $R \geq 27$ be given, let $1 \leq p \leq j \leq k$, and let $X \in \mathbf D_j$. The \emph{sequence localised to $X$, of the order $p$}, is
\begin{equation}\label{eq:localised-sequence}
    \mathbb T \supset \Omega_1 \cap \Omega_1(X) \supset \dots \supset \Omega_p \cap \Omega_p(X) \supset \emptyset \supset \dots \supset \emptyset\;,
\end{equation}
with the domains $\Omega_n(X)$ of Definition~\ref{def:minimal-sequence}. It is a sequence of domains with the block size $M_{\mathrm c}$, the separation $27$, and the support $X$, and its determining set is denoted by $\mathbf B_p(X; \{ \Omega_n \})$. For a gauge field $U$ and a configuration $\mathbf U$ on the bonds of $X$, the configuration $U_{p,X}(\mathbf U; U)$ and its current are defined as in Definition~\ref{def:derived-minimiser}, with this determining set in the place of $\mathbf B_n(X)$.
\end{definition}

The separation of \eqref{eq:localised-sequence} is that of the two sequences which are intersected: a point of $\Omega_{n+1} \cap \Omega_{n+1}(X)$ is at a distance greater than $27M_{\mathrm c}a_n$ from the complement of $\Omega_n$ and from that of $\Omega_n(X)$. The determining set is $\mathbf B_p(X) \cup \{ \Gamma_i \}_{i<p}$ in the notation of \cite[p.~261]{BalabanCRE}: in the interior of $X$ a point of the layer $i < p$ of the given sequence remains in the layer $i$, the points of the layers $i \geq p$ are in the layer $p$, and near the boundary of $X$ the layers are those of the minimal sequence.

\begin{definition}[The Space of Pairs of Several Scales]\label{def:pair-space-layers}
Let the sequence of domains, $j$, and $X$ be as in Definition~\ref{def:localised-sequence}, let $\tilde\alpha_0 = (\alpha_{0,n})_{0 \leq n \leq j}$ and $\tilde\alpha_1 = (\alpha_{1,n})_{0 \leq n \leq j}$ be sequences of positive numbers, let $0 < \beta < 1$ and $B_{\mathrm r} \geq 1$, and put
\begin{equation}\label{eq:layer-factor}
    c_{j,n} \coloneqq 1 - \beta\bigl(1 - 2^{-(j-n)}\bigr)\;,
    \qquad
    X_n \coloneqq X \cap (\Omega_n \setminus \Omega_{n+1}) \quad (n < j)\;,
    \qquad
    X_j \coloneqq X \cap \Omega_j\;.
\end{equation}
The set $\tilde{\mathfrak U}^c_j(X; \tilde\alpha_0, \tilde\alpha_1)$ consists of the pairs $(\mathbf U, \mathbf J)$ on the bonds of $X$ of the form $\mathbf U = e^{i\eta A'}U$, with a gauge field $U$ and a $1$-form $A'$ with values in $\mathfrak g^c$, which satisfy the following conditions for $0 \leq n \leq j$.
\begin{enumerate}[label=(\roman*)]
\item On the plaquettes and the bonds of $X_n$,
\begin{equation}\label{eq:pair-space-layers}
    \abs{U(\partial p) - \id},\ \abs{\mathbf U(\partial p) - \id} \leq c_{j,n}\alpha_{0,n}\Bigl(\frac{\eta}{a_n}\Bigr)^2\;,
    \qquad
    a_n^3\abs{\mathbf J(b)} \leq c_{j,n}\alpha_{0,n}\;.
\end{equation}
\item For $1 \leq p \leq j$ the configurations $U_{p,X}(\mathbf U; U)$ and $U_{p,X}(U; U)$ of Definition~\ref{def:localised-sequence} are defined, and each of them, written $\mathbf U_p$, satisfies on the plaquettes and the bonds of $X_n$
\begin{equation}\label{eq:pair-space-layers-derived}
    \abs{\mathbf U_p(\partial p) - \id} \leq c_{j,n}\alpha_{0,n}\Bigl(\frac{\eta}{a_n}\Bigr)^2\;,
    \qquad
    a_p^3\abs{J(\mathbf U_p)(b)} \leq c_{j,n}\alpha_{0,n}\Bigl(\frac{a_p}{a_n}\Bigr)^{e}\;,
\end{equation}
with $e = 3$ for $n \leq p$ and $e = 2$ for $n > p$.
\item For every cube $\square$ which, for $n < j$, is a cube of the class of the sequence of the index $n$ (Definition~\ref{def:regular-pairs}(i)), and, for $n = j$, is a union of cubes of $\pi_j$ contained in $\Omega_j$ with a side $10Ma_j \leq s_\square \leq 20Ma_j$, there is a gauge transformation $v$ on $\square \cap X$ with values in $G$ such that $U^v = e^{i\eta A}$ on the bonds of $\square \cap X$ with $a_n\abs A \leq (s_\square/a_n)B_{\mathrm r}\alpha_{0,n}$ and $a_n^2\abs{\nabla_{\id}A} \leq (s_\square/a_n)B_{\mathrm r}\alpha_{0,n}$.
\item $a_n\abs{A'} \leq c_{j,n}\alpha_{1,n}$ and $a_n^2\abs{\nabla_UA'} \leq c_{j,n}\alpha_{1,n}$ on $X_n$.
\end{enumerate}
\end{definition}

These are the conditions \cite[(2.34)--(2.39), p.~261]{BalabanCRE}; the exponents in \eqref{eq:pair-space-layers-derived} are those of \cite[(2.37)]{BalabanCRE}, where the current of $\mathbf U_p$ is taken in the units of the scale $p$, and for $n > p$ they agree with \eqref{eq:pair-space-derived}. The conditions are stated there for the layers $n \geq 1$. We include $n = 0$, the part of $X$ outside $\Omega_1$, because the derived configurations of (ii) are defined only if the data of $U$ are regular on the whole support $X$. The set is invariant under the gauge transformations with values in $G$, and not under those with values in $G^c$. The factor $c_{j,n}$ decreases from $1$ at $n = j$ to more than $1 - \beta$: a term is analytic in a domain which shrinks slightly on the layers below its own scale, and this is the room in which the operations of a step move the configurations at which they evaluate the older terms.

\subsection{The Terms of an Effective Action}\label{sec:terms}

An effective action is a sum of local terms, each a function of the background in a localisation domain, and what the induction carries from step to step is a short list of properties of these functions. We state the list for the terms which a step creates in the region of small fields \cite[(1.7)--(1.19), pp.~261--263]{BalabanRG1}, in the form with a base point which \cite[(2.26), (2.27), p.~259]{BalabanCRE} introduces. Since the sets of Definition~\ref{def:pair-space} are closed, we call a function on such a set \emph{holomorphic} if it is the restriction of a holomorphic function on an open subset of the complex manifold of the pairs on the bonds of $X$, that is, of $(G^c)^{\mathsf b} \times (\mathfrak g^c)^{\mathsf b}$ with $\mathsf b$ the set of the positive bonds of $X$, which contains the set.

\begin{definition}[Regular Terms]\label{def:regular-terms}
Let $1 \leq j \leq k$, and let $\alpha_0$, $\alpha_1$, $E_0$, and $\kappa$ be positive numbers. A \emph{family of regular terms of the scale $j$} with these constants is a family of functions $\mathbf E(X, \cdot\,, z)$, one for each $X \in \mathbf D_j$ and each point $z \in T^{(j)} \cap X$, with the following properties.
\begin{enumerate}[label=(\roman*)]
\item \emph{Locality and analyticity.} $\mathbf E(X, \cdot\,, z)$ is a holomorphic function on $\mathfrak U^c_j(X; \alpha_0, \alpha_1)$, a set of pairs on the bonds of $X$.
\item \emph{Gauge invariance.} $\mathbf E(X, (\mathbf U^u, R(u)\mathbf J), z) = \mathbf E(X, (\mathbf U, \mathbf J), z)$ for all gauge transformations $u$ of the sites of $X$ with values in $G^c$.
\item \emph{The bound.} $\abs{\mathbf E(X, (\mathbf U, \mathbf J), z)} \leq E_0e^{-\kappa d_j(X)}$ on $\mathfrak U^c_j(X; \alpha_0, \alpha_1)$.
\item \emph{Covariance.} Let $r$ be an isometry of the torus which maps $T^{(j)}$ onto itself, let $U$ be a gauge field on $T_\eta$, and let $rU$ be the gauge field with $(rU)(rb) \coloneqq U(b)$ for the oriented bonds $b$ of $T_\eta$. If, for $U$ and for $rU$ in the place of $U$, the pair $(U|_X, J(U)|_X)$ lies in $\mathfrak U^c_j(X; \alpha_0, \alpha_1)$ for every $X \in \mathbf D_j$, then
\begin{equation}\label{eq:regular-covariance}
    \mathbf E(rU, rz) = \mathbf E(U, z)\;,
    \qquad
    \mathbf E(U, z) \coloneqq \sum_{X \in \mathbf D_j,\ X \ni z}\mathbf E\bigl(X, (U|_X, J(U)|_X), z\bigr)\;.
\end{equation}
\end{enumerate}
For a gauge field $U$, a domain $\Lambda \in \mathbf D_j$, and a point $z \in T^{(j)} \cap \Lambda$ we write
\begin{equation}\label{eq:regular-sum}
    \mathbf E(\Lambda, U, z) \coloneqq \sum_{X \in \mathbf D_j,\ z \in X \subset \Lambda}\mathbf E\bigl(X, (U|_X, J(U)|_X), z\bigr)\;.
\end{equation}
\end{definition}

The current in \eqref{eq:regular-covariance} and \eqref{eq:regular-sum} is the current of $U$ on the torus, restricted to the bonds of $X$; at a bond near the boundary of $X$ it depends on $U$ outside $X$, and the pair is a point of the space, at which the term is evaluated. The sums are finite, and by (iii) and Lemma~\ref{lem:cube-count} they are bounded uniformly in the torus for $\kappa$ large. The isometries of (iv) are generated by the translations by the vectors of $a_j\mathbb Z^4$, the permutations of the coordinates, and the reflections; they map $T_\eta$ onto itself, but not the grid $\pi_j$, and for this reason the covariance is a property of the sum over all domains and not of its terms \cite[(2.29), p.~260]{BalabanCRE}.

\subsection{The Fluctuation Field of a Step}\label{sec:fluctuation-field}

A step integrates the field of a lattice $T^{(n)}$ with its averages on $T^{(n+1)}$ held fixed, in a region in which the field is regular. The integration variable is the deviation of the field from the configuration of least action with the given averages, and three things have to be known about it: coordinates in which the constraint is linear, with the Jacobian of the passage to them; the quadratic form of the action in these coordinates, which defines the Gaussian measure of the step; and the remainder of the action \cite[Sect.~2]{BalabanRG1}, \cite[Sect.~E]{BalabanPBG}. Throughout this subsection the mean transports are used, $\mathfrak t = \mathrm m$.

\begin{definition}[The Setting of a Step]\label{def:step-setting}
Let $0 \leq n < k$, and let a sequence of domains $\mathbb T \supset \Omega_1 \supset \dots \supset \Omega_{n+1}$, with $\Omega_j = \emptyset$ for $j > n + 1$, with a block size $M \geq M_{\mathrm c}$ and a separation $R \geq 27$ be given. We write $\mathcal O \coloneqq \Omega_{n+1}$ and call the given sequence the \emph{new} one. The \emph{old} sequence is $\mathbb T \supset \Omega_1 \supset \dots \supset \Omega_n$; its layers are those of the new sequence for $j < n$, and its layer of the scale $n$ is $\Omega_n^{(n)}$.
\begin{enumerate}[label=(\roman*)]
\item $\mathsf b(\mathcal O)$ is the set of the positive bonds of $T^{(n)}$ with at least one end in $\mathcal O$, and $\mathsf c(\mathcal O)$ is the set of the positive bonds of $T^{(n+1)}$ with at least one end in $\mathcal O$, the bonds of the layer $\Gamma_{n+1}$ of the new sequence.
\item For $c \in \mathsf c(\mathcal O)$ the \emph{central bond} $b_0(c)$ is the bond of the path $[c]$ of Definition~\ref{def:block-paths} which goes from $B(c_-)$ to $B(c_+)$, and $\mathsf b^*(\mathcal O) \coloneqq \mathsf b(\mathcal O) \setminus \{ b_0(c) \,|\, c \in \mathsf c(\mathcal O) \}$.
\item Let $V$ be data on the determining set of the new sequence which are regular with a constant for which Theorem~\ref{thm:background-field} holds, let $U_{n+1}$ be the minimal configuration for $V$ in the axial gauge, and
\begin{equation}\label{eq:evaluation-field}
    V^{(n)} \coloneqq M^n(U_{n+1})\;,
\end{equation}
a gauge field on the bonds of $T^{(n)}$ with both ends in $\Omega_n$.
\end{enumerate}
\end{definition}

The old sequence has the data which equal $V$ on the layers $j < n$ and are a gauge field $V_n$ on the bonds of $T^{(n)}$ in $\Omega_n$; outside $\mathsf b(\mathcal O)$ this field is given by $V$, and on $\mathsf b(\mathcal O)$ it is the variable of the step. The configuration $U_{n+1}$ has the averages $V^{(n)}$ on the scale $n$ and is critical for the old sequence, since a $1$-form with vanishing averages of the order $n$ has vanishing averages of the order $n + 1$; so it is the minimal configuration of the old sequence for the data with $V_n = V^{(n)}$ (Theorem~\ref{thm:minimiser}(ii)), and $V^{(n)}$ is the field about which the step expands \cite[(2.3), p.~265]{BalabanRG1}, \cite[(3.10), p.~266]{BalabanCRE}. Since $U_{n+1}$ is in the axial gauge, $V^{(n)}(y, x) = \id$ for $y \in \Gamma_{n+1}$ and $x \in B(y)$.

\begin{definition}[The Average on the Unit Lattice]\label{def:unit-average}
In the setting of Definition~\ref{def:step-setting}, let $B$ be a function on $\mathsf b(\mathcal O)$ with values in $\mathfrak g^c$, extended by zero to the other bonds of $T^{(n)}$.
\begin{enumerate}[label=(\roman*)]
\item The linear map $\tilde Q$ and the function $\tilde C$, with values in the functions on $\mathsf c(\mathcal O)$, are defined by
\begin{equation}\label{eq:unit-average}
    \frac{1}{i}\operatorname{Log}\Bigl(M\bigl(e^{iB}V^{(n)}\bigr)(c)\,M\bigl(V^{(n)}\bigr)(c)^{-1}\Bigr) = L(\tilde QB)(c) + \tilde C(B)(c)\;,
\end{equation}
with $\tilde Q$ linear and $\tilde C$ vanishing to the second order; that is, $\tilde Q = Q^{\mathrm{pl}}(V^{(n)})$ in the notation of Definition~\ref{def:relative-average}, applied to functions on bonds without a change of their dimension. The part of the second order of $\tilde C$ is denoted by $\tilde C^{(2)}$.
\item For $c \in \mathsf c(\mathcal O)$ let $q_c$ be the linear map of $\mathfrak g^c$ with $(\tilde QB)(c) = q_cB(b_0(c))$ for the functions $B$ which vanish on the bonds other than $b_0(c)$. For a function $\mu$ on $\mathsf c(\mathcal O)$ let $h\mu$ be the function on $\mathsf b(\mathcal O)$ with $(h\mu)(b_0(c)) \coloneqq (Lq_c)^{-1}\mu(c)$ which vanishes on $\mathsf b^*(\mathcal O)$, so that $L\tilde Qh = I$.
\item For a function $\tilde B$ on $\mathsf b^*(\mathcal O)$, extended by zero to the central bonds, $C\tilde B \coloneqq \tilde B - hL\tilde Q\tilde B$, so that $\tilde QC\tilde B = 0$ and $C\tilde B = \tilde B$ on $\mathsf b^*(\mathcal O)$.
\end{enumerate}
\end{definition}

The central bond of $c$ lies on the path $[c]$ and on no other path $[x, x']$ of a bond of $T^{(n+1)}$, and it lies in no block, so that it enters the average of $c$ alone, and it enters it through the straight transport and the loops of $c$. At the trivial configuration $q_c = L^{-5}$.

\begin{proposition}[The Linearising Transformation on the Unit Lattice]\label{prop:unit-linearising}
There are numbers $a_7 > 0$ and $C_7$, which depend on $L$ and $G$ only, such that in the setting of Definition~\ref{def:step-setting}:
\begin{enumerate}[label=(\roman*)]
\item the maps $q_c$ are invertible, $\norm{L^5q_c - I} \leq \frac{1}{2}$, and $q_c$ maps $\mathfrak g$ to itself;
\item for every function $B$ on $\mathsf b(\mathcal O)$ with values in $\mathfrak g^c$ and $\abs B \leq a_7$ there is exactly one function $\tilde D(B)$ on $\mathsf c(\mathcal O)$ with $\abs{\tilde D(B)} \leq a_7$ and
\begin{equation}\label{eq:unit-linearising}
    \tilde D(B) = \tilde C\bigl(B - h\tilde D(B)\bigr)\;;
\end{equation}
it is holomorphic in $B$, real for real $B$, $\tilde D(B)(c)$ depends only on $B$ on the bonds with both ends in $B(c_-) \cup B(c_+)$, $\abs{\tilde D(B)(c)} \leq C_7\abs B^2$, and its part of the second order is $\tilde C^{(2)}(B)$;
\item the map $B \mapsto B' \coloneqq B - h\tilde D(B)$ is injective on the set of the $B$ with $\abs B \leq a_7$, its image contains the set of the $B'$ with $\abs{B'} \leq \frac{1}{2}a_7$, and $L\tilde QB' + \tilde C(B') = L\tilde QB$.
\end{enumerate}
\end{proposition}

This is the construction of \cite[p.~267]{BalabanRG1}. The next statement gives the coordinates of the step and their Jacobian. In the papers the passage to them is a computation with $\delta$-functions; here it is a statement on the image of a measure under the average.

\begin{definition}[The Density of the Haar Measure]\label{def:haar-density}
Let $\dd V$ be the Haar measure of $G$ of total mass one and $\dd X$ the Lebesgue measure of $\mathfrak g$ which belongs to the scalar product of Definition~\ref{def:group}. The function $\sigma$ on the domain of the chart in $\mathfrak g$ is defined by $\int f(V)\,\dd V = \int f(e^{iX})\sigma(X)\,\dd X$ for the continuous functions $f$ on $G$ which are supported in the image of the chart, and $\sigma_0 \coloneqq \sigma(0)$. On the functions on a finite set of bonds with values in $G$ or in $\mathfrak g$ we use the product measures.
\end{definition}

\begin{theorem}[The Coordinates of the Fluctuation]\label{thm:fluctuation-coordinates}
There is a number $a_8 \in (0, \frac{1}{2}a_7]$, which depends on $L$ and $G$ only, with the following property. In the setting of Definition~\ref{def:step-setting}, let the data on the layers $j \leq n$ be fixed, let $\mathcal W$ be a set of gauge fields $W$ on $\mathsf c(\mathcal O)$ such that the data which equal $W$ on the layer $n + 1$ are regular as in Definition~\ref{def:step-setting}(iii), and write $V^{(n)}_W$, $h_W$, $C_W$, $\tilde D_W$ for the objects of Definitions~\ref{def:step-setting} and~\ref{def:unit-average} at these data. For $W \in \mathcal W$ and a function $\tilde B$ on $\mathsf b^*(\mathcal O)$ with values in $\mathfrak g$ and $\abs{\tilde B} \leq a_8$ put
\begin{equation}\label{eq:fluctuation-coordinates}
    B'_W(\tilde B) \coloneqq C_W\tilde B - h_W\tilde D_W(C_W\tilde B)\;,
    \qquad
    \Psi(W, \tilde B) \coloneqq e^{iB'_W(\tilde B)}\,V^{(n)}_W\;,
\end{equation}
a gauge field on $\mathsf b(\mathcal O)$. Then:
\begin{enumerate}[label=(\roman*)]
\item $M(\Psi(W, \tilde B))(c) = W(c)$ for $c \in \mathsf c(\mathcal O)$, the average being taken with the fixed field on the bonds outside $\mathsf b(\mathcal O)$;
\item $\Psi$ is injective, and its image contains every gauge field $V_n$ on $\mathsf b(\mathcal O)$ such that the averages $W$ of $V_n$ on $\mathsf c(\mathcal O)$ lie in $\mathcal W$ and $\abs{\frac{1}{i}\log(V_n(b)V^{(n)}_W(b)^{-1})} \leq \frac{1}{4}a_8$ for all $b \in \mathsf b(\mathcal O)$;
\item for every bounded Borel function $F$ of the gauge fields on $\mathsf b(\mathcal O)$ which vanishes outside the image of $\Psi$, and every bounded Borel function $f$ of the gauge fields on $\mathsf c(\mathcal O)$,
\begin{equation}\label{eq:fluctuation-jacobian}
    \int f\bigl(M(V_n)\bigr)F(V_n)\,\dd V_n = \int_{\mathcal W}f(W)\int F\bigl(\Psi(W, \tilde B)\bigr)\,\mathcal J(W, \tilde B)\,\dd\tilde B\,\dd W\;,
\end{equation}
where the inner integral is over the $\tilde B$ with $\abs{\tilde B} \leq a_8$ and
\begin{equation}\label{eq:jacobian}
    \mathcal J(W, \tilde B) \coloneqq \prod_{c \in \mathsf c(\mathcal O)}\frac{1}{\sigma_0\abs{\det(Lq_c)}}\;\prod_{b \in \mathsf b(\mathcal O)}\sigma\bigl(B'_W(\tilde B)(b)\bigr)\;\det\Bigl(I - h_W\frac{\partial\tilde D_W}{\partial B}(C_W\tilde B)\Bigr)\;,
\end{equation}
the determinants being those of real linear maps of $\mathfrak g$ and of the functions on $\mathsf b(\mathcal O)$ with values in $\mathfrak g$.
\end{enumerate}
\end{theorem}

In the notation of the papers \eqref{eq:fluctuation-jacobian} reads $\int\dd V_n\,\delta(M(V_n)W^{-1})F(V_n) = \int\dd\tilde B\,\mathcal J F$, and the three factors of \eqref{eq:jacobian} are the factor $z^{(n)}$ of the elimination of the central bonds, the density of the Haar measure, and the Jacobian of the linearising transformation \cite[(2.10)--(2.12), pp.~267--268]{BalabanRG1}, \cite[(1.14), (1.15), p.~249]{BalabanCRE}. The scaling $\tilde B = g_nA$ of the papers multiplies the right side by $g_n^{\dim\mathfrak g\,\#\mathsf b^*(\mathcal O)}$.

The gauge is fixed in the blocks by a weight and not by a condition, so that the symmetry of the lattice is kept \cite[(1.5)--(1.7), p.~247]{BalabanCRE}.

\begin{definition}[The Gauge-Fixing Function]\label{def:gauge-fixing}
For $g > 0$ and $0 < \epsilon \leq r_G$ let
\begin{equation}\label{eq:gauge-normalisation}
    z(g, \epsilon) \coloneqq \int_G\chi\bigl(\{ \abs{v - \id} < \epsilon \}\bigr)\exp\Bigl(-\frac{1}{g^2}\bigl(1 - \re\tr v\bigr)\Bigr)\,\dd v\;.
\end{equation}
For a set $\mathsf P \subset T^{(n+1)}$ and a gauge field $V_n$ on $T^{(n)}$ whose transports in the blocks of the points of $\mathsf P$ are defined,
\begin{equation}\label{eq:gauge-fixing-function}
\begin{aligned}
    \mathbf G(\mathsf P, V_n) &\coloneqq \sum_{y \in \mathsf P}\sum_{x \in B(y),\, x \neq y}\bigl(1 - \re\tr V_n(y, x)\bigr)\;,\\
    \chi_{\mathrm{Ax}}(\mathsf P, V_n) &\coloneqq \prod_{y \in \mathsf P}\prod_{x \in B(y),\, x \neq y}\chi\bigl(\{ \abs{V_n(y, x) - \id} < \epsilon \}\bigr)\;.
\end{aligned}
\end{equation}
\end{definition}

\begin{lemma}[The Insertion of the Gauge Fixing]\label{lem:gauge-fixing}
Let $\mathsf P$ and $V_n$ be as in Definition~\ref{def:gauge-fixing}, and suppose that the transports $V^u_n(y, x)$ are defined for all gauge transformations $u$ of $T^{(n)}$ which equal $\id$ outside the set $\mathsf P' \coloneqq \bigcup_{y \in \mathsf P}B(y) \setminus \{ y \}$. Then
\begin{equation}\label{eq:gauge-fixing-identity}
    \int\chi_{\mathrm{Ax}}(\mathsf P, V^u_n)\exp\Bigl(-\frac{1}{g^2}\mathbf G(\mathsf P, V^u_n)\Bigr)\prod_{x \in \mathsf P'}\dd u(x) = z(g, \epsilon)^{(L^4-1)\#\mathsf P}\;.
\end{equation}
\end{lemma}

\begin{proof}
By the covariance of the transports, $V^u_n(y, x) = u(y)V_n(y, x)u(x)^{-1} = V_n(y, x)u(x)^{-1}$ for $y \in \mathsf P$ and $x \in B(y) \setminus \{ y \}$, since $u(y) = \id$. The integrand is therefore the product over the points $x \in \mathsf P'$ of a function of $V_n(y, x)u(x)^{-1}$, and each point lies in exactly one block. The Haar measure is invariant under $u(x) \mapsto (V_n(y, x)u(x)^{-1})$, so that each factor integrates to \eqref{eq:gauge-normalisation}, and there are $(L^4 - 1)\#\mathsf P$ factors.
\end{proof}

The hypothesis on the transports holds when the plaquettes of $V_n$ in the blocks are small, by Lemma~\ref{lem:group-mean}, because the means are taken over holonomies which differ by plaquettes, whatever the gauge. The quadratic form of the step consists of three parts: the second variation of the action along the minimisers, the term of the second order which the linearising transformation takes from the term of the first order, and the gauge-fixing function \cite[(1.5), p.~261; (2.11), p.~267]{BalabanRG1}.

\begin{definition}[The Form and the Covariance of the Unit Lattice]\label{def:unit-form}
In the setting of Definition~\ref{def:step-setting}, let $H_1$, $G_1$, and the forms of Definition~\ref{def:landau-operators} be those of the old sequence at the pair $(U_{n+1}, J(U_{n+1}))$, or at a pair $(\mathbf U, \mathbf J)$ for which they exist, and let a function $B$ on $\mathsf b(\mathcal O)$ be regarded as the function on the bonds of the layers of the old sequence which equals $a_n^{-1}B$ on $\mathsf b(\mathcal O)$ and zero elsewhere.
\begin{enumerate}[label=(\roman*)]
\item $\mathbf G^{(2)}(B)$ is the quadratic form with $\mathbf G(\Gamma_{n+1}, e^{iB}V^{(n)}) = \frac{1}{2}\mathbf G^{(2)}(B) + O(B^3)$.
\item The \emph{form of the unit lattice} is
\begin{equation}\label{eq:unit-form}
    \langle B, \Delta^{(n)}B\rangle \coloneqq \bigl\langle H_1B, \bigl(\Delta(\mathbf U) - \Delta'_\pi - \Delta^{(2)}_\pi\bigr)H_1B\bigr\rangle - 2\bigl\langle H_1h\tilde C^{(2)}(B), \mathbf J\bigr\rangle + \mathbf G^{(2)}(B)\;.
\end{equation}
\item For a set $\Lambda \subset \mathcal O$ which is a union of cubes of the grid of side $Ma_{n+1}$, let $\mathsf b^*(\Lambda)$ be the set of the bonds of $\mathsf b^*(\mathcal O)$ with both ends in $\Lambda$. The \emph{covariance of the unit lattice} in $\Lambda$ is the inverse $C^{(n)}(\Lambda)$ of the form $\tilde B \mapsto \langle C\tilde B, \Delta^{(n)}C\tilde B\rangle$ on the functions on $\mathsf b^*(\Lambda)$ with values in $\mathfrak g$, whenever this form is positive definite, and $\mu_{C^{(n)}(\Lambda)}$ is the Gaussian probability measure on these functions with this covariance.
\end{enumerate}
\end{definition}

By Theorem~\ref{thm:landau-operators}(iv) the first term of \eqref{eq:unit-form} equals $\langle B, ((QG_1Q^*)^{-1} - \mathsf a)B\rangle$ at $\mathbf J = J(\mathbf U)$, which is the form \cite[(3.156), p.~428]{BalabanPBG}.

\begin{theorem}[The Covariance of the Unit Lattice]\label{thm:unit-lattice}
There are numbers $\gamma_0 > 0$, $\delta_3 > 0$, $a_9 > 0$, and $B_7$, which depend on $L$ and $G$ only, with the following properties. In the setting of Definition~\ref{def:step-setting}, let the data be regular with a constant $\varepsilon_1$ such that $M\varepsilon_1 \leq a_9$, and let $\Lambda$ be as in Definition~\ref{def:unit-form}(iii).
\begin{enumerate}[label=(\roman*)]
\item \emph{Positivity.} $\langle C\tilde B, \Delta^{(n)}C\tilde B\rangle \geq \gamma_0\sum_{b \in \mathsf b^*(\Lambda)}\tr\tilde B(b)^2$ for the functions $\tilde B$ on $\mathsf b^*(\Lambda)$ with values in $\mathfrak g$.
\item \emph{Decay.} The kernels of $C^*\Delta^{(n)}C$ and of $C^{(n)}(\Lambda)$ satisfy $\abs{\mathcal K(b, b')} \leq B_7e^{-\delta_3\abs{b_- - b'_-}/a_n}$.
\item \emph{Analyticity.} The form $\Delta^{(n)}$ and, for $\Lambda$ fixed, the covariance are holomorphic functions of the pair, as in Theorem~\ref{thm:landau-operators}(i), on the regular pairs with constants $M\alpha_0 + \alpha_1 + \gamma \leq a_9$, and \emph{(ii)} holds there.
\item \emph{Covariance.} Under a gauge transformation $u$ of $T_\eta$ with values in $G$ the kernels are conjugated by the values of $u$ at the sites of $T^{(n)}$, and under an isometry of the torus which preserves $T^{(n+1)}$ and the domains they are transported.
\item \emph{Walks.} $C^{(n)}(\Lambda)$ has an expansion in walks with the properties of Theorem~\ref{thm:walk-structure}.
\end{enumerate}
\end{theorem}

This is \cite[Theorem~3.15, p.~432]{BalabanPBG} for the form with the gauge-fixing weight, with the lower bound which is asserted on \cite[p.~428]{BalabanPBG} and proved at the trivial configuration in \cite[Lemma~2.4]{BalabanPRT2}.

The last ingredient is the action of the old background as a function of the fluctuation. It is the expansion of Lemma~\ref{lem:functional} at the base $U_{n+1}$, composed with the function of the data.

\begin{proposition}[The Action in the Fluctuation]\label{prop:action-fluctuation}
In the setting of Definition~\ref{def:step-setting}, let $B'$ be a function on $\mathsf b(\mathcal O)$ with values in $\mathfrak g$ and $\abs{B'} \leq a_8$, let $V_n \coloneqq e^{iB'}V^{(n)}$ on $\mathsf b(\mathcal O)$ and $V_n \coloneqq V^{(n)}$ on the other bonds of $T^{(n)}$ in $\Omega_n$, and let $U_n(V_n)$ be a minimal configuration of the old sequence for the data with this field on the scale $n$. Let $\mathfrak b$ be the constraint field \eqref{eq:constraint-field} of these data relative to those with $V^{(n)}$, and let the functions of Section~\ref{sec:critical-point} be those of the old sequence at the pair $(U_{n+1}, J(U_{n+1}))$. Then
\begin{equation}\label{eq:action-fluctuation}
    A\bigl(U_n(V_n)\bigr) = A(U_{n+1}) + \langle H_1\mathfrak b, J(U_{n+1})\rangle + \frac{1}{2}\bigl\langle H_1\mathfrak b, \bigl(\Delta(U_{n+1}) - \Delta'_\pi - \Delta^{(2)}_\pi\bigr)H_1\mathfrak b\bigr\rangle + \mathsf W(\mathfrak b)\;,
\end{equation}
where
\begin{equation}\label{eq:fluctuation-potential}
    \mathsf W(\mathfrak b) \coloneqq \frac{1}{2}\bigl\langle\mathcal A_1, \bigl(\Delta(U_{n+1}) - \Delta'_\pi - \Delta^{(2)}_\pi\bigr)\mathcal A_1\bigr\rangle + \mathsf V_\pi\bigl(\mathcal A_1 + H_1\mathfrak b\bigr)\;,
    \qquad
    \mathcal A_1 \coloneqq \mathcal A_1(H_1\mathfrak b)\;,
\end{equation}
is a holomorphic function of $\mathfrak b$ which vanishes to the third order at $\mathfrak b = 0$. If $B' = B - h\tilde D(B)$ with $\tilde QB = 0$, then $\langle H_1\mathfrak b, J(U_{n+1})\rangle$ vanishes to the second order in $B$.
\end{proposition}

This is \cite[(2.6)--(2.8), p.~266]{BalabanRG1}, where $\mathsf W$ is written $V(H_1B')$, and the last statement is the vanishing of the term of the first order after the linearisation \cite[p.~267]{BalabanRG1}. On the bonds with both ends in $\mathcal O$ the constraint field is $a_n^{-1}B'$; on the bonds of $\mathsf b(\mathcal O)$ with one end outside $\mathcal O$ it is the same, since both ends of such a bond lie in the layer of the scale $n$ of the old sequence.

\subsection{The Old Terms at the New Background}\label{sec:old-terms}

After the change of variables the terms of the old effective action are evaluated at the old background, which is the new background multiplied by the exponential of a field $\eta\mathbf A$ built from the fluctuation, and each term has to be expanded in this field. A term is a function of a pair, and the second entry of the pair at the old background is its current, which contains the second differences $D^*D\mathbf A$. For the fields which occur, $D^*D\mathbf A$ is not bounded, but the covariant Laplacian is; the two differ by a term which vanishes in the Landau gauge and by local terms, and the term is evaluated at the pair in which the Laplacian stands in the place of $D^*D$ \cite[(3.10)--(3.14), p.~272]{BalabanRG1}.

\begin{lemma}[The Current of a Product]\label{lem:current-expansion}
There are numbers $c_6 > 0$ and $C_8$, which depend on $G$ only, with the following property. Let $\mathbf U$ be a configuration and $\mathbf A$ a $1$-form with values in $\mathfrak g^c$ on $T_\eta$, and let $b$ be a bond such that $\abs{\mathbf U(\partial p) - \id} \leq c_6$ and $\eta\abs{\mathbf A} \leq c_6$ on the plaquettes which contain $b$ and on their bonds. Then
\begin{equation}\label{eq:current-expansion}
    J\bigl(e^{i\eta\mathbf A}\mathbf U\bigr)(b) = J(\mathbf U)(b) + \bigl(D^*_{\mathbf U}D_{\mathbf U}\mathbf A\bigr)(b) + \mathbf F(\mathbf U, \mathbf A)(b)\;,
\end{equation}
where $\mathbf F(\mathbf U, \mathbf A)(b)$ is a holomorphic function of the values of $\mathbf U$ and of $\eta\mathbf A$ on the bonds of the plaquettes which contain $b$, and, with the suprema over these plaquettes and bonds,
\begin{equation}\label{eq:current-remainder}
    \abs{\mathbf F(\mathbf U, \mathbf A)(b)} \leq C_8\bigl(\abs{\mathbf A} + \eta\abs{\nabla_{\mathbf U}\mathbf A}\bigr)\bigl(\abs{\nabla_{\mathbf U}\mathbf A} + \eta^{-2}\abs{\mathbf U(\partial p) - \id}\bigr)\;.
\end{equation}
Moreover the operator $\mathsf K(\mathbf U) \coloneqq D^*_{\mathbf U}D_{\mathbf U} + D_{\mathbf U}D^*_{\mathbf U} - \Delta_{\mathbf U}$ on $1$-forms is of the order zero: $(\mathsf K(\mathbf U)\mathbf A)(b)$ depends on $\mathbf A$ on the bonds of the plaquettes which contain $b$ only, and $\abs{(\mathsf K(\mathbf U)\mathbf A)(b)} \leq C_8\eta^{-2}\abs{\mathbf U(\partial p) - \id}\abs{\mathbf A}$.
\end{lemma}

The expansion is \cite[(3.11)]{BalabanRG1} with \cite[(1.43)--(1.54)]{BalabanREG}, and the statement on $\mathsf K$ is \cite[(135), p.~298]{BalabanVAR}.

\begin{definition}[The Formal Current]\label{def:formal-current}
Let a sequence of domains be given, with its operator $P = I - R$ of Definition~\ref{def:operators} at a configuration $\mathbf U$, and let $(\mathbf U, \mathbf J)$ be a pair and $\mathbf A$ a $1$-form with values in $\mathfrak g^c$. The \emph{formal current} of $e^{i\eta\mathbf A}\mathbf U$ is
\begin{equation}\label{eq:formal-current}
    \mathbf J_{\mathbf A} \coloneqq \mathbf J + \Delta_{\mathbf U}\mathbf A - DPD^*\mathbf A + \mathsf K(\mathbf U)\mathbf A + \mathbf F(\mathbf U, \mathbf A)\;.
\end{equation}
\end{definition}

Since $D^*D + DRD^* = \Delta_{\mathbf U} - DPD^* + \mathsf K(\mathbf U)$, the formal current at $\mathbf J = J(\mathbf U)$ is $J(e^{i\eta\mathbf A}\mathbf U) + DRD^*\mathbf A$, and it is the current of $e^{i\eta\mathbf A}\mathbf U$ when $\mathbf A$ is in the Landau gauge. The function $\mathbf F_1$ of \cite[(3.13)]{BalabanRG1}, for which no formula is printed, is $\mathsf K(\mathbf U)\mathbf A + \mathbf F(\mathbf U, \mathbf A)$, and $P_1 = DPD^*$.

\begin{proposition}[The Spaces of Pairs Are Stable]\label{prop:space-stability}
For all $\alpha_0, \alpha_1 \leq a_6$ there is a number $\alpha_2 > 0$, which depends on $\alpha_0$, $\alpha_1$, $L$, $G$, $M$, and $B_{\mathrm r}$ only, with the following property. Let $1 \leq j \leq k$, $X \in \mathbf D_j$, let $(\mathbf U, \mathbf J) \in \mathfrak U^c_j(X; \frac{1}{2}\alpha_0, \frac{1}{2}\alpha_1)$, and let $\mathbf A$ be a $1$-form on the bonds of $X$ with values in $\mathfrak g^c$ and
\begin{equation}\label{eq:shift-domain}
    a_j\abs{\mathbf A},\quad
    a_j^2\abs{\nabla_{\mathbf U}\mathbf A},\quad
    a_j^3\abs{\Delta_{\mathbf U}\mathbf A},\quad
    a_j^3\abs{DPD^*\mathbf A} \leq \alpha_2
\end{equation}
on $X$, where $P$ is the operator of a sequence of domains for which it is defined on $X$. Then $(e^{i\eta\mathbf A}\mathbf U, \mathbf J_{\mathbf A}) \in \mathfrak U^c_j(X; \alpha_0, \alpha_1)$. Consequently, for a term $\mathbf E(X, \cdot\,, z)$ of a family of regular terms with the constants $\alpha_0$, $\alpha_1$, $E_0$, $\kappa$, the function
\begin{equation}\label{eq:shifted-term}
    (\mathbf U, \mathbf J, \mathbf A) \longmapsto \mathbf E\bigl(X, (e^{i\eta\mathbf A}\mathbf U, \mathbf J_{\mathbf A}), z\bigr)
\end{equation}
is holomorphic on the set of these triples and bounded there by $E_0e^{-\kappa d_j(X)}$.
\end{proposition}

This is the assertion of \cite[p.~272]{BalabanRG1} that the function (3.13) is analytic for $\mathbf A$ satisfying (3.14), of which it is said that the restrictions ``can be easily obtained from the definition of the spaces''. For the conditions (i)--(iii) of Definition~\ref{def:pair-space} this is so, by Lemma~\ref{lem:current-expansion}; the condition (iv) requires the stability of the derived configurations under the shift. In the constructions the field $\mathbf A$ is a function of the pair, through the functions $\mathcal H$ of Definition~\ref{def:background-field} and their derivatives, and the bound of its Laplacian is the reason for which the second representation \eqref{eq:background-field} is used there.

\subsection{The Action of the Fluctuation as a Sum of Local Terms}\label{sec:localised-action}

In the coordinates of Theorem~\ref{thm:fluctuation-coordinates}, scaled by the coupling, the weight of the fluctuation field is the Gaussian weight of the form $\Delta^{(n)}$ times the exponential of a remainder. The remainder consists of the terms of the third and higher orders of the action and of the gauge-fixing function, of the logarithm of the Jacobian, and of the differences of the old terms between the old background and the new one. Each of its parts is a function of the fluctuation and of the new background on the whole lattice, through the operators $H$, $H_0$, $H_1$, $\tilde G$, $\mathfrak G$. It is made a sum of local functions by interpolating the walk expansions of these operators \cite[Sect.~1]{BalabanRG2}.

\begin{definition}[The Remainder of the Fluctuation]\label{def:fluctuation-remainder}
In the setting of Definition~\ref{def:step-setting}, let $g > 0$, let $A$ be a function on $\mathsf b^*(\mathcal O)$ with values in $\mathfrak g^c$, and put $B \coloneqq gCA$ and $B' \coloneqq B - h\tilde D(B)$. With the operators at a pair $(\mathbf U, \mathbf J)$, with $\Delta_1 \coloneqq \Delta(\mathbf U) - \Delta'_\pi - \Delta^{(2)}_\pi$, $\tilde D_3 \coloneqq \tilde D - \tilde C^{(2)}$, and $\mathbf G_3 \coloneqq \mathbf G - \frac{1}{2}\mathbf G^{(2)}$,
\begin{equation}\label{eq:fluctuation-remainder}
\begin{aligned}
    \mathbf P^{(n)}(g; A) \coloneqq{}& \log\det\Bigl(I - h\frac{\partial\tilde D}{\partial B}(B)\Bigr) + \sum_{b \in \mathsf b(\mathcal O)}\log\frac{\sigma(B'(b))}{\sigma_0}\\
    &+ \frac{1}{g^2}\Bigl[\langle H_1h\tilde D_3(B), \mathbf J\rangle + \langle H_1B, \Delta_1H_1h\tilde D(B)\rangle - \frac{1}{2}\langle H_1h\tilde D(B), \Delta_1H_1h\tilde D(B)\rangle\Bigr]\\
    &- \frac{1}{g^2}\Bigl[\mathsf W(B') + \mathbf G_3(B)\Bigr]\;,
\end{aligned}
\end{equation}
with the function $\mathsf W$ of \eqref{eq:fluctuation-potential} and the conventions of Definition~\ref{def:unit-form} for the functions on $\mathsf b(\mathcal O)$.
\end{definition}

By Proposition~\ref{prop:action-fluctuation}, by $\langle H_1B, J\rangle = 0$ for $\tilde QB = 0$, and since the gauge-fixing function does not depend on the central bonds, at the pair $(U_{n+1}, J(U_{n+1}))$ and for real $A$
\begin{equation}\label{eq:fluctuation-weight}
\begin{aligned}
    -\frac{1}{g^2}\Bigl[A\bigl(U_n(e^{iB'}V^{(n)})\bigr) - A(U_{n+1}) + \mathbf G\bigl(\Gamma_{n+1}, e^{iB'}V^{(n)}\bigr)\Bigr] &+ \log\frac{\mathcal J}{\mathcal J|_{A=0}}\\
    &= -\frac{1}{2}\langle A, C^*\Delta^{(n)}CA\rangle + \mathbf P^{(n)}(g; A)\;,
\end{aligned}
\end{equation}
which is \cite[(2.10)--(2.12), pp.~267--268]{BalabanRG1}; there the last term of the second line of \eqref{eq:fluctuation-remainder} is printed without the factor $\frac{1}{2}$.

\begin{proposition}[The Local Terms of the Fluctuation Action]\label{prop:localised-action}
There is a number $q \geq 0$, which depends on $L$ and $G$ only, and for every $\kappa \geq 1$ there are numbers $C_{10}(\kappa)$, $M(\kappa)$, and $\varepsilon^{\mathrm f}(\kappa, M) > 0$ with the following property. Let $M \geq M(\kappa)$, let the old terms of the scales $j \leq n$ be families of regular terms with constants $(\alpha_{0,j}, \alpha_{1,j}, E_0, \kappa)$, let $\Lambda \subset \mathcal O$ be as in Definition~\ref{def:unit-form}(iii), and let $\varepsilon^{\mathrm f} \leq \varepsilon^{\mathrm f}(\kappa, M)$. For every domain $Y \in \mathbf D_n$ which meets $\Lambda$ there is a function $\mathbf V(Y; (\mathbf U, \mathbf J), A)$ of the pairs on $Y$ and of the functions $A$ on the bonds of $\mathsf b^*(\Lambda)$ in $Y$, such that:
\begin{enumerate}[label=(\roman*)]
\item at the new background and for real $A$ with $g\abs A < \varepsilon^{\mathrm f}$, the sum of $\mathbf P^{(n)}(g; A)$ and of the differences of the old terms between the backgrounds $U_n(e^{iB'}V^{(n)})$ and $U_{n+1}$ equals $\sum_Y\mathbf V(Y; (U_{n+1}, J(U_{n+1})), A)$;
\item for $Y \subset \Lambda$ the function $\mathbf V(Y; \cdot)$ is holomorphic on the product of $\mathfrak U^c_{n+1}(Y; (1 + \beta)\alpha_{0,n+1}, (1 + \beta)\alpha_{1,n+1})$ and of the set of the $A$ with values in $\mathfrak g^c$ and $g\abs A < \varepsilon^{\mathrm f}$, it is invariant under the simultaneous gauge transformations of the pair and of $A$ with values in $G^c$ near $G$, and it does not depend on $\Lambda$, on the domains of the sequence outside $Y$, or on the torus;
\item $\mathbf V(Y; \cdot) = \frac{1}{2}\langle A, \mathcal Q(Y; A)A\rangle + \mathbf V''(Y; \cdot)$, where the kernel of the operator $\mathcal Q(Y; A)$ and the function $\mathbf V''$ satisfy
\begin{equation}\label{eq:localised-action-bounds}
\begin{aligned}
    \abs{\mathcal Q(Y; A)(b, b')} &\leq C_{10}(\kappa)\varepsilon^{\mathrm f}M^qe^{-2\kappa(d_n(Y) + \#(Y))}\;,\\
    \abs{\mathbf V''(Y; \cdot)} &\leq C_{10}(\kappa)E_0\varepsilon^{\mathrm f}M^qe^{-\frac{4}{5}\kappa d_n(Y)}\;.
\end{aligned}
\end{equation}
\end{enumerate}
For the domains $Y$ which are not contained in $\Lambda$ the same holds with the set $\tilde{\mathfrak U}^c_{n+1}(Y; (1 + \beta)\tilde\alpha_0, (1 + \beta)\tilde\alpha_1)$ in the place of the space of pairs and with the functions depending also on the fluctuation field on $Y \setminus \Lambda$, an outer variable.
\end{proposition}

For a sequence all of whose domains are the torus this is \cite[Lemma~2, p.~11]{BalabanRG2}, with $\varepsilon^{\mathrm f}$ for the threshold $\varepsilon_1$ of the fluctuation field there; the bounds \eqref{eq:localised-action-bounds} are \cite[(1.36), (1.43)]{BalabanRG2} with the exponents fixed, and the quadratic part is separated because the terms of $\mathbf P^{(n)}$ which carry the factor $g^{-2}$ are bounded only by a multiple of the square of the field. For a sequence with several layers the statement is the one which \cite[pp.~275--276, (3.43)]{BalabanCRE} describes, where no formulae are printed. The functions $\mathbf V(Y; \cdot)$ are not determined by (i): they are produced by a definite procedure, the interpolation in parameters attached to the cubes of the walks, and the comparison of two constructions on different fine lattices uses the procedure and not only its result.

\subsection{The Cluster Expansion}\label{sec:cluster-expansion}

The integral of the exponential of a sum of local terms is a sum over families of pairwise separated domains of products of local quantities, and the logarithm of such a sum is again a sum of local quantities when these are small. We state the second fact for an abstract system of activities, and then the result for the fluctuation integral.

\begin{definition}[Polymers]\label{def:polymers}
Let a grid be given. A \emph{polymer} is a nonempty connected union of cubes of the grid. Two polymers are \emph{compatible} if they share neither a cube nor a wall of a cube; for polymers $Z_1, \dots, Z_m$ let $\zeta(Z_i, Z_j) \coloneqq 1$ if $Z_i$ and $Z_j$ are compatible and $0$ otherwise, and
\begin{equation}\label{eq:ursell}
    \rho^T(Z_1, \dots, Z_m) \coloneqq \sum_{\mathsf g}\prod_{\{ i, j \} \in \mathsf g}\bigl(\zeta(Z_i, Z_j) - 1\bigr)\;,
\end{equation}
the sum over the connected graphs $\mathsf g$ on $\{ 1, \dots, m \}$, with $\rho^T(Z_1) \coloneqq 1$.
\end{definition}

\begin{theorem}[The Logarithm of a Polymer Sum]\label{thm:polymer-logarithm}
There are numbers $\kappa_2 \geq 1$ and $c_7 > 0$ with the following property. Let $\mathcal P$ be a finite set of polymers of a grid, let $\kappa \geq \kappa_2$ and $0 < \epsilon \leq c_7$, and let $\mathsf H \colon \mathcal P \to \mathbb C$ satisfy $\abs{\mathsf H(Z)} \leq \epsilon e^{-\kappa d(Z)}$. Then
\begin{equation}\label{eq:polymer-sum}
    \Xi \coloneqq 1 + \sum_{m \geq 1}\frac{1}{m!}\sum_{(Z_1, \dots, Z_m)}\prod_{i < j}\zeta(Z_i, Z_j)\prod_{i}\mathsf H(Z_i)
\end{equation}
does not vanish, and
\begin{equation}\label{eq:polymer-logarithm}
    \log\Xi = \sum_X\mathsf E(X)\;,
    \qquad
    \mathsf E(X) \coloneqq \sum_{m \geq 1}\frac{1}{m!}\sum_{(Z_1, \dots, Z_m)\colon \bigcup_iZ_i = X}\rho^T(Z_1, \dots, Z_m)\prod_i\mathsf H(Z_i)\;,
\end{equation}
where the sums run over the sequences of polymers of $\mathcal P$, the outer sum of the first identity over the connected unions $X$ of cubes, the series converge absolutely, and $\abs{\mathsf E(X)} \leq 2\epsilon e^{-(\kappa - \kappa_2)d(X)}$. If $\mathsf H$ depends holomorphically on parameters, with the bound uniform, then so do $\log\Xi$ and the $\mathsf E(X)$.
\end{theorem}

In \cite[(2.11)--(2.13), (2.39)--(2.41)]{BalabanRG2} this is taken from the literature on cluster expansions; it is the criterion of Koteck\'y and Preiss for polymers with the weight $e^{c\#(Z)}$, which the linear size controls by Lemma~\ref{lem:cube-count}.

\begin{theorem}[The Fluctuation Integral]\label{thm:fluctuation-integral}
Let $\delta_L \coloneqq \frac{1}{10}(L - 2)$. For every $\kappa \geq \kappa_2/\delta_L$ there are numbers $M'(\kappa)$, and for $M \geq M'(\kappa)$ numbers $\gamma^{\mathrm f} > 0$ and $C_{11}$, which depend on $L$, $G$, $\kappa$, $M$, and $E_0$, with the following property. In the setting of Proposition~\ref{prop:localised-action}, let $0 < g \leq \gamma^{\mathrm f}$, let the threshold $\varepsilon^{\mathrm f}$ satisfy $C_{11}\varepsilon^{\mathrm f} \leq 1$ and $(\varepsilon^{\mathrm f}/g)^2 \geq C_{11}$, let $\chi^{(n)}(A)$ be the product over $b \in \mathsf b^*(\Lambda)$ of functions $\chi_b(g\abs{A(b)})$ with values in $[0, 1]$ which equal one for $g\abs{A(b)} \leq \frac{1}{2}\varepsilon^{\mathrm f}$ and vanish for $g\abs{A(b)} \geq \varepsilon^{\mathrm f}$, and let the outer fluctuation field vanish. Then for every polymer $Z$ of the grid $\pi_{n+1}$ which meets $\Lambda$ there is a function $\mathsf H(Z; (\mathbf U, \mathbf J))$ such that:
\begin{enumerate}[label=(\roman*)]
\item at the new background,
\begin{equation}\label{eq:fluctuation-polymers}
    \int\chi^{(n)}(A)\exp\Bigl(\sum_Y\mathbf V(Y; A)\Bigr)\,\dd\mu_{C^{(n)}(\Lambda)}(A) = 1 + \sum_{m \geq 1}\frac{1}{m!}\sum_{(Z_1, \dots, Z_m)}\prod_{i < j}\zeta(Z_i, Z_j)\prod_i\mathsf H(Z_i)\;;
\end{equation}
\item for $Z \subset \Lambda$ the function $\mathsf H(Z; \cdot)$ is holomorphic on $\mathfrak U^c_{n+1}(Z; \alpha_{0,n+1}, \alpha_{1,n+1})$, invariant under the gauge transformations with values in $G^c$, and independent of $\Lambda$, of the domains of the sequence outside $Z$, and of the torus;
\item $\abs{\mathsf H(Z; \cdot)} \leq C_{11}\varepsilon^{\mathrm f}e^{-(1 + \delta_L)\kappa d_{n+1}(Z)}$.
\end{enumerate}
Consequently the logarithm of the integral is the sum of the functions $\mathsf E(X)$ of \eqref{eq:polymer-logarithm}, which for $X \subset \Lambda$ have the properties \emph{(ii)} and satisfy $\abs{\mathsf E(X)} \leq 2C_{11}\varepsilon^{\mathrm f}e^{-\kappa d_{n+1}(X)}$.
\end{theorem}

This is \cite[(2.1)--(2.13) and Lemma~3, p.~20]{BalabanRG2}: the Mayer expansion of the exponential, the expansion of the characteristic functions, the conditioning of the Gaussian measure on the complement of a domain, the interpolation of the covariances and of their square root, and the resummation into activities of the polymers of the next scale; the gain from the rate $\frac{4}{5}\kappa$ of the local terms at the scale $n$ to $(1 + \delta_L)\kappa$ at the scale $n + 1$ is the inequality $2d_n(Z) \geq Ld_{n+1}(Z')$ of \cite[(2.36)]{BalabanRG2}, with the choice of the losses in \cite[p.~21]{BalabanRG2}. The condition $(\varepsilon^{\mathrm f}/g)^2 \geq C_{11}$ makes the Gaussian weight of a field at the threshold small, and it is the reason for which the threshold of the fluctuation is a power of the logarithm of the coupling times the coupling. The functions $\chi_b$ may be characteristic functions.

\subsection{The Step in the Region of Small Fields}\label{sec:small-field-step}

The statements of this section combine to the theorem that a step reproduces the form of the regular part of the effective action \cite[Theorem~3, p.~264]{BalabanRG1}. We state it for the terms which lie inside the region of small fields of a term of the density, which is the generality of \cite[pp.~272--279]{BalabanCRE}.

\begin{theorem}[The Regular Terms of a Step]\label{thm:small-field-step}
Let $\mathfrak g$ be simple. There are numbers $\kappa_3$, and for $\kappa \geq \kappa_3$ numbers $M''(\kappa)$, and for $M \geq M''(\kappa)$ positive numbers $\gamma$, $\bar\alpha_0$, $\bar\alpha_1$, $E_0$, which depend on $L$, $G$, $\kappa$, and $M$, with the following property. Let $0 \leq n < k$, let $g_0, \dots, g_n \in (0, \gamma]$, let $\alpha_{0,j} \leq \bar\alpha_0$ and $\alpha_{1,j} \leq \bar\alpha_1$ for $j \leq n + 1$ and let the threshold $\varepsilon^{\mathrm f}_n$ of the fluctuation field be as in Theorem~\ref{thm:fluctuation-integral}, with the relations between these numbers which are collected in Section~10. Suppose that for $1 \leq j \leq n$ families $\mathbf E^{(j)}$ of regular terms of the scale $j$ with the constants $(\alpha_{0,j}, \alpha_{1,j}, E_0, \kappa)$ and numbers $\beta_j$ are given, with $\beta_j = \beta(\mathbf E^{(j)})$ up to the error of Proposition~\ref{prop:polarisation}(iv). Then there is a family $\mathbf E^{(n+1)}$ of regular terms of the scale $n + 1$ with the constants $(\alpha_{0,n+1}, \alpha_{1,n+1}, E_0, \kappa)$ such that:
\begin{enumerate}[label=(\roman*)]
\item for the sequence all of whose domains are the torus and for regular data $V$ on $T^{(n+1)}$, the sum $\sum_z\mathbf E^{(n+1)}(U_{n+1}, z)$ over $z \in T^{(n+1)}$ equals, up to a constant, the logarithm of the first factor of \eqref{eq:jacobian}, plus the logarithm of the normalisation $\int e^{-\frac{1}{2}\langle A, C^*\Delta^{(n)}CA\rangle}\dd A$ of the Gaussian measure, plus the logarithm of the integral over the fluctuation field, with the measure $\mu_{C^{(n)}(\mathbb T)}$, of $\chi^{(n)}(A)$ times the exponential of
\begin{equation}\label{eq:small-field-integrand}
    \mathbf P^{(n)}(g_n; A) + \sum_{j=1}^{n}\sum_{z \in T^{(j)}}\Bigl[\mathbf E^{(j)}(U, z) - \beta_jA(h^{(j)}_z, U)\Bigr]_{U = U_{n+1}}^{U = U_n(e^{iB'}V^{(n)})}\;,
\end{equation}
where $h^{(j)}_z$ is the partition of unity \eqref{eq:base-point-partition} and $A(h, U)$ is the action with the plaquette at $x$ weighted by $h(x)$;
\item for a term of the density with a region $\Lambda$ of small fields of the step, the same integral over the fluctuation field in $\Lambda$ has the logarithm $\sum_z\sum_X\mathbf E^{(n+1)}(X, \cdot, z)$, summed over the domains $X \subset \Lambda$ and their base points, plus terms which are indexed by domains that meet the complement of $\Lambda$;
\item the terms $\mathbf E^{(n+1)}(X, \cdot\,, z)$ do not depend on the term of the density, on $\Lambda$, or on the torus, and they depend continuously on $g_0, \dots, g_n$ and holomorphically on $g_n$ in a sector, when the thresholds are kept fixed.
\end{enumerate}
\end{theorem}

Here the partition of unity of the base points is
\begin{equation}\label{eq:base-point-partition}
    h^{(j)}_z(x) \coloneqq \prod_{\mu}\max\Bigl\{ 1 - \frac{\abs{x_\mu - z_\mu}}{a_j}, 0 \Bigr\}\;,
    \qquad
    z \in T^{(j)}\;,\ x \in T_\eta\;,
\end{equation}
of \cite[(3.40), p.~275]{BalabanCRE}. The theorem contains the cancellation which defines the coefficients: by Proposition~\ref{prop:polarisation} and the Ward identities the part of $\mathbf E^{(j)}(U, z)$ of the second, third, and fourth order in the field of the scale $j$ which does not decay with the scale is $\beta(\mathbf E^{(j)})$ times the corresponding part of $A(h^{(j)}_z, U)$ \cite[(3.57)--(3.67), pp.~281--283]{BalabanCRE}, so that the bracket in \eqref{eq:small-field-integrand} is a sum of terms which are small in the ratio of the scales $j$ and $n$. The continuity in the couplings is not stated in the papers; it holds because the couplings enter the integrand of the step only through the factors $g_n$, the old terms, and the coefficients, once the thresholds no longer depend on them.

\subsection{Gauge Invariance and the Coefficient of the Action}\label{sec:coefficient}

Among the terms of the expansion of an old term in the new fluctuation there are terms of the second, the third, and the fourth order in the field whose bounds do not improve with the scale. Gauge invariance and the symmetry of the lattice reduce all of them to a multiple of the action, with one coefficient, which is removed by changing the coupling \cite[Sects.~4, 5]{BalabanRG1}, \cite[(3.49)--(3.67)]{BalabanCRE}. We state the identities which express the gauge invariance, define the coefficient, and state the form of the second derivative. In this subsection the scale $j$ is fixed, and a \emph{unit field} is a function on the positive bonds of $T^{(j)}$ with values in $\mathfrak g^c$, extended to the negative bonds by $B(b^{-1}) = -B(b)$.

\begin{lemma}[Ward Identities]\label{lem:ward}
Let $\mathbf E$ be a holomorphic function on a neighbourhood of the trivial configuration in the configurations on the bonds of a finite subset of $T^{(j)}$, invariant under the gauge transformations with values in $G$, and write $\mathbf E^{(n)}$ for the derivative of the order $n$ of $B \mapsto \mathbf E(e^{iB})$ at $B = 0$, a symmetric $n$-linear form on unit fields. Then $\mathbf E$ is invariant under the gauge transformations with values in $G^c$ near $\id$, and for all unit fields $B_1$ and all functions $\lambda$ on the sites with values in $\mathfrak g^c$, with $(\partial\lambda)(b) \coloneqq \lambda(b_+) - \lambda(b_-)$:
\begin{enumerate}[label=(\roman*)]
\item $\langle\mathbf E^{(1)}, i[\lambda(b_-), B_1(b)]\rangle = \langle\mathbf E^{(2)}, \partial\lambda, B_1\rangle + \frac{1}{2}\langle\mathbf E^{(1)}, i[B_1, \partial\lambda]\rangle$;
\item if the centre of $\mathfrak g$ is trivial, then $\mathbf E^{(1)} = 0$, and consequently $\langle\mathbf E^{(2)}, \partial\lambda, B_1\rangle = 0$.
\end{enumerate}
\end{lemma}

These are \cite[(4.13)--(4.15), p.~284]{BalabanRG1}; the identities of the third and of the fourth order, which the proof of Theorem~\ref{thm:small-field-step} uses, are derived in the same way from \cite[(4.9)]{BalabanRG1}. The second derivative of a term at the trivial data is the kernel from which the coefficient is read off.

\begin{definition}[The Polarisation and the Coefficient]\label{def:polarisation}
Let $\mathbf E$ be a family of regular terms of the scale $j$ on the torus $\mathbb T$, and let the sequence of domains be the one with $\Omega_1 = \dots = \Omega_j = \mathbb T$. For a unit field $B$ with values in $\mathfrak g$ and $\abs B \leq C_{\mathrm d}\varepsilon_1$ let $U_j(e^{iB})$ be the configuration \eqref{eq:moved-minimiser} at the trivial base, $V_0 = \id$ and $U_0 = \id$.
\begin{enumerate}[label=(\roman*)]
\item The \emph{polarisation} of the family with the base point $z \in T^{(j)}$ is the bilinear form on unit fields
\begin{equation}\label{eq:polarisation}
    \Pi_z(B_1, B_2) \coloneqq \frac{\partial^2}{\partial s\,\partial t}\,\mathbf E\bigl(U_j(e^{i(sB_1 + tB_2)}), z\bigr)\Big|_{s=t=0}\;,
\end{equation}
with the sum \eqref{eq:regular-covariance}, and $\Pi \coloneqq \sum_z\Pi_z$.
\item If the centre of $\mathfrak g$ is trivial and $\mathfrak g$ is simple, then $\Pi(B_1, B_2) = \sum_{b, b'}\Pi(b, b')\tr B_1(b)B_2(b')$ with real numbers $\Pi(b, b')$, and we write $\Pi_{\mu\nu}(x) \coloneqq \Pi(\langle 0, a_je_\mu\rangle, \langle x, x + a_je_\nu\rangle)$ for $x \in T^{(j)}$.
\item The \emph{coefficient} of the family is
\begin{equation}\label{eq:coefficient}
    \beta(\mathbf E) \coloneqq \sum_{x \in T^{(j)}}\Pi_{12}(x)\,\frac{x_1x_2}{a_j^2}\;,
\end{equation}
the coordinates of $x$ being taken in $(-\frac{1}{2}\ell_{\mathbb T}, \frac{1}{2}\ell_{\mathbb T}]$.
\end{enumerate}
\end{definition}

In (ii) the form $\Pi(b, b')$ on $\mathfrak g \times \mathfrak g$ is invariant under the adjoint action, because the terms and the minimiser are covariant under constant gauge transformations (Lemma~\ref{lem:problem-covariance}), and an invariant bilinear form on a simple Lie algebra is a multiple of the trace form. The definition is \cite[(1.20)--(1.22), p.~264]{BalabanRG1} and \cite[(3.50), (3.57), (3.61)]{BalabanCRE} on a torus. In the papers the coefficient is the limit of \eqref{eq:coefficient} as the torus grows, and it is used as a number which does not depend on the torus; the dependence of \eqref{eq:coefficient} on the torus is exponentially small in its side, and it has to be controlled by the statement on the cluster expansion.

\begin{proposition}[The Form of the Polarisation]\label{prop:polarisation}
Let $\mathfrak g$ be simple, and let $\mathbf E$ be a family of regular terms of the scale $j$ with the constants $\alpha_0$, $\alpha_1$, $E_0$, $\kappa$, with $\kappa \geq \kappa_1$ for a number $\kappa_1$ which depends on $L$ and $M$ only. Then, with a number $\delta_4 > 0$ which depends on $L$, $G$, $M$, and $\kappa$, and a number $C_9$ which depends on these and on $\alpha_0$, $\alpha_1$:
\begin{enumerate}[label=(\roman*)]
\item $\abs{\Pi(b, b')} \leq C_9E_0e^{-\delta_4\abs{b_- - b'_-}/a_j}$;
\item $\Pi(rb, rb') = \Pi(b, b')$ for the isometries $r$ of the torus which preserve $T^{(j)}$, and $\Pi(b, b') = \Pi(b', b)$;
\item $\sum_b\Pi(b, b')(\partial\lambda)(b) = 0$ for every real function $\lambda$ on $T^{(j)}$;
\item for all $\mu, \nu, \kappa', \lambda'$, each up to an error of at most $C_9E_0e^{-\delta_4\ell_{\mathbb T}/(4a_j)}$,
\begin{equation}\label{eq:polarisation-moments}
\begin{gathered}
    \sum_{x \in T^{(j)}}\Pi_{\mu\nu}(x)\,\frac{x_{\kappa'}x_{\lambda'}}{a_j^2} = \beta(\mathbf E)\bigl(\delta_{\mu\kappa'}\delta_{\nu\lambda'} + \delta_{\mu\lambda'}\delta_{\nu\kappa'} - 2\delta_{\mu\nu}\delta_{\kappa'\lambda'}\bigr)\;,\\
    \sum_{x \in T^{(j)}}\Pi_{\mu\nu}(x)\,\frac{x_{\kappa'}}{a_j} = 0\;,
    \qquad
    \sum_{x \in T^{(j)}}\Pi_{\mu\nu}(x) = 0\;.
\end{gathered}
\end{equation}
\end{enumerate}
Here $\Pi(b, b')$ is extended to the negative bonds by $\Pi(b^{-1}, b') = -\Pi(b, b') = \Pi(b, b'^{-1})$.
\end{proposition}

The parts (i)--(iii) are \cite[(5.4)--(5.10), pp.~292--293]{BalabanRG1}, and (iv) is the content of \cite[(5.16)]{BalabanRG1} which the renormalisation uses, in the form of \cite[(3.60)--(3.64)]{BalabanCRE}: the tensor of the second moments is invariant under the reflections and the permutations of the coordinates and is annihilated by the contraction which (iii) gives, and such a tensor is a multiple of the one on the right, which is the tensor of the second moments of the kernel of the quadratic part of the action \cite[(4.41)--(4.43)]{BalabanRG1}.

\section{The Step on a Term}\label{sec:step}

The density of the field after $n$ steps is a sum of terms, and a term is labelled by the regions in which, at each of the scales, the steps could be performed. The next step treats each term separately. It decides, by decompositions of unity, in which part of the remaining region of small fields the field of the scale $n$ is regular enough for the fluctuation integral of Section~\ref{sec:small-field}; there it passes to the coordinates of the fluctuation and integrates; in the rest it leaves the field as an integration variable of the new term. The operations are those of \cite[Sect.~3]{BalabanCRE}, which repeat those of the first step \cite[Sect.~1]{BalabanCRE}. We define them in the order in which they are applied, with the test functions in the soft form which the comparison of two constructions requires.

\subsection{Parameters, Cubes, and Labels}\label{sec:labels}

\begin{definition}[Parameters and Cubes]\label{def:parameters}
A \emph{system of parameters} consists of powers $M$ and $M_2$ of $L$ with $M_{\mathrm c} \leq M_2$ and $L^2M_2 \leq M$; for each scale $0 \leq j \leq k$ of a power $R_j$ of $L$, with $R_{j+1} \leq R_j \leq LR_{j+1}$, and of positive numbers $g_j$, $\varepsilon_j$, $\vartheta_j$, $\alpha_{0,j}$, $\alpha_{1,j}$; and of a number $\kappa_s \in (0, 1]$. For $0 \leq j \leq k$:
\begin{enumerate}[label=(\roman*)]
\item $\mathsf r_j$ is the grid of side $MR_ja_j$, $\mathsf s_j$ the grid of side $LM_2R_ja_j$, and $\mathsf m_j$ the grid of side $M_2R_ja_j$;
\item for a union $X$ of cubes of $\mathsf r_j$ the operations $X^{\sim i}$ are those of the grid $\mathsf r_j$, and for a union $X$ of cubes of $\mathsf s_j$ they are those of the grid $\mathsf m_j$ (Definition~\ref{def:localisation}(ii));
\item for a union $X$ of cubes of $\mathsf m_j$ with $j \geq 1$, $X^*$ is the set of the positive bonds of $T^{(j-1)}$ with at least one end in $X$ which are not central bonds $b_0(c)$ of bonds $c$ of $T^{(j)}$ with an end in $X$.
\end{enumerate}
\end{definition}

The number $g_j$ is the coupling of the scale $j$, $R_j$ fixes the size of the cubes on which the field of the scale $j$ is tested, $\varepsilon_j$ is the threshold of the plaquettes and $\vartheta_j$ that of the deviation from a reference configuration and of the fluctuation field, and $\alpha_{0,j}$, $\alpha_{1,j}$ are the radii of the spaces of pairs. In \cite[(2.4), (2.5), (2.28); p.~265]{BalabanCRE} they are functions of the coupling: $\varepsilon_j = g_jA_0(\log g_j^{-2})^{p_0}$, $R_j$ the smallest power of $L$ not less than $(\log g_j^{-2})^r$, $\vartheta_j = (A_1/A_0)\varepsilon_j$, which is written $\delta_j$ there. Here they are independent data, and the relations which the theorems need are stated where they are used. By the condition on the numbers $R_j$ the grids $\mathsf r_j$, $\mathsf s_j$, $\mathsf m_j$ of all scales are nested: a cube of a grid of a coarser scale, or of a larger side at the same scale, is a union of cubes of the other.

\begin{definition}[Labels]\label{def:label}
Let $0 \leq n \leq k$. A \emph{label of the order $n$} consists of sets $\Omega_j$, $\Lambda_j$, $S_j$ for $1 \leq j \leq n$ such that
\begin{enumerate}[label=(\roman*)]
\item $\Omega_j$ and $\Lambda_j$ are unions of cubes of $\pi_j$, and $\Omega_1 \supset \Lambda_1 \supset \Omega_2 \supset \dots \supset \Omega_n \supset \Lambda_n$;
\item the sequence $\mathbb T \supset \Omega_1 \supset \dots \supset \Omega_n$, continued by empty sets, is a sequence of domains with the block size $M$ and a separation of at least $27$;
\item $S_j$ is a union of cubes of $\mathsf s_j$ contained in $\Omega_j$ and in the closure of the complement of $\Lambda_j$.
\end{enumerate}
The \emph{region of large fields} of the scale $j$ is the closure $Z_j$ of the complement of $\Lambda_j$, with $Z_0 \coloneqq \emptyset$, and the determining set of the label is that of the sequence (ii). The label of the order $0$ is empty.
\end{definition}

These are the sequences \cite[(2.1)--(2.3), pp.~254--255]{BalabanCRE}. The fields of a label are a gauge field $V_j$ on the bonds of each layer $\Gamma_j$, $0 \leq j \leq n$, which together are data $V$ on the determining set, and the fluctuation fields of the earlier steps on the regions $\Omega_j \setminus \Lambda_j$.

\subsection{Local Minimisers and Test Functions}\label{sec:test-functions}

The regularity of a field of the scale $j$ near a cube is tested on the configuration of least action which has this field as its average of the order $j$ near the cube and which, close to the boundary of a neighbourhood of the cube, is tied to the field itself on the finer scales \cite[(2.13)--(2.17), pp.~256--257]{BalabanCRE}.

\begin{definition}[The Minimiser of a Cube]\label{def:cube-minimiser}
Let $1 \leq j \leq k$.
\begin{enumerate}[label=(\roman*)]
\item For a gauge field $V_j$ on $T^{(j)}$ the \emph{straight extension} is the gauge field $Q^{s*}_jV_j$ on $T_\eta$ which on a positive bond $b$ equals $\id$ if both ends of $b$ lie in one block $B^j(y)$, $y \in T^{(j)}$, and equals $V_j(c)$ if $b_- \in B^j(c_-)$ and $b_+ \in B^j(c_+)$ for a bond $c$ of $T^{(j)}$.
\item For a cube $\square$ of $\mathsf s_j$ let $W \coloneqq \square^{\sim4}$, and let $V_{\square}$ be the data on the determining set $\mathbf B_j(W)$ of the minimal sequence of $W$ (Definition~\ref{def:minimal-sequence}) which on the layer $i$ equal $M^i(Q^{s*}_jV_j)$. If these data are regular with a constant $\varepsilon_1 \leq a^{\mathrm v}_1$, then $U_{j,\square}(V_j)$ is the minimal configuration for them on the support $W$ which is in the axial gauge (Theorem~\ref{thm:support}), and
\begin{equation}\label{eq:plaquette-test}
    f_\square(V_j) \coloneqq \sup_{p \subset \square^\sim}\Bigl(\frac{a_j}{\eta}\Bigr)^2\abs{U_{j,\square}(V_j)(\partial p) - \id}\;,
\end{equation}
the supremum over the plaquettes of $T_\eta$ in $\square^\sim$; otherwise $f_\square(V_j) \coloneqq \infty$. For $j = 0$ and a gauge field $V_0$ on $T_\eta$, $f_\square(V_0)$ is the supremum of $\abs{V_0(\partial p) - \id}$ over the plaquettes in $\square^\sim$.
\end{enumerate}
\end{definition}

The average of the order $j$ of the straight extension is $V_j$, and its averages of the lower orders are straight extensions of $V_j$ to the intermediate lattices, by the definition of the average: the transports inside the blocks are trivial, and the loops of Definition~\ref{def:block-paths} have trivial holonomy. So the data $V_\square$ equal $V_j$ on the top layer and are determined by $V_j$ on the bonds of $T^{(j)}$ which meet $W$.

\begin{definition}[Soft Test Functions]\label{def:soft-test}
Let $\bar\chi(x) \coloneqq \min\{ 1, \max\{ 0, (1 + \kappa_s - x)/\kappa_s \} \}$ for $x \in [0, \infty]$. For a finite family of functions $f_\square$ with values in $[0, \infty]$, indexed by the cubes of a grid in a region $D$, a number $t > 0$, and a union $\mathsf P \subset D$ of cubes,
\begin{equation}\label{eq:soft-term}
    \chi\bigl(D \setminus \mathsf P; f, t\bigr) \coloneqq \prod_{\square \subset D,\ \square \not\subset \mathsf P}\bar\chi\Bigl(\frac{f_\square}{t}\Bigr)\;,
    \qquad
    \chi^c\bigl(\mathsf P; f, t\bigr) \coloneqq \prod_{\square \subset \mathsf P}\Bigl(1 - \bar\chi\Bigl(\frac{f_\square}{t}\Bigr)\Bigr)\;,
\end{equation}
so that $\sum_{\mathsf P}\chi(D \setminus \mathsf P; f, t)\chi^c(\mathsf P; f, t) = 1$, the sum over all unions of cubes in $D$. The \emph{test function of the step $j$} on a union $X$ of cubes of $\mathsf s_j$ is $\chi_j(X) \coloneqq \chi(X; f, \varepsilon_j)$ with the functions \eqref{eq:plaquette-test}.
\end{definition}

A factor $\bar\chi(f_\square/t)$ vanishes unless $f_\square < (1 + \kappa_s)t$ and equals one for $f_\square \leq t$, and a factor $1 - \bar\chi(f_\square/t)$ vanishes unless $f_\square > t$. For $\kappa_s \to 0$ these are the characteristic functions of \cite[(2.17), (3.2)]{BalabanCRE}. The identity is the expansion of the product of the sums $\bar\chi + (1 - \bar\chi) = 1$ over the cubes.

\subsection{The Regions of a Step}\label{sec:step-regions}

We fix $0 \leq n < k$ and a label of the order $n$, with its region $Z_n$ of large fields, and describe how the step $n + 1$ divides the torus. All cubes are those of the scale $n + 1$, and the fields are a gauge field $V_n$ on $T^{(n)}$ and its average $V_{n+1} = M(V_n)$.

\begin{definition}[The Regions of a Step]\label{def:step-regions}
For a union $X$ of cubes of $\mathsf s_{n+1}$ let $X'$ denote the union of the cubes of $\mathsf r_{n+1}$ which meet $X$; likewise $Z'_n$ is the union of the cubes of $\mathsf r_{n+1}$ which meet $Z_n$.
\begin{enumerate}[label=(\roman*)]
\item \emph{The plaquettes of the new field.} $D_1 \coloneqq \mathbb T \setminus Z'^{\sim4}_n$, and $\mathsf P$ ranges over the unions of cubes of $\mathsf s_{n+1}$ in $D_1$, with the functions $f_{\square'}(V_{n+1})$ of \eqref{eq:plaquette-test} at the scale $n + 1$ and the threshold $\varepsilon_{n+1}$.
\item \emph{The gauge fixing.} $D_2 \coloneqq \mathbb T \setminus (\mathsf P'^\sim \cup Z'^{\sim5}_n)$ and $\mathsf P^1 \coloneqq D_2 \cap T^{(n+1)}$.
\item \emph{The deviation from the reference field.} $D_3 \coloneqq D_2^{\sim-1}$. For a cube $\square'$ of $\mathsf s_{n+1}$ in $D_3$ with $f_{\square'}(V_{n+1}) < \infty$ let $V^{(n)}_{\square'} \coloneqq M^n(U_{n+1,\square'}(V_{n+1}))$ and
\begin{equation}\label{eq:deviation-test}
    f'_{\square'}(V_n) \coloneqq \sup_{b \in (\square'^{\sim2})^*}\abs{V_n(b)V^{(n)}_{\square'}(b)^{-1} - \id}\;,
\end{equation}
and $f'_{\square'} \coloneqq \infty$ if $f_{\square'}(V_{n+1}) = \infty$. $\mathsf Q$ ranges over the unions of cubes of $\mathsf s_{n+1}$ in $D_3$, with the functions $f'$ and the threshold $2\vartheta_n$.
\item \emph{The new region of small fields of the background.}
\begin{equation}\label{eq:new-omega}
    \Omega_{n+1} \coloneqq \mathbb T \setminus \bigl(\mathsf Q'^{\sim2} \cup (\mathbb T \setminus D_2)^{\sim3}\bigr)\;.
\end{equation}
\item \emph{The fluctuation field.} With the coordinates $A$ of the fluctuation on $\mathsf b^*(\Omega_{n+1})$ (Theorem~\ref{thm:fluctuation-coordinates}, scaled by $g_n$), $\mathsf R$ ranges over the unions of cubes of $\mathsf s_{n+1}$ in $\Omega^{\sim-1}_{n+1}$, with the functions $f''_{\square'}(A) \coloneqq \sup_{b \in (\square'^{\sim2})^*}g_n\abs{A(b)}$, the threshold $\vartheta_n$, and the characteristic functions of $\{ f'' < \vartheta_n \}$ and of its complement in the place of the soft functions.
\item \emph{The new region of small fields.} $\Lambda_{n+1}$ is obtained from $\Omega^{\sim-2}_{n+1} \setminus \mathsf R'^\sim$ by adding to its complement, for every connected component of the complement which lies in a cube of side $100MR_{n+1}a_{n+1}$, the remaining points of the smallest union of cubes of $\mathsf r_{n+1}$ of the form of a product of intervals which contains the component; and
\begin{equation}\label{eq:new-s}
    S_{n+1} \coloneqq \bigl(\Omega_{n+1} \setminus \Omega^{\sim-1}_{n+1}\bigr) \cup \mathsf R\;.
\end{equation}
\end{enumerate}
\end{definition}

These are the constructions \cite[(3.2)--(3.5), (3.16), (3.20)]{BalabanCRE}. The region $\Omega_{n+1}$ is a union of cubes of $\mathsf r_{n+1}$ at a distance of at least two such cubes from the regions in which a test function of (i) or (iii) has failed, from the region in which the gauge is not fixed, and at a distance of at least eight from $Z_n$; with $\Lambda_{n+1}$ and $S_{n+1}$ it extends the label to a label of the order $n + 1$ when $R_{n+1} \geq L$, since $8LR_{n+1} > 27$.

\begin{lemma}[The Old Test Functions on the New Region]\label{lem:old-tests}
There are numbers $c_8 > 0$ and $C_{12}$, which depend on $L$ and $G$ only, with the following property. Suppose that $\vartheta_n \leq c_8\varepsilon_n$, $\varepsilon_{n+1} \leq 2\varepsilon_n \leq c_8$, $\kappa_s \leq \frac{1}{4}$, and $C_{12}e^{-\delta_2M_2R_n} \leq \vartheta_n/\varepsilon_n$. Let $V_n$ be a gauge field such that the factors of the test functions of Definition~\ref{def:step-regions}(i) and (iii) for the cubes of $\Omega^{\sim2}_{n+1}$ do not vanish. Then:
\begin{enumerate}[label=(\roman*)]
\item $f_\square(V_n) \leq \varepsilon_n$ for every cube $\square$ of $\mathsf s_n$ in $\Omega^\sim_{n+1}$, so that the factors $\bar\chi(f_\square/\varepsilon_n)$ of the test function $\chi_n$ of the label equal one there, and $\abs{V_n(y, x) - \id} < \varepsilon_n$ for $y \in T^{(n+1)} \cap \Omega^\sim_{n+1}$ and $x \in B(y)$ when $V_n$ is in the axial gauge of the step;
\item the data which consist of the data of the label on the layers $j < n$, of $V_n$ on the bonds of $T^{(n)}$ in $\Omega_n \setminus \Omega_{n+1}$, and of $V_{n+1}$ on $\mathsf c(\Omega_{n+1})$ are regular with the constant $2L^2(1 + \kappa_s)\varepsilon_n$ on a support of the new sequence, the minimal configuration $U_{n+1}$ for them satisfies $\abs{U_{n+1}(\partial p) - \id} \leq 2B_3\varepsilon_{n+1}(\eta/a_{n+1})^2$ on $\Omega_{n+1}$, and $\abs{V_n(b)V^{(n)}(b)^{-1} - \id} \leq C_{12}\varepsilon_n$ for $b \in \mathsf b(\Omega_{n+1})$, with $V^{(n)}$ of \eqref{eq:evaluation-field};
\item for a cube $\square'$ of $\mathsf s_{n+1}$ in $\Omega^{\sim-1}_{n+1}$ and $b \in (\square'^{\sim2})^*$, $\abs{V^{(n)}(b)V^{(n)}_{\square'}(b)^{-1} - \id} \leq C_{12}e^{-\delta_2M_2R_n}\varepsilon_n$, so that $f'_{\square'} \leq 2\vartheta_n$ whenever $f''_{\square'}(A) < \vartheta_n$.
\end{enumerate}
\end{lemma}

These are \cite[(3.6)--(3.9), (3.12)--(3.14), (3.17)--(3.19), pp.~266--269]{BalabanCRE}: the local minimisers of the two scales and the global one are compared by Theorem~\ref{thm:background-field}, whose function $\mathcal H$ has an argument supported near the boundary of the cube, at a distance of at least $M_2R_n$ in the units of the scale, and by Lemma~\ref{lem:axial-smallness}.

\subsection{The Identity of the Step}\label{sec:step-identity}

We can now state what the step does to the integrand of a term. The statement is an identity between integrals for an arbitrary gauge-invariant factor in the integrand, so that it applies to the density and to any observable which is carried along with it.

\begin{definition}[The Weights of a Step]\label{def:step-weights}
In the setting of Definition~\ref{def:step-regions}, with the number $z_n \coloneqq z(g_n, \varepsilon_n)$ of \eqref{eq:gauge-normalisation}:
\begin{enumerate}[label=(\roman*)]
\item for a region $\Omega_{n+1}$, the function $\zeta(\Omega^c_{n+1})$ of $V_n$ and $V_{n+1}$ is the sum, over the pairs $(\mathsf P, \mathsf Q)$ for which \eqref{eq:new-omega} is $\Omega_{n+1}$, of the product of the factors of
\[
    \chi(D_1 \setminus \mathsf P; f, \varepsilon_{n+1})\,\chi^c(\mathsf P; f, \varepsilon_{n+1})\,\chi(D_3 \setminus \mathsf Q; f', 2\vartheta_n)\,\chi^c(\mathsf Q; f', 2\vartheta_n)
\]
which belong to cubes not contained in $\Omega_{n+1}$, and of
\[
    z_n^{-(L^4-1)\#(\mathsf P^1 \setminus \Omega_{n+1})}\,\chi_{\mathrm{Ax}}(\mathsf P^1 \setminus \Omega_{n+1}, V_n)\exp\Bigl(-\frac{1}{g_n^2}\mathbf G(\mathsf P^1 \setminus \Omega_{n+1}, V_n)\Bigr)\;;
\]
\item $\chi'_n(X) \coloneqq \chi(X; f', 2\vartheta_n)$ for a union $X$ of cubes of $\mathsf s_{n+1}$ in $\Omega_{n+1}$, and $\chi^{(n)}(X)$, $\chi^{(n)c}(X)$ are the products over the cubes of $X$ of the characteristic functions of Definition~\ref{def:step-regions}(v);
\item the constant of the step and the factor of the central bonds are
\begin{equation}\label{eq:step-constant}
\begin{aligned}
    E^{(n)}_0(\Omega_{n+1}) &\coloneqq \#\mathsf b^*(\Omega_{n+1})\bigl(\dim\mathfrak g\,\log g_n + \log\sigma_0\bigr) - (L^4 - 1)\#(\Omega_{n+1} \cap T^{(n+1)})\log z_n\;,\\
    z^{(n)} &\coloneqq \prod_{c \in \mathsf c(\Omega_{n+1})}\abs{\det(Lq_c)}^{-1}\;,
\end{aligned}
\end{equation}
with the maps $q_c$ of Definition~\ref{def:unit-average} at the data of Lemma~\ref{lem:old-tests}(ii).
\end{enumerate}
\end{definition}

\begin{theorem}[The Identity of the Step]\label{thm:step-identity}
Let the hypotheses of Lemma~\ref{lem:old-tests} and of Theorems~\ref{thm:fluctuation-coordinates} and~\ref{thm:background-field} hold for the parameters, with $C_{12}\varepsilon_n \leq \frac{1}{4}a_8$. Fix a label of the order $n$ and the fields of the label other than $V_n$ on $\Omega_n$. Let $\Phi$ be a bounded Borel function of the gauge fields $V_n$ on the bonds of $T^{(n)}$ in $\Omega_n$ which is invariant under the gauge transformations of $T^{(n)}$ with values in $G$ that equal $\id$ at the points of $T^{(n+1)}$ and outside $\Omega_n$, and let $f$ be a bounded Borel function of the gauge fields on the bonds of $T^{(n+1)}$ in $\Omega_n$. Then
\begin{equation}\label{eq:step-identity}
\begin{aligned}
    \int f\bigl(M(V_n)\bigr)&\,\chi_n(\Omega_n)\,\Phi(V_n)\,\dd V_n\\
    ={}& \sum_{\Omega_{n+1}}\sum_{\mathsf R}\int\!\!\int f\,\chi_{n+1}(\Omega_{n+1})\,\zeta(\Omega^c_{n+1})\,\chi_n(\Omega_n \setminus \Omega^\sim_{n+1})\,e^{E^{(n)}_0(\Omega_{n+1})}z^{(n)}\\
    &\times\int\chi^{(n)}\bigl(\Omega^{\sim-1}_{n+1} \setminus \mathsf R\bigr)\,\chi^{(n)c}(\mathsf R)\,\chi'_n(S_{n+1})\,e^{-g_n^{-2}\mathbf G(\Gamma_{n+1}, V_n)}\\
    &\qquad\times\frac{\mathcal J}{\mathcal J|_{A=0}}\,\Phi\bigl(V_n\bigr)\,\dd A\,\dd W\,\dd V^{\mathrm{out}}_n\;,
\end{aligned}
\end{equation}
where $V^{\mathrm{out}}_n$ is the field on the bonds of $T^{(n)}$ in $\Omega_n$ with no end in $\Omega_{n+1}$, $W$ is the field on $\mathsf c(\Omega_{n+1})$, $A$ is a function on $\mathsf b^*(\Omega_{n+1})$ with values in $\mathfrak g$, the argument of $f$ on the right is the field which equals $W$ on $\mathsf c(\Omega_{n+1})$ and $M(V^{\mathrm{out}}_n)$ on the other bonds, and the field $V_n$ on the right is $V^{\mathrm{out}}_n$ outside $\mathsf b(\Omega_{n+1})$ and
\begin{equation}\label{eq:step-field}
    V_n = \exp\bigl(i\bigl[g_nCA - h\tilde D(g_nCA)\bigr]\bigr)\,V^{(n)}
\end{equation}
on $\mathsf b(\Omega_{n+1})$, with the objects of Section~\ref{sec:fluctuation-field} for the new sequence and the data of Lemma~\ref{lem:old-tests}(ii). The sums run over the regions \eqref{eq:new-omega} and over the unions $\mathsf R$ of cubes of $\mathsf s_{n+1}$ in $\Omega^{\sim-1}_{n+1}$, and the integral over $A$ over the set on which the test functions of the fluctuation do not vanish, which lies in the domain of Theorem~\ref{thm:fluctuation-coordinates}.
\end{theorem}

This is \cite[(3.15), (3.16), pp.~267--268]{BalabanCRE} for a general factor. The steps of the proof are the decompositions of unity (i), (iii), and (v) of Definition~\ref{def:step-regions}; the insertion of the gauge fixing by Lemma~\ref{lem:gauge-fixing}, which uses the invariance of $\Phi$ and of the other factors; Lemma~\ref{lem:old-tests}, which removes the factors of the old test functions and of the gauge fixing on $\Omega^\sim_{n+1}$ and those of the deviation on the cubes outside $S_{n+1}$; and Theorem~\ref{thm:fluctuation-coordinates} at fixed outer fields, by which the integral over the field on $\mathsf b(\Omega_{n+1})$ becomes the integral over $W$ and $A$. When $\Phi$ contains the factor $\exp(-g_n^{-2}A(U_n(V_n)))$ of the density, with the background of the label, \eqref{eq:fluctuation-weight} writes the product of this factor, of the gauge-fixing weight on $\Omega_{n+1}$, and of the Jacobian as
\begin{equation}\label{eq:step-gaussian}
    \exp\Bigl(-\frac{1}{g_n^2}A(U_{n+1}) - \frac{1}{2}\langle A, C^*\Delta^{(n)}CA\rangle + \mathbf P^{(n)}(g_n; A)\Bigr)\;,
\end{equation}
and the remaining factors of $\Phi$ are evaluated at the field \eqref{eq:step-field}; this is the form \cite[(3.15)]{BalabanCRE}.

\subsection{The Fluctuation near the Boundary of the New Region}\label{sec:boundary-change}

One factor of \eqref{eq:step-identity} is not yet local. The function $\chi'_n(S_{n+1})$ compares the field \eqref{eq:step-field} with the reference fields of the cubes of $S_{n+1}$, and the field \eqref{eq:step-field} contains $V^{(n)}$, which depends on the new field $W$ in the whole of $\Omega_{n+1}$. A term of the next density, however, may depend on $W$ in the new region of small fields through the background alone. On the bonds near $S_{n+1}$ the field $V^{(n)}$ is therefore replaced by the corresponding field of a minimiser which is determined by the fields in the complement of $\Lambda_{n+1}$, and the difference is absorbed by a change of the variables $A(b)$ \cite[pp.~251--252, 269]{BalabanCRE}.

\begin{definition}[The Boundary Field and the Change of Variables]\label{def:boundary-field}
In the setting of Theorem~\ref{thm:step-identity}, with the sets $\Lambda_{n+1}$ and $S_{n+1}$ of Definition~\ref{def:step-regions}:
\begin{enumerate}[label=(\roman*)]
\item $\mathsf Y$ is the closure of $\Lambda_n \setminus \Lambda_{n+1}$, and $U_{\mathsf Y}$ is the minimal configuration, on the support $\mathsf Y$ and in the axial gauge, of the sequence $\{ \Omega_i \cap \Omega_i(\mathsf Y) \}_{i \leq n+1}$ of the new label localised to $\mathsf Y$ (Definition~\ref{def:localised-sequence}), for the data which on the bonds of the layer $i$ are $M^i(Q^{s*}_{n+1}V_{n+1})$ in $\Omega_{n+1}$ and $M^i(Q^{s*}_nV_n)$ elsewhere; $V^{(n)}_{\mathsf Y} \coloneqq M^n(U_{\mathsf Y})$.
\item $\mathsf S$ is the set of the bonds of $\mathsf b^*(\Omega_{n+1})$ with an end in $S_{n+1}$, and
\begin{equation}\label{eq:boundary-field}
    \mathbf H^{(n)}(b) \coloneqq \frac{1}{i}\operatorname{Log}\bigl(V^{(n)}(b)\,V^{(n)}_{\mathsf Y}(b)^{-1}\bigr)\;,
    \qquad
    b \in \mathsf S\;.
\end{equation}
\item Let $\varphi_1$ be a $C^\infty$ function on $[0, \infty)$ with values in $[0, 1]$ which vanishes on $[0, \frac{4}{3}]$ and equals one on $[\frac{5}{3}, \infty)$, fixed once and for all. For $b \in \mathsf S$ with an end in $\mathsf R^{\sim2}$ put $\omega_b(A') \coloneqq \varphi_1(g_n\abs{A'}/\vartheta_n)$, and $\omega_b \coloneqq 1$ for the other bonds of $\mathsf S$. For $A' \in \mathfrak g$ such that the logarithm is defined,
\begin{equation}\label{eq:boundary-change}
    \psi_b(A') \coloneqq \frac{1}{ig_n}\operatorname{Log}\Bigl(e^{ig_nA'}e^{-i\omega_b(A')\mathbf H^{(n)}(b)}\Bigr)\;,
\end{equation}
and $\Psi_{n+1}$ is the map of the functions on $\mathsf b^*(\Omega_{n+1})$ which acts as $\psi_b$ on the variable of each bond $b \in \mathsf S$ and as the identity on the others.
\item For a cube $\square'$ of $\mathsf s_{n+1}$ in $S_{n+1}$,
\begin{equation}\label{eq:local-deviation-test}
    f'^{\,\mathrm{loc}}_{\square'}(A') \coloneqq \sup_{b \in (\square'^{\sim2})^*}\bigl\lvert e^{ig_nA'(b)}\,V^{(n)}_{\mathsf Y}(b)\,V^{(n)}_{\square'}(b)^{-1} - \id\bigr\rvert\;,
\end{equation}
and $\chi'^{\,\mathrm{loc}}_n(S_{n+1}) \coloneqq \chi(S_{n+1}; f'^{\,\mathrm{loc}}, 2\vartheta_n)$.
\end{enumerate}
\end{definition}

The configuration $U_{\mathsf Y}$ depends on $V_n$ and on $W$ on the bonds which meet $\mathsf Y$ only. In the interior of $\mathsf Y$ the localised sequence has the layers of the new label, of the scale $n + 1$ in $\Omega_{n+1}$ and of the scale $n$ outside, and near the two parts of the boundary of $\mathsf Y$ it has the layers of the minimal sequences; this is the determining set $\mathbf B(\Lambda^c_{k+1} \cap \Lambda_k) \cup \mathbf B_{k+1}$ of \cite[p.~269]{BalabanCRE}. The function $\varphi_1$ is the bump function of \cite[p.~252]{BalabanCRE}, which is specified there on two sets only; here it is a fixed function of the unscaled field $g_nA'$ in the units of the threshold, so that its support does not move with the coupling.

\begin{proposition}[The Change of Variables near the Boundary]\label{prop:boundary-change}
There are numbers $c_{10} > 0$ and $C_{13}$, which depend on $L$ and $G$ only, with the following property. Let the hypotheses of Theorem~\ref{thm:step-identity} hold, with $\vartheta_n \leq c_{10}$ and $C^2_{13}\varepsilon_ne^{-\delta_2M_2R_{n+1}} \leq \vartheta_n$, and fix a term of \eqref{eq:step-identity} and fields $V^{\mathrm{out}}_n$ and $W$ at which its test functions do not vanish. Then:
\begin{enumerate}[label=(\roman*)]
\item $\abs{\mathbf H^{(n)}(b)} \leq C_{13}\varepsilon_ne^{-\delta_2M_2R_{n+1}}$ for $b \in \mathsf S$;
\item $\Psi_{n+1}$ is a $C^\infty$ diffeomorphism of an open set of functions $A'$ onto an open set which contains every function $A$ at which the characteristic functions of the fluctuation and the function $\chi'_n(S_{n+1})$ of the term do not vanish;
\item at $A = \Psi_{n+1}(A')$ the characteristic functions of the fluctuation have the same values as at $A'$, and where they do not vanish the function $\chi'_n(S_{n+1})$, evaluated at the field \eqref{eq:step-field}, equals $\chi'^{\,\mathrm{loc}}_n(S_{n+1})$ at $A'$.
\end{enumerate}
Consequently, for every bounded Borel function $F$ of $A$,
\begin{equation}\label{eq:boundary-change-identity}
    \int\chi^{(n)}\chi^{(n)c}(\mathsf R)\,\chi'_n(S_{n+1})\,F(A)\,\dd A = \int\chi^{(n)}\chi^{(n)c}(\mathsf R)\,\chi'^{\,\mathrm{loc}}_n(S_{n+1})\,F\bigl(\Psi_{n+1}(A')\bigr)\,\abs{\det D\Psi_{n+1}(A')}\,\dd A'\;,
\end{equation}
with the characteristic functions of the term on the two sides.
\end{proposition}

This is the operation of \cite[p.~269]{BalabanCRE}, which repeats \cite[(1.21)--(1.24), pp.~251--252]{BalabanCRE}. On a bond with $\omega_b = 1$ the field \eqref{eq:step-field} is $e^{ig_nA}V^{(n)} = e^{ig_nA'}V^{(n)}_{\mathsf Y}$, by \eqref{eq:boundary-field} and \eqref{eq:boundary-change}. Where $\omega_b < 1$ we have $g_n\abs{A'(b)} < \frac{5}{3}\vartheta_n$, and by (i) the tested quantities of both functions are below $2\vartheta_n$ at this bond, where the soft functions equal one; where $\omega_b > 0$ we have $g_n\abs{A'(b)} > \frac{4}{3}\vartheta_n$, and both $A(b)$ and $A'(b)$ exceed the threshold $\vartheta_n/g_n$ of the fluctuation. After the change of variables the term depends on $W$ in $\Lambda_{n+1}$ through the configuration $U_{n+1}$ and the operators formed with it, and no longer through a test function. The factor $\abs{\det D\Psi_{n+1}}$ and the change of the other factors of the integrand under $\Psi_{n+1}$ are, in the papers, a new part $V^{(k)}(S_{k+1}, A_k, \mathbf H^{(k)})$ of the action.

\subsection{The Operation of a Step and the Density after It}\label{sec:step-density}

After the change of variables the fluctuation field in $\Lambda_{n+1}$ occurs in no test function except its own, and it can be integrated. The Gaussian weight couples it to the fluctuation field on the bonds outside $\Lambda_{n+1}$, which remains a variable of the new term, and the coupling is removed by completing the square.

\begin{lemma}[The Splitting of the Gaussian Weight]\label{lem:gaussian-splitting}
Let $\mathsf B = \mathsf B^{\mathrm{in}} \sqcup \mathsf B^{\mathrm{out}}$ be a finite set, let $\Gamma$ be a real symmetric form on the functions on $\mathsf B$ with values in $\mathfrak g$ whose restriction to the functions supported in $\mathsf B^{\mathrm{in}}$ is positive definite, and let $C^{\mathrm{in}}$ be the inverse of this restriction, extended by zero to the other functions. Then, for $A = A^{\mathrm{in}} + A^{\mathrm{out}}$ and $a \coloneqq C^{\mathrm{in}}\Gamma A^{\mathrm{out}}$,
\begin{equation}\label{eq:gaussian-splitting}
    \langle A, \Gamma A\rangle = \bigl\langle A^{\mathrm{in}} + a, \Gamma\bigl(A^{\mathrm{in}} + a\bigr)\bigr\rangle + \bigl\langle A^{\mathrm{out}}, \bigl(\Gamma - \Gamma C^{\mathrm{in}}\Gamma\bigr)A^{\mathrm{out}}\bigr\rangle\;,
\end{equation}
and $\int\exp(-\frac{1}{2}\langle A, \Gamma A\rangle)\,\dd A^{\mathrm{in}} = \det(2\pi C^{\mathrm{in}})^{1/2}\exp(-\frac{1}{2}\langle A^{\mathrm{out}}, (\Gamma - \Gamma C^{\mathrm{in}}\Gamma)A^{\mathrm{out}}\rangle)$.
\end{lemma}

\begin{proof}
The function $a$ is supported in $\mathsf B^{\mathrm{in}}$, and $C^{\mathrm{in}}$ inverts $\Gamma$ on such functions, so that $\langle A^{\mathrm{in}}, \Gamma a\rangle = \langle A^{\mathrm{in}}, \Gamma A^{\mathrm{out}}\rangle$ and $\langle a, \Gamma a\rangle = \langle A^{\mathrm{out}}, \Gamma C^{\mathrm{in}}\Gamma A^{\mathrm{out}}\rangle$. Hence the first term on the right of \eqref{eq:gaussian-splitting} is $\langle A^{\mathrm{in}}, \Gamma A^{\mathrm{in}}\rangle + 2\langle A^{\mathrm{in}}, \Gamma A^{\mathrm{out}}\rangle + \langle A^{\mathrm{out}}, \Gamma C^{\mathrm{in}}\Gamma A^{\mathrm{out}}\rangle$, and adding the second term gives the expansion of $\langle A, \Gamma A\rangle$. The integral is that of a Gaussian function of $A^{\mathrm{in}} + a$ with the covariance $C^{\mathrm{in}}$, and the Lebesgue measure is invariant under the translation by $a$.
\end{proof}

\begin{definition}[The Operation of a Step]\label{def:step-operation}
In the setting of Theorem~\ref{thm:step-identity} and Definition~\ref{def:boundary-field}, fix $\Omega_{n+1}$ and $\mathsf R$, hence $\Lambda_{n+1}$ and $S_{n+1}$. Let $\Gamma \coloneqq C^*\Delta^{(n)}C$ on the functions on $\mathsf b^*(\Omega_{n+1})$, let $\mathsf B^{\mathrm{in}} \coloneqq \mathsf b^*(\Lambda_{n+1})$, so that $C^{\mathrm{in}} = C^{(n)}(\Lambda_{n+1})$, and write $A = A^{\mathrm{in}} + A^{\mathrm{out}}$. Let $\mathbf T_n$ be the operation of a label of the order $n$, with $\mathbf T_0$ the identity. For a function $F$ of the fields of the label, of $V^{\mathrm{out}}_n$, of $W$, and of $A^{\mathrm{out}}$,
\begin{equation}\label{eq:step-operation}
\begin{aligned}
    \mathbf T^{(n+1)}[F] \coloneqq \int\zeta(\Omega^c_{n+1})\,&\chi_n(\Omega_n \setminus \Omega^\sim_{n+1})\,z^{(n)}_{\mathrm{out}}\,\chi^{(n)}\bigl((\Omega^{\sim-1}_{n+1} \setminus \mathsf R) \setminus \Lambda_{n+1}\bigr)\,\chi^{(n)c}(\mathsf R)\,\chi'^{\,\mathrm{loc}}_n(S_{n+1})\\
    &\times\exp\Bigl(-\frac{1}{2}\bigl\langle A^{\mathrm{out}}, \bigl(\Gamma - \Gamma C^{(n)}(\Lambda_{n+1})\Gamma\bigr)A^{\mathrm{out}}\bigr\rangle\Bigr)\,F\,\dd A^{\mathrm{out}}\;,
\end{aligned}
\end{equation}
where $z^{(n)}_{\mathrm{out}}$ is the product of the factors of $z^{(n)}$ of the bonds $c$ with an end outside $\Lambda_{n+1}$, and $\mathbf T_{n+1}[F]$ is the image, under the map which sends the fields to the gauge field $V_{n+1}$ of Theorem~\ref{thm:step-identity}, of the measure $\mathbf T_n[\mathbf T^{(n+1)}[F]\,\dd V^{\mathrm{out}}_n]\,\dd W$.
\end{definition}

This is \cite[(3.23), (3.24), p.~270]{BalabanCRE}. The new fields of the label are the field $V_n$ on the bonds of the layer of the scale $n$, which is integrated with its average as the field of the next scale outside $\Omega_{n+1}$, and the fluctuation field on the bonds of $\Omega_{n+1}$ which are not in $\Lambda_{n+1}$. By Lemma~\ref{lem:gaussian-splitting} what remains of the Gaussian weight \eqref{eq:step-gaussian} is the weight of the field $A^{\mathrm{in}} + C^{(n)}(\Lambda_{n+1})\Gamma A^{\mathrm{out}}$ with the covariance $C^{(n)}(\Lambda_{n+1})$, and the integral over $A^{\mathrm{in}}$ is the fluctuation integral of Section~\ref{sec:small-field} in the region $\Lambda_{n+1}$, with further terms in the exponent: the boundary terms and the terms of the operation of Section~\ref{sec:r-operation} of the earlier steps at the shifted background, the terms which contain the outer field $A^{\mathrm{out}}$, and the new part of the action of Proposition~\ref{prop:boundary-change}. Its logarithm is divided into three parts \cite[(3.26)--(3.47), pp.~271--279]{BalabanCRE}: the terms whose domains lie in $\Lambda_{n+1}$ and which come from the regular terms are those of Theorem~\ref{thm:small-field-step}; the terms with domains in $\Lambda_{n+1}$ which come from the terms of the operation of the scales $j$ with $L^{-(n-j)}MR_j > L$ form a new family of such terms; and all others are new boundary terms. The result is the following statement, in which the classes are those of Section~\ref{sec:inductive}.

\begin{theorem}[The Step on a Density]\label{thm:step-density}
Let the constants and an admissible system of parameters be as in Theorem~\ref{thm:inductive}, let $0 \leq n < k$, and let $\varrho_n$ be a density of the order $n$ with the constants of that theorem. Let $\mathbf E^{(n+1)}$ be the family of regular terms of Theorem~\ref{thm:small-field-step} for the regular terms of $\varrho_n$, and suppose that its coefficient $\beta_{n+1}$ satisfies $g_n^{-2} - \beta_{n+1} = g_{n+1}^{-2}$. Then the image $T\varrho_n$ of $\varrho_n$ under the average $V_n \mapsto M(V_n)$ is a density of the order $n + 1$ with the same constants, in which:
\begin{enumerate}[label=(\roman*)]
\item the labels are the labels of $\varrho_n$ extended by the sets $\Omega_{n+1}$, $\Lambda_{n+1}$, $S_{n+1}$ of Definition~\ref{def:step-regions}, the test functions are those of Definition~\ref{def:soft-test}, and the operations are those of Definition~\ref{def:step-operation};
\item the families of the scales $j \leq n$ are those of $\varrho_n$, the family of regular terms of the scale $n + 1$ is $\mathbf E^{(n+1)}$, and the coefficient of the scale $n + 1$ is $\beta_{n+1}$;
\item the families $\mathbf R^{(n+1)}$ and $\mathbf B^{(n+1)}$ are holomorphic on the spaces with the radii $(1 + \beta)\alpha_{0,j}$, $(1 + \beta)\alpha_{1,j}$ in the place of $\alpha_{0,j}$, $\alpha_{1,j}$, they satisfy their bounds with $(1 + 4\beta)\kappa$ in the place of $\kappa$, the bound of $\mathbf R^{(n+1)}$ holds with $\frac{1}{2}\varrho_{n+1}$, and the terms $\mathbf R^{(n+1)}(X, \cdot)$ do not depend on the label, on $k$, or on the torus.
\end{enumerate}
\end{theorem}

This is the first half of the inductive step of \cite[Theorem~1, p.~262]{BalabanCRE}, proved in \cite[Sect.~3]{BalabanCRE}; the improved constants of (iii) are those of \cite[p.~262]{BalabanCRE}, and they are the room which the operation of the next section needs. We call a density of the order $n + 1$ with the properties (iii) \emph{fresh}.

\section{The Operation on Small Regions of Large Fields}\label{sec:r-operation}

On a region in which a test function has failed the action supplies a small factor, of the order of $\exp(-(\varepsilon_j/g_j)^2)$ per cube, and the terms of the effective action near the region are only bounded. In four dimensions the coupling decreases too slowly for one such factor to pay for the bounded terms of all the later steps. The remedy is to integrate a region of large fields out once it has become small on the scale of the step and no new large field has appeared in it for many steps: the density is replaced by one with the same integral in which the region is a region of small fields, and the price of the replacement, a ratio of two integrals over the region, is a product of local factors close to one, whose logarithm joins the effective action \cite{BalabanLFR1,BalabanLFR2}. The operation has three parts. Preparatory integrations, which are identities, shrink a treated region to a box of a few blocks of the last scale, and they decide, by further decompositions of unity, whether the region is treated at all. The field of the last scale in the box is then integrated out and replaced by a normalised weight, which changes the density and keeps its integral. Finally the result is brought back to the form of a density, with new terms in the effective action.

\subsection{The Regions Which Are Treated}\label{sec:treated-regions}

Throughout this section a system of parameters and an order $1 \leq n \leq k$ are fixed, together with the data of the operation: an integer $N \geq 1$, the \emph{depth}; numbers $\delta'_j > 0$ for $0 \leq j \leq n$, the thresholds of the fluctuation fields of the operation; a number $\beta \in (0, \frac{1}{4}]$; and a number $M_0 > 0$. We assume $L \geq 28$, and we put $L_0 \coloneqq 2$ and $h \coloneqq n - N$. Differences of unions of cubes of a grid are understood as unions of cubes.

\begin{definition}[Small and Old Components]\label{def:treated-components}
Let a label of the order $n$ be given. Suppose that $h \geq 1$ and that there is an integer $N_0$ with $2 \leq N_0 < N$ and $L^{N_0-1} = R_{n-N_0+1}$; let $N_0$ be the smallest one, and $n_0 \coloneqq n - N_0$. A connected component $X$ of the region $Z_n$ of large fields is \emph{small and old} if
\begin{enumerate}[label=(\roman*)]
\item $X$ is contained in a cube of side $100MR_na_n$ which is a union of cubes of $\mathsf r_n$;
\item for $h \leq j \leq n$ the set $Z_j \cap X$ is a product of intervals, contained in a cube of side $100MR_ja_j$ which is a union of cubes of $\mathsf r_j$;
\item for $h \leq j < n$, with $W_j$ the union of the cubes of $\mathsf r_{j+1}$ which meet $Z_j \cap X$ and with the operations of the grid $\mathsf r_{j+1}$,
\begin{equation}\label{eq:old-component}
    X \setminus \Omega_{j+1} = W_j^{\sim8}\;,
    \qquad
    Z_{j+1} \cap X = W_j^{\sim10}\;,
    \qquad
    S_{j+1} \cap X = X \cap \bigl(\Omega_{j+1} \setminus \Omega^{\sim-1}_{j+1}\bigr)\;.
\end{equation}
\end{enumerate}
If $h < 1$ or there is no such $N_0$, no component is small and old.
\end{definition}

These are the conditions (i), (ii) of \cite[p.~177]{BalabanLFR1}. By Definition~\ref{def:step-regions} the equalities \eqref{eq:old-component} hold exactly when no test function of the step $j + 1$ has failed near $X$, so that (iii) says that no new large field was created in $X$ in the last $N$ steps. In \cite{BalabanLFR1} the regions $Z_j \cap X$ of these steps are products of intervals by the rule of Definition~\ref{def:step-regions}(vi); we require it of the sets themselves, which excludes a region with several components. The number $N_0$ is that of \cite[p.~179]{BalabanLFR1}: the grid $\mathsf r_{n_0+1}$ has the side $MR_{n_0+1}a_{n_0+1} = Ma_n$, so that it is the grid $\pi_n$.

\begin{definition}[The Reduced Regions]\label{def:reduced-regions}
Let $X$ be small and old. With the operations of the grid $\mathsf r_j$ in the first line and of $\mathsf r_{n_0+1} = \pi_n$ in the second,
\begin{equation}\label{eq:reduced-regions}
\begin{aligned}
    Z''_j &\coloneqq Z_j \cap X \quad (j \leq h)\;,
    &\qquad
    Z''_j &\coloneqq W_{j-1}^{\sim3} \quad (h < j \leq n_0)\;,\\
    Z''_{n_0+1} &\coloneqq W_{n_0}^{\sim1}\;,
    &\qquad
    Z''_n &\coloneqq W_{n_0}^{\sim3}\;,
\end{aligned}
\end{equation}
and for $n_0 + 1 < j < n$, in decreasing order, $Z''_j$ is the union of the cubes of $\pi_{j+1}$ in $Z''_{j+1}$ which touch no cube of $\pi_{j+1}$ outside $Z''_{j+1}$. Further $\Omega''_j \coloneqq \Omega_j$ for $j \leq h$ and
\begin{equation}\label{eq:reduced-domains}
    \Omega''_j \coloneqq (\Omega_j \setminus X) \cup (X \setminus Z''_j) \quad (h < j \leq n)\;,
    \qquad
    \Lambda_X \coloneqq W_{n_0}^{\sim4}\;,
    \qquad
    \Lambda^\circ_X \coloneqq \Lambda_X \setminus W_h^{\sim1}\;,
\end{equation}
with the operations of $\pi_n$ for the \emph{box} $\Lambda_X$ and of $\mathsf r_{h+1}$ for $W_h$.
\end{definition}

By \eqref{eq:old-component} these are the sets $Z''_j = (\Omega^{\sim5}_j)^c \cap X$, $Z''_{n_0+1} = (\Omega^{\sim7}_{n_0+1})^c \cap X$, $Z''_n = (\Omega^{\sim5}_{n_0+1})^c \cap X$, and $\Lambda = (\Omega^{\sim4}_{n_0+1})^c \cap X$ of \cite[(1.10)--(1.12), (1.73)]{BalabanLFR1}, written as enlargements of the regions of large fields, and with a definite choice of the sets between $Z''_{n_0+1}$ and $Z''_n$, which are only required there to form an admissible sequence.

\begin{lemma}[The Geometry of the Reduced Regions]\label{lem:reduced-regions}
Let $X$ be small and old. Then:
\begin{enumerate}[label=(\roman*)]
\item $Z''_h \subset Z''_{h+1} \subset \dots \subset Z''_n \subset \Lambda_X \subset X$, and each of these sets with an index $j > h$ is a product of intervals;
\item $\Lambda_X$ is a union of cubes of $\pi_n$ contained in a cube of side $100Ma_n$, at a distance of at least $6Ma_n$ from the complement of $X$ and of at least $4Ma_n$ from $\Omega_n$;
\item $\{ \Omega''_j \}$ is a sequence of domains with the block size $M$ and the separation $27$, which agrees with $\{ \Omega_j \}$ outside $X$ and for $j \leq h$, and $\Omega''_j \supset \Omega_j$.
\end{enumerate}
\end{lemma}

\begin{proof}
(i) An enlargement of a product of intervals by layers of cubes is a product of intervals, and so is the union of the cubes of a coarser grid which meet it. For $h < j \leq n_0$ we have $Z''_j = W_{j-1}^{\sim3} \subset W_{j-1}^{\sim10} = Z_j \cap X \subset W_j \subset W_j^{\sim3}$, by \eqref{eq:old-component}, which gives the inclusions up to $Z''_{n_0} \subset W_{n_0} \subset Z''_{n_0+1}$, and $Z''_h = Z_h \cap X \subset W_h \subset Z''_{h+1}$. The sets $Z''_j$ with $n_0 + 1 < j < n$ are obtained from $Z''_n$ by removing layers of the widths $Ma_n, Ma_{n-1}, \dots, Ma_{j+1}$, whose sum is less than $\frac{L}{L-1}Ma_n$, while $Z''_{n_0+1}$ lies in $Z''_n$ at the distance $2Ma_n$ from its complement; since $2 - \frac{L}{L-1} \geq L^{-1}$, each of them contains $Z''_{n_0+1}$ at a distance of at least $Ma_{n-1} \geq Ma_{j}$ from its complement. Finally $Z''_n = W_{n_0}^{\sim3} \subset W_{n_0}^{\sim4} = \Lambda_X \subset W_{n_0}^{\sim10} \subset X$.

(ii) $\Lambda_X \subset W_{n_0}^{\sim10} = Z_{n_0+1} \cap X$, which lies in a cube of side $100MR_{n_0+1}a_{n_0+1} = 100Ma_n$ and in $X$, and six layers of cubes of $\pi_n$ separate $\Lambda_X$ from its complement. By \eqref{eq:old-component}, $X \setminus \Omega_{n_0+1} = W_{n_0}^{\sim8}$, so that four layers separate $\Lambda_X$ from $\Omega_{n_0+1} \supset \Omega_n$.

(iii) Let $h \leq j < n$. Outside $X$ the two sequences agree, and in $X$ the complement of $\Omega''_j$ is $Z''_j$ for $j > h$ and $X \setminus \Omega_h \subset Z''_h$ for $j = h$, while $\Omega''_{j+1} \cap X = X \setminus Z''_{j+1}$. For $h < j < n_0$ the sets $Z''_j \subset W_j$ and $X \setminus Z''_{j+1} = X \setminus W_j^{\sim3}$ are separated by three layers of cubes of $\mathsf r_{j+1}$, of the width $3MR_{j+1}a_{j+1} \geq 3LMa_j$, and the same holds for $j = h$ and for $j = n_0$, with one layer of the width $Ma_n \geq LMa_{n_0}$ in the last case. For $n_0 < j < n$ they are separated by a layer of cubes of $\pi_{j+1}$, of the width $LMa_j$, by the proof of (i). Since $L \geq 28$, the distance exceeds $27Ma_j$ in every case. The inclusion $\Omega''_j \supset \Omega_j$ holds because $Z''_j \subset W_{j-1}^{\sim8} = X \setminus \Omega_j$ for $h < j \leq n_0$ and $Z''_j \subset W_{n_0}^{\sim3} \subset X \setminus \Omega_{n_0+1}$ for $j > n_0$.
\end{proof}

The preparatory integrations pass from the sequence of the label to the sequence $\{ \Omega''_j \}$ one scale at a time, and each of them has its own background \cite[(1.14)--(1.21)]{BalabanLFR1}.

\begin{definition}[The Successive Sequences and Backgrounds]\label{def:successive-sequences}
Let $X$ be small and old, $0 \leq m \leq N$, and $j \coloneqq h + m$.
\begin{enumerate}[label=(\roman*)]
\item The \emph{$m$-th sequence} is $\Omega^{(m)}_i \coloneqq \Omega''_i$ for $i \leq j$ and $\Omega^{(m)}_i \coloneqq \Omega_i$ for $j < i \leq n$, and $\mathbf B^{(m)}$ is its determining set; $\mathbf B^{(0)}$ is that of the label and $\mathbf B^{(N)}$ that of $\{ \Omega''_i \}$.
\item The data $V^{(m)}$ on $\mathbf B^{(m)}$ are defined recursively: $V^{(0)}$ are the data of the label, and $V^{(m+1)}$ equals $V^{(m)}$ on the bonds which belong to both determining sets and the average $M(V^{(m)}_j)$ of the field of the layer $j$ on the bonds of the layer $j + 1$ of $\mathbf B^{(m+1)}$ which do not belong to $\mathbf B^{(m)}$.
\item $U^{(m)}_n$ is the minimal configuration for $V^{(m)}$ in the axial gauge, and $U^{(m)}_{n,X}$ is the minimal configuration, on the support $X$ and in the axial gauge, of the sequence $\{ \Omega^{(m)}_i \cap \Omega_i(X) \}_{i \leq n}$ localised to $X$ (Definition~\ref{def:localised-sequence}), for the data which are $V^{(m)}$ on the bonds of its layers which belong to $\mathbf B^{(m)}$ and $M^i(Q^{s*}_nV_n)$ on the other bonds of its layers $i$.
\item $D_j \coloneqq X \cap (\Omega''_{j+1} \setminus \Omega_{j+1})$ for $j < n$, and $D^*_j$ is the set of the positive bonds of the layer $j$ of $\mathbf B^{(m)}$ with an end in $D_j$ which are not central bonds $b_0(c)$ of bonds $c$ of $T^{(j+1)}$.
\end{enumerate}
\end{definition}

The $m$-th sequence differs from the one before it in the domain of the scale $j$ alone, which grows by $D_{j-1}$: the field of the scale $j - 1$ in $D_{j-1}$ is replaced by its average. The other bonds of the layers of the localised sequence lie near the boundary of $X$, in $\Omega_n$, where the field of the label is $V_n$; this is the determining set $\mathbf B^{(n)}_k(Z)$ of \cite[(1.19), (1.20)]{BalabanLFR1}.

\subsection{The Preparatory Integrations}\label{sec:preparatory}

For $j = h, \dots, n - 1$ in turn the field of the scale $j$ in $D_j$ is integrated with its average held fixed, as in a step, together with the fluctuation field of the step $j + 1$ of the label in $X$. Before each integration the field is tested, so that the integral is one of small fields, and the test functions of the label on the scales $j$ and $j + 1$ are removed; a component in which a test fails acquires a new region of large fields and is not treated. The tests of the backgrounds form a cascade, in which each function is removed on the support of the next, and whose thresholds leave, at each stage, the room which the removal of the previous one consumes \cite[(1.22)--(1.24), pp.~181--182]{BalabanLFR1}.

\begin{definition}[The Cascade of Test Functions]\label{def:cascade}
Let $X$ be small and old, $0 \leq m \leq N$, and $j \coloneqq h + m$. For a plaquette $p$ of $T_\eta$ in $X$ the \emph{threshold} $T_m(p)$ is the smallest of the following numbers which apply to $p$, the regions being closed, and $T_m(p) \coloneqq \infty$ if none applies.
\begin{enumerate}[label=(\roman*)]
\item For $h \leq i < j$ and $p \subset Z''_{i+1} \setminus Z''_i$, where for $i = h$ the set $Z''_h$ is replaced by $X \setminus \Omega_h$:
\[
    \bigl(1 - \beta\bigl(1 - 2^{-(j-i+1)}\bigr)\bigr)L_0^{2\max\{0,\, i-n_0-1\}}\varepsilon_i\Bigl(\frac{\eta}{a_i}\Bigr)^2\;.
\]
\item If $j \leq n_0$, for $p \subset (X \setminus \Omega_{j+1}) \setminus Z''_j$, with the same convention for $j = h$: $(1 - \frac{1}{2}\beta)\varepsilon_j(\eta/a_j)^2$.
\item If $j > n_0$, for $n_0 \leq l \leq j$ and $p \subset (X \cap (\Omega_l \setminus \Omega_{l+1})) \setminus Z''_j$, where $\Omega_{n_0}$ is read as $X$ and $\Omega_{n+1}$ as the empty set:
\[
    c\,\bigl(1 - \tfrac{1}{2}\beta\bigr)L_0^{2\max\{0,\, j-l-1\}}\varepsilon_j\Bigl(\frac{\eta}{a_j}\Bigr)^2\;,
\]
with $c = 3$ if $j = n$ and $l \geq n - 1$, and $c = 1$ otherwise.
\end{enumerate}
Let $\kappa_c \coloneqq \kappa_s/(N + 1)$ and $\bar\chi_c(x) \coloneqq \min\{ 1, \max\{ 0, (1 + \kappa_c - x)/\kappa_c \} \}$. For a cube $\square$ of $\mathsf s_i$ in $X$, $h \leq i \leq n$, let $f^{(m)}_\square$ be the supremum of $\abs{U^{(m)}_{n,X}(\partial p) - \id}/T_m(p)$ over the plaquettes $p$ whose corner with the smallest coordinates lies in $\square$ and whose threshold is one of the numbers above with the factor $\varepsilon_i$, and $f^{(m)}_\square \coloneqq \infty$ if $U^{(m)}_{n,X}$ is not defined. The $m$-th function of the cascade is
\begin{equation}\label{eq:cascade}
    \chi^{(m)}_X \coloneqq \prod_\square\bar\chi_c\Bigl(\frac{f^{(m)}_\square}{(1 + \kappa_c)^{N-m}}\Bigr)\;.
\end{equation}
\end{definition}

For $\kappa_s \to 0$ these are the characteristic functions $\chi^{(n)}_k$ of \cite[(1.22)--(1.24)]{BalabanLFR1}, with $j \leq n_0$ covering (1.22), (1.23), and their generalisations, and the widths are nested as the comparison of two constructions requires: the function $\chi^{(m)}_X$ equals one where every $f^{(m)}_\square$ is at most $(1 + \kappa_c)^{N-m}$ and vanishes unless every $f^{(m)}_\square$ is less than $(1 + \kappa_c)^{N-m+1}$, and $(1 + \kappa_c)^{N+1} \leq e^{\kappa_s}$.

\begin{definition}[The Tests of a Preparatory Integration]\label{def:preparatory-tests}
Let $X$ be small and old, $0 \leq m < N$, and $j \coloneqq h + m$.
\begin{enumerate}[label=(\roman*)]
\item \emph{The deviation from the reference field.} $V^{(j)}_X \coloneqq M^j(U^{(m+1)}_{n,X})$, and $\chi'_j \coloneqq \bar\chi(f'_j/(2\delta'_j))$ with $f'_j \coloneqq \sup_{b \in D^*_j}\abs{V^{(m)}_j(b)V^{(j)}_X(b)^{-1} - \id}$.
\item \emph{The fluctuation field of the label.} $\chi^{(j)}_{\mathrm{cpl}}$ is the product, over the cubes $\square'$ of $\mathsf s_{j+1}$ in $X \cap (\Omega_{j+1} \setminus \Omega^{\sim-1}_{j+1})$, of the characteristic functions of $\{ f''_{\square'}(A_j) < \vartheta_j \}$, with the functions of Definition~\ref{def:step-regions}(v) of the scale $j$ and the fluctuation field $A_j$ of the label.
\item \emph{The fluctuation field of the integration.} $V^{(j)} \coloneqq M^j(U^{(m+1)}_n)$, and $B_j$ is the function on $D^*_j$ with values in $\mathfrak g$ for which
\begin{equation}\label{eq:preparatory-field}
    V^{(m)}_j = \exp\bigl(i\bigl[g_jCB_j - h\tilde D(g_jCB_j)\bigr]\bigr)\,V^{(j)}
\end{equation}
on the bonds of the layer $j$ of $\mathbf B^{(m)}$ with an end in $D_j$, with the objects of Section~\ref{sec:fluctuation-field} for the pair of the $m$-th and the $(m + 1)$-st sequence and the region $D_j$ in the place of $\mathcal O$; $\chi'^{(j)}$ is the characteristic function of $\{ g_j\abs{B_j(b)} < \delta'_j \text{ for all } b \in D^*_j \}$.
\end{enumerate}
\end{definition}

These are the functions \cite[(1.27), (1.28), (1.55)]{BalabanLFR1}. In (iii) the setting of Definition~\ref{def:step-setting} is used for two sequences which differ in one domain and have further layers of higher scales; the definitions and the statements of Section~\ref{sec:fluctuation-field} are local and apply to it with $\mathsf b(\mathcal O)$ replaced by the bonds of the layer $j$ with an end in $D_j$.

The terms which an integration produces are analytic on spaces of pairs which shrink with $m$, as the thresholds of the cascade do \cite[(1.64)--(1.69), pp.~190--191]{BalabanLFR1}.

\begin{definition}[The Spaces of the Preparatory Integrations]\label{def:preparatory-spaces}
Let $X$ be small and old, $0 \leq m \leq N$, $j \coloneqq h + m$, and for $1 \leq i \leq n$
\[
    b_i \coloneqq 1 - \beta\sum_{l=h+1}^{j}2^{-\abs{i-l}} - \beta\sum_{q=i+1}^{n}2^{-(q-i)}\;.
\]
For a union $X'$ of cubes of $\pi_n$ the set $\tilde{\mathfrak U}^{(m)c}_n(X'; \tilde\alpha_0, \tilde\alpha_1)$ is defined as the set of Definition~\ref{def:pair-space-layers} for the $m$-th sequence, with the following scales and factors in the place of the scale $n$ of a layer and of the factor $c_{j,n}$:
\begin{enumerate}[label=(\roman*)]
\item on $X' \cap (\Omega''_i \setminus \Omega''_{i+1})$ for $i < j$, and on $X' \cap (\Omega''_j \setminus \Omega_{\max\{j, n_0\}+1})$ for $i = j$: the scale $i$ and the factor $b_iL_0^{2\max\{0,\, i-n_0\}}$;
\item on $X' \cap (\Omega_i \setminus \Omega_{i+1})$ for $n_0 < i \leq j$: the scale $j$ and the factor $b_jL_0^{2(j-i)}$;
\item on $X' \cap (\Omega_i \setminus \Omega_{i+1})$ for $j < i \leq n$: the scale $i$ and the factor $b_i$, multiplied by $3$ for $i = n$;
\end{enumerate}
the bounds of Definition~\ref{def:pair-space-layers}(iii) being multiplied by the same powers of $L_0$.
\end{definition}

Since $\beta \leq \frac{1}{4}$ the numbers $b_i$ are positive, and since $L_0^2 \leq \frac{1}{3}L$ the sets decrease as $m$ increases \cite[p.~191]{BalabanLFR1}.

\begin{proposition}[A Preparatory Integration]\label{prop:preparatory}
Let the constants of the construction be chosen (Section~\ref{sec:labels}), let the system of parameters be admissible, with the thresholds $\delta'_j$ and the depth of Definition~\ref{def:admissible}, and let $\varrho$ be a fresh density of the order $n$. Let $X$ be a small and old component of a label, $0 \leq m < N$, and $j \coloneqq h + m$. Then:
\begin{enumerate}[label=(\roman*)]
\item at the fields of the label at which $\chi^{(m+1)}_X\,\chi'_j\,\chi^{(j)}_{\mathrm{cpl}}$ does not vanish, the function $\chi^{(m)}_X$ and the factors of the test functions of the label which belong to the cubes of $\mathsf s_j$ in $X \cap (\Omega^{\sim4}_j \setminus \Omega^\sim_{j+1})$, to the cubes of $\mathsf s_{j+1}$ in $X \cap (\Omega^{\sim4}_{j+1} \setminus \Omega^\sim_{j+2})$ if $j + 1 < n$, to the gauge fixing of the step $j + 1$ in $X$, and to the deviations of the step $j + 1$ in $X$ are equal to one;
\item at the fields at which, in addition, $\chi'^{(j)}$ does not vanish, the integral over $B_j$ and over the field $A_j$ of the label in $X$ of the factors of the term which depend on them is the exponential of a constant and of a sum $\sum_{X'}\mathbf C^{(m+1)}(X', U^{(m+1)}_n)$ over the unions $X'$ of cubes which, for some $i$ with $j + 1 \leq i \leq n$, belong to $\mathbf D_i$, lie in the complement of $\Omega_{i+1}$, meet $\Omega_i$, and meet $(Z_{j+1} \cap X) \setminus Z''_{j+1}$; each $\mathbf C^{(m+1)}(X', \cdot)$ is holomorphic on $\tilde{\mathfrak U}^{(m+1)c}_n(X'; \tilde\alpha_0, \tilde\alpha_1)$, invariant under the gauge transformations with values in $G$, and bounded there by $B_0e^{-(1+3\beta)\kappa d_i(X')}$.
\end{enumerate}
\end{proposition}

This is \cite[(1.29), (1.60)--(1.63), (1.70)]{BalabanLFR1}. The integral of (ii) is a fluctuation integral of the kind of Section~\ref{sec:step}, in a region which lies in the region of large fields, so that all its terms are boundary terms. Where $\chi'^{(j)}$ vanishes the dependence of $\chi'_j$ on the background is removed by the change of variables of Definition~\ref{def:boundary-field}, with the field $\frac{1}{i}\operatorname{Log}(V^{(j)}(V^{(j)}_X)^{-1})$, the bump function $\varphi_1(g_j\abs{B'}/\delta'_j)$, and all the bonds of $D^*_j$ \cite[(1.59)]{BalabanLFR1}. After the $N$ integrations the fields of the label in $X$ are the data $V^{(N)}$ on the layers of $\{ \Omega''_i \}$ and the fields of the label of the order $h$ in $Z_h \cap X$, the background is $U''_n \coloneqq U^{(N)}_n$, and the effective action has the new boundary terms $\mathbf C^{(1)}, \dots, \mathbf C^{(N)}$ \cite[(1.71), (1.72)]{BalabanLFR1}.

\subsection{The Minimiser with Free Data}\label{sec:free-data}

The region is integrated out about the configuration of least action in which the field of the last scale is prescribed outside a box and free inside it. Let $1 \leq n \leq k$, let $Z$ be a union of cubes of $\mathsf r_n$, and for a gauge field $V_n$ on the bonds of $T^{(n)}$ in $Z$ let $U_{n,Z}(V_n)$ be the minimal configuration, on the support $Z$ and in the axial gauge, of the minimal sequence of $Z$ of the order $n$ for the data which on the layer $i$ are $M^i(Q^{s*}_nV_n)$, as in Definition~\ref{def:cube-minimiser}; it is defined when these data are regular with a constant $\varepsilon \leq a^{\mathrm v}_1$.

\begin{proposition}[The Minimiser with Free Data]\label{prop:free-data}
There are numbers $c_9 > 0$ and $B_8$, which depend on $L$ and $G$ only, with the following property. Let $\Lambda \subset Z$ be a product of intervals which is a union of cubes of the grid of side $Ma_n$, contained in a cube of side $100Ma_n$ and at a distance of at least $Ma_n$ from the complement of $Z$, and let $V_n$ be given on the bonds of $T^{(n)}$ in $Z$ with no end in $\Lambda$, with $\abs{V_n(\partial p') - \id} \leq \varepsilon \leq c_9M^{-5}$ for the plaquettes $p'$ of $T^{(n)}$ in $Z$ with no corner in $\Lambda$. Then:
\begin{enumerate}[label=(\roman*)]
\item among the extensions of $V_n$ to the bonds with an end in $\Lambda$ with plaquettes of at most $a^{\mathrm v}_1$, the critical points of $V_n \mapsto A_Z(U_{n,Z}(V_n))$ form exactly one orbit of the gauge transformations of the points of $T^{(n)}$ in $\Lambda$, and they are minima;
\item an element $V_\Lambda$ of this orbit satisfies $\abs{V_\Lambda(\partial p') - \id} \leq B_8M^5\varepsilon$ for the plaquettes $p'$ of $T^{(n)}$ with a corner in $\Lambda$;
\item in the axial gauge of $\Lambda$, $V_\Lambda$ extends to a holomorphic function of the outer field $e^{iB'}V_n$ for functions $B'$ with values in $\mathfrak g^c$ and $\abs{B'} \leq \varepsilon$.
\end{enumerate}
\end{proposition}

This is \cite[Proposition~1, p.~194]{BalabanLFR1}, proved in \cite[p.~359]{BalabanLFR2} for the threshold of the step by a contraction for the equation of the critical point in the axial gauge of the box; the linear part of that equation is invertible by the following inequality, in which no average is involved.

\begin{lemma}[The Axial Gauge of a Box]\label{lem:axial-poincare}
Let $\Lambda = \prod_\mu\{ a_\mu, \dots, b_\mu \} \subset \mathbb Z^4$ with $b_\mu - a_\mu < \ell$ for all $\mu$, and let $B$ be a real function on the bonds of $\mathbb Z^4$ with both ends in $\Lambda$ which vanishes on the bonds of the tree which consists, for every point $x \in \Lambda$, of the path from $(a_0, a_1, a_2, a_3)$ to $x$ which changes the coordinates in the order $3, 2, 1, 0$. Then
\begin{equation}\label{eq:axial-poincare}
    \sum_b\abs{B(b)}^2 \leq 4\ell^5\sum_p\abs{(\partial B)(p)}^2\;,
\end{equation}
the sums over the bonds and the plaquettes of $\Lambda$, with $(\partial B)(p)$ the oriented sum of $B$ around $p$.
\end{lemma}

This is \cite[(1.8), p.~358]{BalabanLFR2}. A bond which is not in the tree closes, with two paths of the tree, a loop which bounds a strip of at most $\ell$ plaquettes in each of at most four planes, and $B$ on the bond is the sum of $\partial B$ over these plaquettes.

\section{The Densities of All Orders}\label{sec:inductive}

\subsection{The Form of a Density}\label{sec:density-form}

\begin{definition}[Terms of the Operation and Boundary Terms]\label{def:other-terms}
Let a system of parameters and a label of the order $n$ be given, let $1 \leq j \leq n$, and let $\kappa$, $B_0$, $\varrho_j$, and $\beta$ be positive numbers, $\beta < 1$.
\begin{enumerate}[label=(\roman*)]
\item A \emph{family of $\mathbf R$-terms of the scale $j$} is a family of functions $\mathbf R(X, \cdot)$, one for each $X \in \mathbf D_j$ which is a union of cubes of $\mathsf r_j$, with the properties (i) and (ii) of Definition~\ref{def:regular-terms} on $\mathfrak U^c_j(X; \alpha_{0,j}, \alpha_{1,j})$, with the bound $\abs{\mathbf R(X, \cdot)} \leq \varrho_je^{-\kappa d_j(X)}$, and with $\mathbf R(rX, (rU, rJ(U))) = \mathbf R(X, (U, J(U)))$ for the isometries $r$ of the torus which preserve the grid $\mathsf r_j$.
\item A \emph{family of boundary terms of the scale $j$} for the label is a family of functions $\mathbf B(X, \cdot)$, one for each $X \in \mathbf D_j$ which meets $\Omega_j$ and $Z_j^\sim$, of the pairs on $X$ and of the fluctuation fields of the label on $X$, holomorphic in the pair on $\tilde{\mathfrak U}^c_j(X; \tilde\alpha_0, \tilde\alpha_1)$ (Definition~\ref{def:pair-space-layers}, with the constant $\beta$), invariant under the simultaneous gauge transformations with values in $G$ of the pair and of the fluctuation fields, bounded by $B_0e^{-\kappa d_j(X)}$, and independent of the fields on the connected components of $X \cap \Omega_1^\sim$ which do not meet $\Omega_j$.
\end{enumerate}
\end{definition}

\begin{definition}[Densities of the Order $n$]\label{def:density}
Let a system of parameters be given. A function $\varrho$ of the gauge fields $V_n$ on $T^{(n)}$ is a \emph{density of the order $n$} with the constants $(\kappa, E_0, B_0, (\varrho_j), \beta)$ if there are numbers $\beta_1, \dots, \beta_n$ and families $\mathbf E^{(j)}$ of regular terms of the scales $j \leq n$ with the constants $(\alpha_{0,j}, \alpha_{1,j}, E_0, \kappa)$, and for every label of the order $n$ families $\mathbf R^{(j)}$ of $\mathbf R$-terms and $\mathbf B^{(j)}$ of boundary terms, a number $E$, and an operation $\mathbf T$, such that
\begin{equation}\label{eq:density}
    \varrho(V_n) = \sum_{\text{labels}}\chi_n(\Omega_n)\,\mathbf T\Bigl[\exp\Bigl(-A\bigl(g_n(\cdot)^{-2}, U_n\bigr) + \mathbf E_n + \mathbf R_n + \mathbf B_n - E\Bigr)\Bigr]\;,
\end{equation}
where:
\begin{enumerate}[label=(\roman*)]
\item $U_n$ is the minimal configuration, on a support, for the data of the label which are regular near $\Omega_1$, and $A(g(\cdot)^{-2}, U)$ is the action \eqref{eq:action} with the plaquette at $x$ weighted by $g(x)^{-2}$, where $g_0(x) \coloneqq g_0$ and $g_{j-1}(x)^{-2} = g_j(x)^{-2} + \beta_j\phi_j(x)$ with functions $\phi_j$ with values in $[0, 1]$, supported in $\Lambda_j$ and equal to one on $\Lambda_j^{\sim-1}$;
\item $\mathbf E_n \coloneqq \sum_{j=1}^{n}\sum_z\bigl[\mathbf E^{(j)}(\Lambda_j, U_n, z) - \mathbf E^{(j)}(\Lambda_j, \id, z)\bigr] - \sum_{j=1}^n\beta_jA(\phi_j, U_n)$, the sum over the points $z$ of $T^{(j)}$ in $\Lambda_j^{\sim-1}$, with the sums \eqref{eq:regular-sum};
\item $\mathbf R_n$ is the sum over $j \leq n$ and over the domains $X \subset \Lambda_j^{\sim-1}$ of $\mathbf R^{(j)}(X, U_n) - \mathbf R^{(j)}(X, \id)$, and $\mathbf B_n$ is the sum over $j \leq n$ and $X$ of the boundary terms at $U_n$ and at the fluctuation fields of the label;
\item $\mathbf T$ is an integral, with respect to a positive measure which depends on the label, over the fields of the label in the region $Z_n$ of large fields, the measure being a product over the connected components of $Z_n$ of measures on the fields in the components.
\end{enumerate}
The operations ${}^\sim$ on $\Lambda_j$ are those of the grid $\mathsf r_j$.
\end{definition}

This is the inductive assumption of \cite[Sect.~2, (2.17)--(2.42)]{BalabanCRE}. As there, the operation $\mathbf T$ is described here by its general properties only; its form is fixed by the steps and by the operation of Section~\ref{sec:r-operation}.

\subsection{The Parameters}\label{sec:admissible}

The parameters of the scales are tied to a sequence of reference couplings, which the running couplings follow within a fixed factor.

\begin{definition}[Admissible Systems of Parameters]\label{def:admissible}
Let $\beta_0 \in (0, \sqrt 2 - 1]$, $\gamma > 0$, exponents $r \geq 2$, $p_0$, $p_1 < p_0$, $q_0$, $q_1$, and constants $A_0$, $A_1$, $C^\alpha_0$, $C^\alpha_1$ be given. A system of parameters, together with numbers $0 < \hat g_0 \leq \hat g_1 \leq \dots \leq \hat g_k \leq \gamma$, is \emph{admissible} if, for $0 \leq m < n \leq k$ and $0 \leq j \leq k$:
\begin{enumerate}[label=(\roman*)]
\item $\hat g_j \leq g_j \leq (\frac{3}{2})^{1/2}\hat g_j$;
\item $g_n \leq (1 + \beta_0)(n - m)^{1/2}g_m$ and $g_m \leq (1 + \beta_0)g_n$, and the same for the numbers $\hat g$;
\item $R_j$ is the smallest power of $L$ which is at least $(\log\hat g_j^{-2})^r$;
\item $\varepsilon_j = A_0\hat g_j(\log\hat g_j^{-2})^{p_0}$, $\vartheta_j = (A_1/A_0)\varepsilon_j$, $\alpha_{0,j} = C^\alpha_0\hat g_j(\log\hat g_j^{-2})^{q_0}$, and $\alpha_{1,j} = C^\alpha_1\hat g_j(\log\hat g_j^{-2})^{q_1}$.
\end{enumerate}
The threshold of the fluctuation fields of the operation of Section~\ref{sec:r-operation} is $\delta'_j \coloneqq A_1\hat g_j(\log\hat g_j^{-2})^{p_1}$, and its depth at the step $j$ is the larger of $N_0 + 1$, with $N_0$ the smallest positive integer with $L^{N_0-1} = R_{j-N_0+1}$, and of the least integer $N'$ with $L^{-N'} \leq (\log\hat g_j^{-2})^{-(p_0+q_1+7r+1)}$.
\end{definition}

For $\hat g_j = g_j$ these are the parameters of \cite[(2.4)--(2.9), (2.28)]{BalabanCRE} and \cite[p.~183]{BalabanLFR1}, where (ii) is derived from the flow of the couplings; with reference couplings they are the second of the two choices of \cite[p.~246]{BalabanCRE}. The depth is not fixed in the papers, which state four conditions on it \cite[pp.~179, 198]{BalabanLFR1}, \cite[pp.~361, 363]{BalabanLFR2}; the choice above satisfies them and is of the order of $\log\log\hat g_j^{-2}$.

\subsection{The Theorem}\label{sec:main-theorem}

\begin{theorem}[The Densities of All Orders]\label{thm:inductive}
Let $\mathfrak g$ be simple. The constants $L$; $M$ and $M_2$; $\kappa$, $E_0$, $B_0$, $\beta$, and $\kappa_s$; the exponents and the constants of Definition~\ref{def:admissible}; and $\gamma$ can be chosen, in this order, such that the following holds for every $k$, every torus, and every admissible system of parameters. Let $\varrho_0(V_0) \coloneqq \exp(-g_0^{-2}A(V_0))$ on the gauge fields of $T_\eta$, and for $0 \leq n < k$ let $\varrho_{n+1}$ be obtained from $\varrho_n$ by the step of Section~\ref{sec:step}, in which the coefficient $\beta_{n+1}$ is that of the family $\mathbf E^{(n+1)}$ of Theorem~\ref{thm:small-field-step}, followed by the operation of Section~\ref{sec:r-operation}, provided that $(g_n^{-2} - \beta_{n+1})^{-1/2}$ is the coupling $g_{n+1}$ of the system. Then:
\begin{enumerate}[label=(\roman*)]
\item $\int f(V_n)\varrho_n(V_n)\,\dd V_n = \int f(M^n(V_0))\varrho_0(V_0)\,\dd V_0$ for $f = 1$, and for every bounded Borel function $f$ of the field on a set of bonds of $T^{(n)}$ at a distance of at least $100MR_ja_j$ from the regions which the operation has treated at the steps $j \leq n$;
\item $\varrho_n$ is a density of the order $n$ with the constants $(\kappa, E_0, B_0, (\varrho_j), \beta)$, where $\varrho_j \coloneqq \exp(-A_0(\log\hat g_j^{-2})^{p_0})$, and its regular terms and coefficients do not depend on the label, on $n$, on $k$, or on the torus;
\item for $1 \leq j \leq n$, every label, every union $\Omega$ of cubes of $\pi_n$ in $\Lambda_j^{\sim-1}$, every function $\phi$ with values in $[0, 1]$ which equals one on $\Omega$ and vanishes outside $\Omega^\sim$, every $\beta' < 1$, and the background $U_n$ of data for which the test functions of the term do not vanish,
\begin{equation}\label{eq:theorem-two}
    \Bigl\lvert\sum_{z \in \Omega}\bigl[\mathbf E^{(j)}(\Lambda_j, U_n, z) - \mathbf E^{(j)}(\Lambda_j, \id, z)\bigr] - \beta_jA(\phi, U_n)\Bigr\rvert \leq E_1(\beta')\sum_{i=j}^{n}L^{-\beta'(i-j)}\#(\Gamma_i \cap \Omega)\;,
\end{equation}
and the corresponding sum of the $\mathbf R$-terms of the scale $j$ is at most $\varrho_j\sum_{i=j}^{n}\#(\Gamma_i \cap \Omega)$, with a number $E_1(\beta')$ which depends on the constants of the theorem only.
\end{enumerate}
\end{theorem}

The second and the third part are \cite[Theorems~1 and~2, pp.~262--263]{BalabanCRE} with \cite[Theorem~1, p.~355]{BalabanLFR2}, for an admissible system of parameters and soft test functions; the first part records what the steps and the operation preserve: a step is the image of the density under the average, by Theorem~\ref{thm:step-identity}, and the operation changes a density only in the regions which it treats \cite[(1.102), p.~201]{BalabanLFR1}. The condition on the couplings says that the running couplings of the construction form, together with the reference couplings, an admissible system; whether they do is a statement about the coefficients $\beta_j$, which is not part of the construction \cite[Theorem~2, p.~259]{BalabanRG1}.

\bibliographystyle{amsplain}
\bibliography{refs}

\end{document}
